% SPDX-FileCopyrightText: 2025-2026 Devon Thiele
% SPDX-License-Identifier: CC-BY-SA-4.0
\RequirePackage{fix-cm}
\documentclass[11pt,twoside]{article}
\usepackage{longtable}

\usepackage[utf8]{inputenc}
\usepackage[T1]{fontenc}
\usepackage{mathpazo}
\usepackage[scaled=0.90]{helvet}
\usepackage{courier}
\usepackage{geometry}
\usepackage{amsmath,amssymb,amsthm}
\usepackage{booktabs}
\usepackage{enumitem}
\usepackage{xcolor}
\usepackage{xr-hyper}
\externaldocument{monograph}
\usepackage{hyperref}
\usepackage[
  activate={true,nocompatibility},
  final,
  tracking=true,
  factor=1100,
  stretch=20,
  shrink=20
]{microtype}
\usepackage{array}
\usepackage[skip=6pt plus 1pt]{parskip}
\usepackage{setspace}
\usepackage{titlesec}
\usepackage{fancyhdr}
\usepackage{seqsplit}
\usepackage{etoolbox}
\usepackage{tikz}
\usepackage{tikz-cd}
\usetikzlibrary{arrows.meta,positioning,fit,backgrounds,calc,shapes.geometric,decorations.pathreplacing,shadings}

\geometry{top=1in, bottom=1in, left=1.15in, right=1.15in}
\setstretch{1.08}

\definecolor{linkblue}{RGB}{20,70,140}
\hypersetup{pdftitle={The Thiele Machine},pdfauthor={Devon Thiele},colorlinks=true,linkcolor=linkblue,citecolor=linkblue,urlcolor=linkblue}

\setlist{itemsep=0.45em, topsep=0.55em, parsep=0pt}
\emergencystretch=12em
\tolerance=2500
\hbadness=4000

% Hyphenation hints for stubborn proper names and compound words.
\hyphenation{
  Thiele
  bisim-u-la-tion
  poly-no-mi-al
  cer-ti-fi-ca-tion
  com-pu-ta-tion
  com-pu-ta-tional
  sub-strate
  con-struc-tion
  cat-e-gor-i-cal
  semi-def-i-nite
  re-al-iz-able
  Tsirelson
  Hilbert
  Kleene
  G\"odel
  Schur
  Landauer
  Bekenstein
}

% Section heading style: bold, ruled underneath
\titleformat{\section}{\large\bfseries}{\thesection}{0.8em}{}[\vspace{2pt}{\titlerule[0.4pt]}]
\titleformat{\subsection}{\normalsize\bfseries}{\thesubsection}{0.75em}{}
\titleformat{\subsubsection}{\normalsize\itshape}{\thesubsubsection}{0.75em}{}
\titlespacing*{\section}{0pt}{2.8ex plus 0.5ex minus 0.2ex}{1.4ex plus 0.2ex}
\titlespacing*{\subsection}{0pt}{1.8ex plus 0.3ex minus 0.1ex}{0.7ex plus 0.1ex}
\titlespacing*{\subsubsection}{0pt}{1.2ex plus 0.2ex}{0.4ex}

\pagestyle{fancy}
\fancyhf{}
\fancyhead[L]{\small\itshape The Thiele Machine}
\fancyhead[R]{\small\itshape Devon Thiele}
\fancyfoot[C]{\small\thepage}
\renewcommand{\headrulewidth}{0.4pt}
\renewcommand{\footrulewidth}{0pt}

\DeclareUnicodeCharacter{03BC}{\ensuremath{\mu}}
\DeclareUnicodeCharacter{03B4}{\ensuremath{\delta}}

% === Theorem / proof / plain-reading visual rhythm =========================
% Theorems get extra breathing room; proofs get a small lead and tail; and
% the "Plain reading" commentary lives in its own indented block with a soft
% blue rule on the left so it reads as authorial gloss, not another theorem.

\newtheoremstyle{thielethm}
  {1.2\baselineskip}     % above
  {0.7\baselineskip}     % below
  {\itshape}             % body font
  {0pt}                  % indent
  {\bfseries}            % head font
  {.}                    % head punctuation
  {0.55em}               % space after head
  {}                     % default head spec

\newtheoremstyle{thieledef}
  {1.2\baselineskip}
  {0.7\baselineskip}
  {\normalfont}
  {0pt}
  {\bfseries}
  {.}
  {0.55em}
  {}

\theoremstyle{thielethm}
\newtheorem{theorem}{Theorem}[section]
\newtheorem{lemma}[theorem]{Lemma}
\newtheorem{corollary}[theorem]{Corollary}
\newtheorem{proposition}[theorem]{Proposition}

\theoremstyle{thieledef}
\newtheorem{definition}{Definition}[section]

\theoremstyle{thieledef}
\newtheorem{remark}{Remark}[section]

% A little extra room around proofs so they don't crowd the surrounding text.
% (Implemented by patching the proof environment definitions below, not via
% AtBeginEnvironment, to avoid list-context interference inside \qed.)

% === TikZ shared styling ===================================================
% A single visual language for every figure: Palatino-compatible text,
% restrained blue accents matching linkblue, thin rules at 0.4-0.6pt to
% rhyme with the section dividers.
\definecolor{thiele@accent}{RGB}{20,70,140}    % matches linkblue
\definecolor{thiele@soft}{RGB}{150,170,200}
\definecolor{thiele@gray}{RGB}{120,120,120}
\definecolor{thiele@red}{RGB}{170,40,40}
\definecolor{thiele@pale}{RGB}{235,240,250}

\tikzset{
  thiele/.style={
    font=\small,
    >={Stealth[length=4pt,width=4pt]},
    line width=0.5pt,
    every node/.style={font=\small}
  },
  tnode/.style={                   % Thiele state node
    draw=thiele@accent, line width=0.55pt,
    rounded corners=2pt,
    inner sep=4pt,
    fill=thiele@pale,
    font=\scriptsize\ttfamily,
    minimum height=2.6em
  },
  cnode/.style={                   % classical state node
    draw=thiele@gray, line width=0.55pt,
    rounded corners=2pt,
    inner sep=4pt,
    fill=white,
    font=\scriptsize\ttfamily,
    minimum height=2.0em
  },
  proj/.style={->, line width=0.55pt, draw=thiele@accent},
  nolift/.style={->, line width=0.7pt, draw=thiele@red, dashed},
  pipe/.style={
    draw=thiele@accent, line width=0.5pt,
    rounded corners=3pt,
    inner sep=5pt,
    fill=thiele@pale,
    align=center,
    font=\small,
    minimum width=11em,
    minimum height=2.0em
  },
  pipearrow/.style={->, line width=0.6pt, draw=thiele@accent},
  trustarrow/.style={->, line width=0.6pt, draw=thiele@red},
  fcaption/.style={font=\footnotesize\itshape, text=thiele@gray}
}

% Plain reading: authorial gloss after a theorem/proof. Indented block,
% with a soft blue left rule and a bold-italic lead-in.
\newenvironment{plainreading}{%
  \par\addvspace{0.85\baselineskip}%
  \begin{list}{}{%
    \setlength{\leftmargin}{1.6em}%
    \setlength{\rightmargin}{0.2em}%
    \setlength{\topsep}{0pt}%
    \setlength{\itemsep}{0pt}%
    \setlength{\parsep}{0.35\baselineskip}%
    \setlength{\labelwidth}{0pt}%
    \setlength{\labelsep}{0.7em}%
  }%
  \item[\textcolor{linkblue!55}{\rule[-0.18ex]{1.4pt}{0.95em}}]%
  \hspace*{-0.35em}\textit{\textbf{Plain reading.}}~\ignorespaces
}{%
  \end{list}%
  \addvspace{0.45\baselineskip}%
}
\theoremstyle{thieledef}
\newtheorem{premise}[definition]{Premise}
\theoremstyle{thielethm}
\newtheorem{conjecture}[theorem]{Conjecture}
\newtheorem{thesis}[theorem]{Thesis}
\theoremstyle{thieleobs}
\newtheorem{observation}[theorem]{Observation}
\theoremstyle{thielethm}

% The book's setting: trim size, Charter body, sans headings, open spacing.
\usepackage{book}

\begin{document}

\hypersetup{pageanchor=false}
\pagestyle{empty}
\vspace*{112pt}
\noindent\BookLabel{Devon Thiele}

\vspace{42pt}
\noindent{\fontsize{34}{38}\selectfont The Thiele\\Machine}

\vspace{7pt}
\noindent{\fontsize{16}{20}\selectfont\itshape The Example, Built}
\clearpage

\vspace*{\fill}
\begingroup\footnotesize\setlength{\parindent}{0pt}\setlength{\parskip}{5pt}
\noindent\textit{The Thiele Machine: The Example, Built}

Devon Thiele

Copyright 2025--2026 Devon Thiele. Licensed under Creative Commons Attribution-ShareAlike 4.0 (CC BY-SA 4.0).

A companion to \textit{The Thiele Machine: No Free Insight}. The source of both, the proofs and the build are public:
\href{https://github.com/sethirus/The-Thiele-Machine}{github.com/sethirus/The-Thiele-Machine}.
\endgroup
\cleardoublepage
\pagestyle{book}
\hypersetup{pageanchor=true}

\phantomsection\section*{What this volume is}

The book \textit{The Thiele Machine: No Free Insight} develops the model and the argument. This volume holds, word for word, what was written for the book and would have doubled its length: the proofs of the theorems the book states about its example machine and its universal programs, the register-by-register programs and the worked runs, the full algebra of the CHSH check, the book's own ledger of its claims and its assumptions, and the works read for the book that its finished text does not name. Theorem and section numbers cited here and not found here are the book's.

\tableofcontents
\clearpage

\section{Proofs for ``What it takes to be certified''}

\paragraph{Proof of the theorem ``Certified is Thiele-complete plus certification'' (book Theorem~\ref{thm:tc-decomposition}).}
\begin{proof}
(1) Clause (a) is (A) together with ``loaded states are clean''. Clause (c) is (C), with the three
record kinds as the marks. Clause (b) says clean states have the record down, which with clean
loads gives (B), and it gives the earned chain. Clause (d) gives (D): the claim's three moves end
with the record up from the loaded state where it holds and down from the one where it fails.
Going back, the same pieces put the four clauses together again. (2) is (1) read left to right.
(3) From any state with the record down, counting $A$ to zero costs nothing and the first mark
then raises the record: a raise for one mark from every clean state, where a certified raise from a
clean state takes three (Theorem~\ref{thm:u34-tc-consequences}(3)).
\end{proof}

\paragraph{Proof of the theorem ``CHECK is a unit-cost test'' (book Theorem~\ref{thm:u34-check-unit}).}
\begin{proof}
At a start the latch is down and the table is empty, so the check passes exactly when
$\mathrm{eval}(\mathsf{Ge}(m), a)$ says yes, that is when $m \le a$; a failing check sets the latch.
The price table charges a CHECK 1.
\end{proof}

\paragraph{Proof of the theorem ``What every certified Thiele-complete machine has'' (book Theorem~\ref{thm:u34-tc-consequences}).}
\begin{proof}
(1) A raise of the record isn't made by a base move, by (a), so it's a record move, which costs
one by (c). Running a program on the base is induction on its steps, each matched by (a).
(2) is (c), added up along the run. (3) Clause (b) finds a CHECK, a COMMIT and a CERTIFY at three
different places, and (2) counts them. (4) is (d): the three moves certify exactly where the
claim holds, the claim fails at some loaded state, and the check can't pass there because (b)
says a passing check means the claim holds; at the loaded state where it holds, the three moves
raise the record. (5) Cut the run just after the move where the record rises. That shorter run
starts clean and ends with the record up, so clause (b) puts the record down right before a
CERTIFY in it and up right after. The record never comes down, so it rises only once, and that
CERTIFY is the move.
(6) Every move is a base move or a record move, and (a) supplies base moves.
\end{proof}

\paragraph{Proof of the theorem ``The language first, and what a certificate then means'' (book Theorem~\ref{thm:u34-tc-over}).}
\begin{proof}
(1) Take the plain interface the language-first one becomes. (2) The checker's soundness and the
keeping proof are the ones clause (b) already supplies. (3) Clause (b) gives the chain; the
checker's soundness proof against $L$ gives $c$'s meaning at the state checked; $L$'s keeping
proof carries it, through states where the thing $c$ is about is unchanged, to the state
committed.
\end{proof}

\paragraph{Proof of the observation ``A clock is a Thiele machine and isn't certified Thiele-complete'' (book Observation~\ref{thm:u34-clock}).}
\begin{proof}
Every move costs one, so any raise costs at least one: the toll holds for any record. The
machine's own moves are the two-counter instructions, which gives the universal base. But a
certified Thiele-complete machine has a move priced 0, by Theorem~\ref{thm:u34-tc-consequences}(6), and this
one has none.
\end{proof}

\paragraph{Proof of the theorem ``Each clause does work'' (book Theorem~\ref{thm:u34-tc-clauses}).}
\begin{proof}
(1) The base moves leave the flag alone, TICK never lowers it, the base costs 0 and TICK costs 1,
which gives (a) and (c). Through any interface, (a) gives a loaded state, which is clean, so its
record is down by (b), and a move that puts the record up from every state raises it alone.
Clause (b) would need a CHECK, a COMMIT and a CERTIFY inside that one-move run, and one move
can't be three. (2) The base gives (a) and costs 0, which with no record moves is (c). No run
ends with the record up, so (b) asks nothing of a run, and a checker that always says no is
sound for any meaning. Clause (d) needs a run from some loaded state that raises the record,
and none does.
\end{proof}

\paragraph{Proof of the theorem ``The four clauses are independent'' (book Theorem~\ref{thm:tc-independent}).}
\begin{proof}
(a) Nothing moves a trapped state's flag, so a run that uses DROP ends with the flag down. A run
from a clean state that ends with the flag up never used DROP; it's a run of this chapter's
machine, and Theorem~\ref{thm:u34-cert-earned} gives its chain. DROP is priced 0 and read as a base
move, so (c) holds, and (d) holds as it does for the plain machine. A state with the flag up has
it down after DROP. That breaks the permanence in (a), and permanence is about the machine alone,
so no interface meets (a). (b) FORGE never lowers the flag, costs one and is read as a record move,
so (a), (c) and (d) hold as before. From a clean start FORGE alone raises the flag, and a run of
one move has no room for a CHECK, a COMMIT and a CERTIFY. Through any interface: (d) gives a
loaded state from which the record rises, so its trap is down; (a) makes that state clean, so its
record is down by (b); and FORGE alone raises it, which (b) forbids. (c) Clauses (a), (b) and (d)
never mention a price, and they hold as for the plain machine. Base moves cost one, so (c) fails,
and every certified Thiele-complete machine has a move priced 0 (Theorem~\ref{thm:u34-tc-consequences}(6)). (d)
is Theorem~\ref{thm:u34-tc-clauses}(2).
\end{proof}

\paragraph{Proof of the theorem ``Any unit'' (book Theorem~\ref{thm:u34-tc-unit}).}
\begin{proof}
(1) At unit one every listed price is unchanged.

(2) Changing the size of the unit changes no state transition, claim or
record. Thus the computing, earned-chain and distinguishing-claim requirements
stay as they were. To multiply the prices by $c$, read the old ledger
multiplied by $c$ too. Every addition is then multiplied by the same amount.

In the other direction, count complete groups of $c$ units in the ledger.
There may be a remainder smaller than $c$ at the start. A step adds either
nothing or exactly one complete group, so it leaves that remainder alone
and increases the group count by the divided price. The equation for that
count is
\[
 \left\lfloor\frac{\ell+cr}{c}\right\rfloor
 =\left\lfloor\frac{\ell}{c}\right\rfloor+r.
\]
It gives the required whole-number ledger on the same states and runs.

(3) The move for an ordinary counter instruction is a base move and costs
zero at every unit. A machine charging for every move has no such move.

(4) Find the passing CHECK, matching COMMIT and raising CERTIFY in the earned
chain. They occupy three different places in the run. Each costs $c$, and
all other charges are nonnegative, giving at least $3c$. The claim was true
at its check and stays true up to its commitment because what it concerns
is unchanged. The record rises at that chain's CERTIFY and never comes down;
changing the price unit changes none of these events.

(5) Doubling the small machine's prices gives unit two by (2). Its CHECK
then costs two. At unit one every move must cost zero or one, so this same
priced machine cannot qualify there through another choice of interface.
\end{proof}

\paragraph{Proof of the theorem ``The small machine, the generic one and the sorted one are certified Thiele-complete'' (book Theorem~\ref{thm:u34-tc-instances}).}
\begin{proof}
(1) Clause (a) is Theorem~\ref{thm:u34-simulation} together with the step line: INC, DEC and HALT
never touch the flag, the flag never comes down, and a start is clean with its trap down. Clause (b) is
Theorem~\ref{thm:u34-cert-earned} with checker soundness (Lemma~\ref{lem:u34-eval}), and a meaning
about $c$ keeps holding while $c$'s value doesn't move. Clause (c) is the price table. Clause (d) is
Theorem~\ref{thm:u34-tight} with the claim ``$A$ is zero,'' which holds at $(0, b)$ and fails at
$(5, b)$. (2) and (3) are the same argument with the property language left open. For (3), read
a number as a list in the list code: 0 is the empty list, and $2^x(2y + 1)$ is $x$ followed by
the list numbered $y$ (every positive number splits that way in exactly one way, and $y$ is
smaller than the number, so the reading ends). The meaning of \emph{sorted} is that each entry
is at most the one after it. The checker reads the list from the front: yes on the empty list
and on a list of one entry, and on $x$ followed by a list that starts with $y$, yes exactly when
$x \le y$ and it says yes on that rest. Induction on the length of the list shows it says yes
exactly when the meaning holds, and a language with one property has an exact equality test. The
sorted claim holds of the list numbered 18, which is $[1, 2]$, and fails of 20, which is
$[2, 1]$. (4) is (1) of Theorem~\ref{thm:u34-tc-consequences} and Section~\ref{sec:toll}'s floor.
(5) The interface over the language is the one of (1), (2) and (3) with the meanings handed over
by the language, and the checker's soundness is Lemma~\ref{lem:u34-eval} and its generic form.
\end{proof}

\paragraph{Proof of the observation ``Every certified Thiele-complete machine hides its record and its ledger from its window'' (book Observation~\ref{thm:u34-tc-window}).}
\begin{proof}
Theorem~\ref{thm:u34-tc-consequences}(4) gives a loaded state and a run from it that raises the
record, and (3) says that run paid at least 3. Say it ends with window (line $\ell$, $x$, $y$).
Two-counter instructions can steer any configuration to any other: empty $A$, fill it to $x + 1$,
empty $B$, fill it to $y$, then one DEC of $A$ that jumps to line $\ell$. Those are base moves,
so from the same loaded state they reach the same window, cost nothing, and leave the record down,
by (a), (b) and (c). That's two runs, one window, two records and two ledgers, which is (2). A function
of the window would have to give both answers on that one window, which is (3). For (1), load
$(x + 1, y)$ and take just the last DEC. (4) is (3) read as a Thiele machine,
whose runs are the machine's.
\end{proof}


\section{Proofs and worked runs for ``Parts I and II, run on the example machine''}

\paragraph{Proof of the theorem ``The small machine runs the eight-state machine'' (book Theorem~\ref{thm:frag-small}).}
\begin{proof}
(1) By cases on the move. On a live state CERTIFY runs, because a committed claim is on the channel;
it raises the flag, moves the line on by one and charges one. INC $B$ moves the line on by one and
charges nothing, and it touches only $B$ and $B$'s version, so the fact about $A$ and the
commitment stand and the state stays live. In a jump, each COMMIT finds its fact in the table and
succeeds, charging one, and the INC and the DEC charge nothing and end at line $a$, whose slot is $a$;
the flag isn't touched. (2) After each compiled move, the state is live again,
so the next compiled move starts with exactly the condition just used in (1).
The window agrees after each move and each move adds its own charge. At the end
of the list the windows agree and the ledger has added the sum of those charges.
(3) The eight-state machine's own
toll (Section~\ref{sec:a2-not-arbitrary}) says a run of its moves that takes the flag from down to
up is charged at least one, and (2) carries that charge across. (4) CHECK ($A \ge 0$) passes at
every value of $A$ and the COMMIT after it passes, which takes line 1 to line 3, slot 3, with two
units of cost paid; then (2) applies. A raised flag was raised by a CERTIFY, and
Theorem~\ref{thm:u34-cert-earned} gives the chain.
\end{proof}

\paragraph{Proof of the theorem ``The small machine's merges'' (book Theorem~\ref{thm:frag-merge-split}).}
\begin{proof}
(1) Only CERTIFY changes the flag, and it costs one. The two states with the flag down and up, and
a committed claim, are the same state after CERTIFY. (2) Take a core at line 1 with $A$ at 1 at
version 0, and the same core with $A$ at 0 at version 1. DEC $A$ 2 takes the first to $A = 0$ at
version 1, line 2, a write; it takes the second to line 2 with $A$ still 0 at version 1, because
DEC on zero is no write. Same state, zero cost.
\end{proof}

\paragraph{Proof of the theorem ``Structural entitlement represented'' (book Theorem~\ref{thm:ent-machine}).}
\begin{proof}
(1) is Theorem~\ref{thm:ent-cs}(1). (2) Clause (b) applied to a run from a clean state that ends
with the record up gives the chain, and a passing check means the claim's meaning holds, which
keeps holding while the thing it's about is unchanged. (3) The covering (M6) and
Theorem~\ref{thm:ent-cs}'s count finish it, with
the record moves, which are the paid steps, in place of the charged steps. (4) is clause (c)
added up. (5) is Theorem~\ref{thm:u34-tc-consequences}(3).
\end{proof}

\paragraph{Proof of the theorem ``A floor with no supplied tree'' (book Theorem~\ref{thm:ent-questions}).}
\begin{proof}
Write each rival's answers on a strip, one box per CHECK in the run.
One yes-or-no question gives two possible strips; a second gives four;
each further question doubles the possibilities. With $c$ questions there
are at most $2^c$ strips. Put rivals with the same strip in the same pile.
The premise limits each pile to $m$ rivals. Hence at most $2^c m$ rivals
can be on the original list.

Now $\lceil\log_2 m\rceil$ bits suffice to number $m$ places inside a pile,
and $c$ bits suffice to name its strip. Together they number every rival
using at most $c+\lceil\log_2 m\rceil$ bits. The least number needed for
$|\Omega|$ rivals cannot be larger. Subtract the inside-pile allowance,
clipping a negative result to zero, to get the claimed bound of $c$.

Each CHECK is a record move. Counting all record moves can only increase
that number, and each contributes exactly one to the ledger while base
moves contribute zero. This gives the two remaining comparisons.
\end{proof}

\paragraph{Proof of the theorem ``Structural entitlement, with the narrowing made by the run'' (book Theorem~\ref{thm:ent-run-machine}).}
\begin{proof}
(1)--(2) Start with the prior list. At each CHECK the actual run passed,
run the preceding instructions from each surviving rival and ask that same
question there. Keep the rivals that pass. Nothing adds a rival, and the
actual start $s_0$ passes every one of its own tests, so it stays throughout.
A missing rival has a first test that removed it. That is the required
failed CHECK, at the stated place in the run.

(3) Name the successive list sizes $q_0,q_1,\ldots,q_n$, with $q_0=|\Omega|$
and $q_n=|\Omega'|$. The before-and-after sizes for test $i$ are consecutive
members of that list. Multiply the after-sizes by $q_0$, and multiply the
before-sizes by $q_n$: each side contains exactly
$q_0q_1\cdots q_n$. This is the displayed product identity.

(4) The share premise says $q_{i-1}\le 2^k q_i$ at every test. Apply it first
to $q_0$, then to $q_1$ in the resulting bound, and keep going. After $n$
tests it gives $q_0\le 2^{kn}q_n$. Numbering $q_n$ final places takes
$\lceil\log_2 q_n\rceil$ bits; another $kn$ bits can number up to $2^{kn}$
times as many. Thus the difference in rounded index lengths is at most
$kn$. Each of these $n$ tests is a CHECK the run actually passed, so it is
among the paid record moves. Exact prices make the number of those moves
the ledger increase. Multiplying that larger count by $k$ preserves the bound.

(5) The run starts clean and ends with its record up. The earned-record
requirement therefore finds its CHECK, COMMIT and raising CERTIFY. Those
three different acts each cost one, so the ledger increased by at least three.

(6) The actual start remained on every list, so every $q_i$ is positive.
The ratios $q_0/q_1,q_1/q_2,\ldots,q_{n-1}/q_n$ multiply to $q_0/q_n$:
each intermediate size cancels once above and once below. Taking base-two
logarithms turns this product into the sum of the per-test logarithms.

To account for rounding, set $u=\log_2 q_0$ and $w=\log_2 q_n$.
The integer $\lceil u-w\rceil+\lceil w\rceil$ is at least
$(u-w)+w=u$, so it is at least the least integer above $u$, namely
$\lceil u\rceil$. Rearranging gives the claimed bound by
$\lceil\log_2(q_0/q_n)\rceil$. Finally, rounding each per-test logarithm
up and adding gives an integer at least as large as their sum. It cannot
be smaller than rounding the sum up just once. This proves the last comparison.
\end{proof}

\paragraph{Proof of the proposition ``One CHECK can take four bits'' (book Proposition~\ref{prop:ent-share-needed}).}
\begin{proof}
$A \ge 15$ holds at exactly one of the sixteen starts. From $(15, 0)$ the CHECK passes and files the
claim, the COMMIT finds it with $A$'s version unchanged, and the CERTIFY raises the record. The
posterior is the one rival that passes the run's only CHECK, so the index bits are
$\lceil \log_2 16 \rceil - \lceil \log_2 1 \rceil = 4$, and the run has three record moves.
\end{proof}

\paragraph{Proof of the theorem ``What the record adds'' (book Theorem~\ref{thm:ent-observed}).}
\begin{proof}
(1) One direction is clause (b): a run from a clean start that ends with the record up contains the
chain. For the other, the chain's CERTIFY leaves the record up, which is part of what the chain
is, and by the last words of clause (a) a record that's up stays up under every move, so whatever the run does after that CERTIFY it ends
with the record up. The floor of 3 is Theorem~\ref{thm:u34-tc-consequences}(3). (2) Read the proof of
Theorem~\ref{thm:ent-machine} again and mark what each conclusion used. Its (1) used (M1) and (M2),
through Theorem~\ref{thm:ent-cs}(1), and nothing about the machine. Its (3) used (M5) and the count
$|\Omega| \le L(T) \cdot |\Omega'|$ of (M6), and (M4) to put the record moves in place of the
tree's depth; the clean start and the raised record never entered. Its (4) is clause (c) added up
along the run, from whatever state the run starts at. The decoded run being shown by some survivor
is in the hypotheses and, like the second half of (M3) above, feeds none of the three. The
conclusions dropped, (2) and (5), are exactly the two that used the clean start and the record.
\end{proof}

\paragraph{Proof of the proposition ``The eight-start member'' (book Proposition~\ref{prop:ent-small}).}
\begin{proof}
From any loaded start, the first CHECK passes exactly when $a$ is even.
If it passes, the second has room in the table and passes exactly when $a\ge2$.
Nothing has changed $A$, so their matching COMMIT and CERTIFY then succeed.
If either check fails, the trap prevents certification. This proves the
claim for every $a,b$. Within the eight rivals, $a$ can only be $0,1,2,3$;
of those, only $2$ is even and at least $2$, whatever $b$. From $(2, 0)$ and $(2, 1)$ the
COMMIT then finds its fact with $A$'s version unchanged, nothing having written $A$, and the CERTIFY
raises the flag; from every other start the first failing CHECK traps the machine, and a trapped
machine raises nothing. So the posterior is those two. The hypotheses: the posterior narrows the
eight strictly (M1); a dropped rival whose run fails a CHECK shows receipts no survivor shows (M2);
the run ends with the record up and its receipts, two passes, are the survivors' (M3); the tree has
depth 2 and the run 4 record moves (M4); two survivors (M5); and the fibres, the four rivals with
$B$ at 0 for the survivor $(2, 0)$ and the four with $B$ at 1 for $(2, 1)$, cover all eight, add up
to eight, and have four members each, the leaves of the tree, so $8 \le 4 \cdot 2$ (M6). Then
$\lceil \log_2 8 \rceil - \lceil \log_2 2 \rceil = 3 - 1 = 2$; the four record moves are two
CHECKs, a COMMIT and a CERTIFY, and the ledger rises by exactly those four by clause (c); and the
chain is in the run. For Theorem~\ref{thm:ent-questions}, the run's questions are its two CHECKs,
and two rivals answer them alike exactly when their $a$'s agree on being even and on reaching 2,
which puts the eight starts in four classes of two, so $m = 2$ and
$\lceil \log_2 8 \rceil - \lceil \log_2 2 \rceil = 2$ is the number of CHECKs. With the two
CHECKs alone there's no CERTIFY, so the record stays down, the ledger rises by 2, and the index
bits are still 2.
\end{proof}

\paragraph{Proof of the theorem ``A search through $n$ bits'' (book Theorem~\ref{thm:ent-search}).}
\begin{proof}
Each question is a CHECK of whether one register holds a given bit, which passes exactly when it
does; a failing CHECK traps. So the program certifies exactly when every answer is right, which is
when $v$ begins with the answers. The counting is the count of bit strings: the $2^m$ fibres of
$2^k$ values each add up to the $2^{k+m}$ of the prior, and none is longer than the $2^k$ leaves of
the tree, so the representation theorem applies with every hypothesis a theorem. The cap is the fact table: it holds sixteen facts,
never evicts, and a CHECK on a full table traps (Theorem~\ref{thm:u34-bounded}), so a seventeenth
CHECK before the CERTIFY ends the run in a trap.
\end{proof}

\paragraph{Proof of the theorem ``The time tax'' (book Theorem~\ref{thm:time-tax}).}
\begin{proof}
(1) Take the starts $(m - 1, b)$ and $(m, b)$. At the $j$th move, $j \le n < m$, at most $j - 1$
pebbles have left $A$, so $A$ is at least $m - j \ge 1$ in both runs, and no DEC finds $A$ empty
in either. The two runs take the same branches and differ only by the one pebble in $A$. They are
at the same line with the same $B$, and they stop at the same moment if they stop at all, though
$A \ge m$ holds of one start and not the other. (2) Run the ladder from $A = a$. Each of the first
$\min(a, m)$ moves takes a pebble off $A$ and jumps two lines on. If $a \ge m$ that's all $m$
moves, ending on line $2m + 1$. If $a < m$, the run reaches line $2a + 1$ after $a$ moves and from
there goes on one line per move (a DEC falling through, then an INC), so after $m$ moves it's on
line $a + m + 1$, below $2m + 1$. It's all DEC and INC, so it costs 0. (3) is the price
table: CHECK ($A \ge m$) passes exactly when $A \ge m$, a failing CHECK traps and the flag stays
down, and the three moves cost three. (4) is the same price table with one instruction: a passing
CHECK goes on to line 2 and a failing one traps and leaves the line at 1. (5) is arithmetic on
(1) to (4).
\end{proof}

\paragraph{The worked run for the time tax, with its table.}
Try $m = 3$ with $A = 2$. The ladder takes a pebble on line 1 and jumps to line 3, takes the second
and jumps to line 5, then finds $A$ empty and falls through to line 6, the INC: after three moves
it sits on line 6, not 7. With $A = 3$ the third DEC succeeds and the run is on line 7. Two
moves in, both runs are on line 5 with $B$ untouched, which is (1).

\begin{center}
\begin{tabular}{lcc}
\toprule
Moves taken & Start with two pebbles & Start with three pebbles\\
\midrule
0 & line 1, $A=2$ & line 1, $A=3$\\
1 & line 3, $A=1$ & line 3, $A=2$\\
2 & line 5, $A=0$ & line 5, $A=1$\\
3 & line 6, $A=0$ & line 7, $A=0$\\
\bottomrule
\end{tabular}

\small The reader in the theorem sees the line and $B$, which stays unchanged
through these three moves. The $A$ entries are shown here to explain the run;
that reader is not allowed to use them. Line 6 still has an instruction:
the decision time is not a claim that both runs have halted.
\end{center}


\section{Proofs, the program and the worked tables for ``The universal machines''}

\paragraph{Proof of the theorem ``No host fact on an exact copy of the guest's counter'' (book Theorem~\ref{thm:u34-nocopy}).}
\begin{proof}
A fixed program mentions finitely many properties, so $Q$ is finite, and the translation sends
the infinitely many $\mathsf{Ge}(0), \mathsf{Ge}(1), \ldots$ into it: two of them land on the
same member. The guest's COMMIT asks for $\mathsf{Ge}(m)$, and its table can only hold
$\mathsf{Ge}(n)$, so it traps (or its CHECK already trapped). The host sees one property twice:
its CHECK passes, since the guest's would, and its COMMIT finds exactly what that CHECK wrote.
\end{proof}

\paragraph{Proof of the theorem ``$U$ keeps pace with its guest'' (book Theorem~\ref{thm:u34-U-sim}).}
\begin{proof}
The one-step part is a case for each thing the guest can do at its line: stop; INC; DEC, taken or
not; CHECK, passing, failing, or with a full table; COMMIT, found or not; CERTIFY, passing or not.
The block invariants and the case analysis in Appendix~\ref{sec:universal-blocks}
show where each matching stretch lands. The relation carries
the bookkeeping that makes the next case go through: each count register equals the number of
slots in use for its counter, each slot in use holds (property number, checked value), a slot's
fact is current exactly when the matching guest fact is, a live slot's mirror is the property's
number plus one and a stale one's is 0, and unused slots, unused mirrors and the dead end are 0.
At loading, the counters and guest line have their specified starting values,
and both tables and channels are empty, so the matching conditions hold before
any step. Now take one guest instruction. The matching stretch just described
ends either at the next matching point or with both machines stopped. In the
first case its slot and mirror conditions are exactly what the following
instruction needs, so you can do the next stretch. Repeat this for the first
$m$ guest steps. Each running stretch takes at least one host step; if neither
machine has stopped, the accumulated host count is therefore at least $m$.
If they stop earlier, both remain stopped with the matching final fields.
That gives the whole-run statement for every $m$.
\end{proof}

\paragraph{Proof of the theorem ``Halting and answers match'' (book Theorem~\ref{thm:u34-U-halting}).}
\begin{proof}
Theorem~\ref{thm:u34-U-sim}: a halted guest is matched by a stopped host, and a host that has
stopped can't be back at line 1 with the guest still running, because line 1 isn't a stop. A
halted machine stays halted, so the two halting points carry the same values.
\end{proof}

\paragraph{Proof of the theorem ``The host's flag is the guest's, and its ledger is exact'' (book Theorem~\ref{thm:u34-U-flag}).}
\begin{proof}
(1) Flags are equal at matching points and never come down. (2) Ledgers are equal at matching
points; inside the stretch of $U$ that runs one guest instruction, at most the one paid site of
that instruction runs, so the host's ledger sits between the guest's before and after it.
\end{proof}

\paragraph{Proof of the theorem ``The host earns its own flag'' (book Theorem~\ref{thm:u34-U-earned}).}
\begin{proof}
By Theorem~\ref{thm:u34-U-flag} the guest's flag rises too; find the guest step where it rises, a
passing CERTIFY, and the related host point just before it. The guest's channel there names a
fact, and the relation says the host's channel names the matching slot fact. The host's own flag
rises at its own first passing CERTIFY, and nothing paid runs between the related point and that
CERTIFY, so the host's channel hasn't moved: it names that slot fact. The host machine's version
of Theorem~\ref{thm:u34-cert-earned} gives the host's CHECK and COMMIT on that slot. The record kept with the slot
says what pair it held and when, and the checker's meaning gives $\mathrm{holds}(p, v)$. The same
theorem on the guest's run, up to the step where its flag rose, gives the guest's chain on the
same $p$ and $c$, and the record says $c$ held $v$ at its CHECK.
\end{proof}

\paragraph{Proof of the theorem ``The host is certified Thiele-complete, and $U$'s halting is undecidable'' (book Theorem~\ref{thm:u34-U-complete}).}
\begin{proof}
(1) The base is registers 0 and 1, run by INC and DEC; the claims are ``$\mathsf{Slot}$ holds of
register $r$''; the rest is the host machine's own versions of the small machine's theorems, and
for non-vacuity, register 0 loaded with 1, which is $\langle 0, 0\rangle$, the pair of
$\mathsf{Zero}$ and 0, against register 0 loaded with 0, which is no pair. Over the claim
language the interface is the same one, with the meanings handed over by the language. (2) is
Theorem~\ref{thm:u34-tc-instances}(4) and Theorem~\ref{thm:u34-U-earned}. (3) is
Theorem~\ref{thm:u34-U-halting}, and the loading map reduces the small machine's halting question,
undecidable by Theorem~\ref{thm:u34-undecidable}, to this one.
\end{proof}

\paragraph{Proof of the theorem ``One host runs every computably presented Thiele machine'' (book Theorem~\ref{thm:u34-presented}).}
\begin{proof}
Two layers, each proved on its own. The bottom layer is the compiler: for every presented
machine, its guest program's run passes through points that code $M$'s states in order, with the
guest's ledger equal to $M$'s ledger plus the surcharge and the guest's flag equal to $M$'s latch,
and the guest's own flag stands on an earned chain whose claim is $M$'s reading program applied to
the coded state where the reading first said yes. Here is why. The guest keeps a few registers
packed into its counters the way Minsky packed them: the coded state, a latch bit, the price
still owed, and a register that counts the reading's steps. It starts at the reading with
nothing owed, and goes round one loop. Run the compiled reading on the coded state, counting its
steps. If it says yes and the latch bit is 0, run CHECK, COMMIT and CERTIFY on the claim
``reading program $r$ says yes of $v$,'' $v$ being the packed registers with the step count as
fuel, set the latch bit, and take three off what's owed, stopping at zero. Pay what's owed in
PAYs. Run the compiled driver, and halt if it says stop. Otherwise run the compiled step and
price for the move it picked, put the new coded state and that price in place, and go round
again. Everything but the PAYs and the one chain is INC and DEC, at no cost. So at the point just
after the PAYs and before the driver runs, once $n$ moves have been made, the coded state is $M$'s state after $n$ moves, the latch bit is $M$'s
latch, and the ledger is $M$'s ledger, except that the move where the latch bit was set was paid
with the chain's three and the larger of $c - 3$ and 0 in PAYs instead of its price $c$ (or, if the reading is yes at $s_0$, the start was charged the chain's three). The
CHECK passes, because the fuel is the reading's own step count, and at such a value the fixed
checker says what the reading says; the COMMIT finds that fact, since nothing wrote $B$ in
between; and the latch bit lets the chain run once. The top layer is $U_P$, proved the way $U$ was
in the previous theorems, with the PAY case added: it keeps pace with any guest of this kind,
halting and flag matched, ledger equal at matching points, and its flag earned on a slot. Compose
the two. The earned part takes one more step: the host's chain gives a guest CHECK that passed, and
a guest run that ends with its flag up has exactly one passing CHECK, the compiler's CHECK of the
reading on counter $B$. So the guest CHECK the host's chain points to is that one, and its claim is $M$'s
reading at the first state where it said yes.
\end{proof}

\paragraph{Proof of the theorem ``Exact, within two, and no better'' (book Theorem~\ref{thm:u34-presented-surcharge}).}
\begin{proof}
(1) and (2): the guest pays $M$'s prices in PAYs, a unit at a time, except at the move where
$M$'s reading first says yes. There it has to run its CHECK, COMMIT and CERTIFY, three units, and
it makes up the rest of the move's price in PAYs, so it pays the larger of the price $c$ and
three. The price is at least one, because the move raised a reading (the toll), so the excess is
at most two, and none when $c \ge 3$.
(3): a raised flag on the host machine costs at least three from a clean start, whatever the
program, and $M$'s ledger is less.
(4): the guest finds the reading yes before $M$ has made a move, runs its CHECK, COMMIT and
CERTIFY there, three units on top of a ledger of nothing, and from then on its flag is up, so the
reading never first says yes again and every move's price is paid in PAYs, exactly.
\end{proof}

\paragraph{Proof of the theorem ``The host runs any guest'' (book Theorem~\ref{thm:u34-universal}).}
\begin{proof}
(1) Before any move, the guest is at $g(h)$. One untrapped guest move
has budget 1, enough for every guest instruction, so it takes exactly the
guest's next step. It leaves the host untrapped, which lets the same argument
apply to the next move. After the first move the guest has taken one step;
after the second, two; and after $n$, exactly $n$. No such move changes the
guest's program. (2) One pass of $U_0$ is a guest step, then
INC $A$ and DEC $A$ back to line 1, so it is (1) once per pass. (3) A host program made only
of $\mathrm{OWN}$ moves is the small machine, core for core, and
Theorem~\ref{thm:u34-halting} does the rest.
\end{proof}

\paragraph{Proof of the observation ``The copy is the record'' (book Observation~\ref{thm:u34-mirror}).}
\begin{proof}
One move at a time. An $\mathrm{OWN}$ move touches neither the mirror nor the guest. A trapped
$\mathrm{GSTEP}$ touches neither. An untrapped one sets the mirror to the old mirror or the
crossing, and the guest's flag is a latch (Theorem~\ref{thm:u34-toll}), so if the old mirror
equalled the old flag, the new mirror equals the new flag. The last sentence is this
together with Theorem~\ref{thm:u34-universal}.
\end{proof}

\paragraph{Proof of the theorem ``The host's step collects the guest's toll'' (book Theorem~\ref{thm:u34-host-toll}).}
\begin{proof}
(1) The guest's flag rises only on a passing CERTIFY (Theorem~\ref{thm:u34-toll}), which costs
1, and the budget let it run, so $b \ge 1$; the mirror was the guest's flag, so it was down,
and the move sets it to the crossing. (2) is (1) at the matching move. (3) Write the host's charge and the guest's ledger increase beside every
move. When a guest step runs, its price is at most the budget the host paid.
When none runs, the guest's increase is zero and the host's charge is
nonnegative. Adding these comparisons along the trace gives
\[
 \mu(g(\mathrm{run}(t,h)))-\mu(g(h))
 \le \mu(\mathrm{run}(t,h))-\mu(h).
\]
Move the two starting readings to the other sides to get the stated form.
\end{proof}

\paragraph{Proof of the theorem ``No unpaid host certification'' (book Theorem~\ref{thm:u34-no-free-host}).}
\begin{proof}
An $\mathrm{OWN}$ move copies the mirror unchanged, and a trapped $\mathrm{GSTEP}$ does too, so
the rise is an untrapped $\mathrm{GSTEP}\ b$ whose guest step crossed; a crossing is a passing
CERTIFY, which costs 1 and needed the budget. The own flag moves by the small machine's step line.
\end{proof}

\paragraph{Proof of the theorem ``The host is a Thiele machine'' (book Theorem~\ref{thm:u34-host-instance}).}
\begin{proof}
The mirror's half is Theorem~\ref{thm:u34-no-free-host}; the own flag's half is the small
machine's toll on the host's core. The trace form is Section~\ref{sec:toll}'s universal floor.
\end{proof}

\paragraph{Proof of the theorem ``Rice's theorem on the host machine'' (book Theorem~\ref{thm:host-rice}).}
\begin{proof}
The reduction above. The prefix before $y$ leaves every register but the two fresh ones as it
found them, and empties those, so when $H$ halts the built program ends in whatever state $y$ ends
in from the same input, with the same ledger, table, channel, trap and flag. When $H$ doesn't halt,
the built program stops on no input. You don't even need to know whether the looping program has
$\Pi$: ask the decider about it, and its answer says which of $y$ and $n$ to put at the end.
\end{proof}

\paragraph{Proof of the theorem ``Rice's theorem for the small machine, from the clean start'' (book Theorem~\ref{thm:guest-rice}).}
\begin{proof}
The road of Theorem~\ref{thm:host-rice}, with the input dropped. Take a two-counter program $H$ with
a start $(a, b)$, and let $y$ have $\Pi$ and $n$ lack it; as before, if the program that loops has
$\Pi$, swap the roles of $y$ and $n$. Build the program that does INC $A$ $a$ times and INC $B$ $b$
times, runs a copy of $H$ with its jumps moved down (a jump to a line where $H$ has no instruction, line 0 or past its end, goes to the line
just after the copy), empties $A$ and then $B$ with a one-line DEC loop
each, and then runs $y$ with its jumps moved down. From the clean start the copy of $H$ begins at
$(a, b)$, which is $H$'s own start, so it stops exactly when $H$ does. If it does, both counters are
emptied and $y$ runs from $(0, 0)$. Every instruction before $y$ is an INC or a DEC: it adds
nothing to the ledger and touches no fact, channel, trap latch or flag. It does move the versions of $A$ and $B$ up, each by
the number of writes that counter took before $y$ started, and the facts $y$ records carry versions.
But a CHECK or a COMMIT reads versions only to ask whether two versions of the same counter are
equal, and adding one fixed amount to both keeps the answer. So $y$ takes the same steps inside the
built program as it takes from the clean start, with the same counters, trap latch, ledger and flag
at every point, and with facts that differ only by that fixed shift in their versions. That shift is
why the fact table and the channel are left out of ``behaves the same from the clean start.'' If $H$ doesn't stop,
neither does the built program: the copy of $H$ never reaches a HALT, and nothing before $y$ can
trap. So the built program behaves the same as $y$ when $H$ stops from $(a, b)$, and the same as the
two-line program that loops, INC $A$ then DEC $A$ back to line 1, when it doesn't. A decider for
$\Pi$ would then decide whether $H$ stops from $(a, b)$, with the consequence stated in
Theorem~\ref{thm:u34-undecidable}; which of $y$ and $n$ goes at the end is the decider's own
answer about the looping program. ``Halts
from the clean start'' is carried by behaving the same from the clean start, held by the empty
program and lacked by the looping one, so it's covered.
\end{proof}

\paragraph{Proof of the theorem ``s-m-n, and a universal program for host programs'' (book Theorem~\ref{thm:host-smn}).}
\begin{proof}
(1) The prefix writes register 2 $c$ times and nothing else, so it moves that register's version from
0 to $c$. A fact records a version, so after the prefix every later version of register 2, and the
version in every fact the run records about register 2, is $c$ more than in a clean run. A CHECK or a
COMMIT reads versions only to ask whether two versions of the same register are equal, and adding $c$
to both keeps the answer, so the two runs take the same steps. The run is $V$'s run from the stated
start, moved down by $c$ lines. (2) is the evaluator. A finite algorithm decodes a number into a program, runs it step by
step on given registers, record layer included, and, searching over the number of steps,
returns the output of the first halting run. It is a finite tape algorithm, and the translation
of Appendix~\ref{sec:counter-constructions} turns any finite tape algorithm into a finite program
of INC and DEC on finitely many counters. The host runs those instructions on registers of its
own, with HALT marking the returned output. That's $W$.
\end{proof}

\paragraph{Proof of the theorem ``Kleene's recursion theorem on the host machine'' (book Theorem~\ref{thm:host-kleene}).}
\begin{proof}
(1) Let $E$ be the following program of INC and DEC, built the same way as $W$ from one finite
evaluator, with $T$'s number $t$ written in. On inputs $x$ and $c$ it computes the number of $\mathrm{spec}(c, P_c)$, where $P_c$ is the
program numbered $c$, runs $T$ on that number to get $d$, and then runs the program numbered $d$ on
$x$. Let $c_0$ be the number of $E$, and let $e = \mathrm{spec}(c_0, E)$. By (1) of
Theorem~\ref{thm:host-smn}, $e$ on input $x$ does what $E$ does on $x$ and $c_0$. And $E$ on
$c_0$ builds the number of $\mathrm{spec}(c_0, P_{c_0})$, which is the number of $e$, since $P_{c_0}$
is $E$. So $e$ hands its own number to $T$, gets $d$, and runs $d$ on $x$. (2) and (3) are (1) with $T$
carrying out $F$ or $f$.
\end{proof}

\paragraph{Proof of the theorem ``The recursion theorem, up to the record'' (book Theorem~\ref{thm:host-kleene-record}).}
\begin{proof}
The construction is Theorem~\ref{thm:host-kleene}'s, with a different evaluator. The first step is to
see that a record is small. From a clean start, a host program that stops has a record I can read off
five numbers, in this order: what it left in register 0; how many facts the table holds; what is
left of the ledger once the facts, the trap and the flag are paid for, which I'll call the replay allowance;
1 if the flag is up and 0 if not; and 1 if the trap latch is up and 0 if not. The record obeys a
short rulebook, and you can read each rule off the price table. A run stops at its first trap, so
its ledger is its passing CHECKs, plus its passing COMMITs, plus its passing CERTIFYs, plus one if
it trapped (only those three can fail, and the failing one is charged). Its passing CHECKs are its
facts, at most sixteen. Its flag is up exactly when some CERTIFY passed. So the allowance is the
passing COMMITs plus the passing CERTIFYs less one for the flag, which is never negative, and the
ledger is the facts, plus that allowance, plus the flag, plus the trap. A COMMIT needs a fact
behind it, so a positive allowance means at least one fact. A CERTIFY needs a commitment, so a
raised flag means a passing COMMIT, and so a positive allowance. And the channel is empty exactly
when no COMMIT passed, which is exactly when the allowance is zero, since a CERTIFY can't pass
before a COMMIT has. The allowance counts repeated CERTIFYs as well as COMMITs; it counts the cost the
replay can spend on commitments, not the original commitment events. Those are the only properties
of a record the replay will use.

The second step is five blocks, one for each number. Each block reads the code of $T$, runs $T$ on the
number of the specialised program, runs the program that comes out on $x$, and returns its number. All of
that happens on plain data: the blocks are programs of INC and DEC, built the way $W$ was, from a
finite evaluator translated into INC and DEC. So if $T$ has marks
of its own, they are simulated and never paid.

The third step is the payment. The evaluator puts 1 in a fresh register, which is
$\langle 0, 0 \rangle$, the pair of $\mathsf{Zero}$ and 0, so it satisfies the host's one property,
and CHECKs it once for every fact; there are at most sixteen, so the table has room. It COMMITs to that claim once for every unit in the replay allowance. It CERTIFYs
once if the flag is up. It copies the first number into register 0. And if the trap latch is up, it
CHECKs a register that holds 0, which fails, since 0 is no pair, and the host traps. Everything
before the payments is INC and DEC and costs nothing, and the payments cost what the record costs, so the ledger is exact. Specialise the
evaluator to its own number, as before.
\end{proof}

\paragraph{Proof of the theorem ``No fixed point of full behaviour'' (book Theorem~\ref{thm:host-no-exact}).}
\begin{proof}
Let $F(p)$ be the one-line program INC $r$, with $r$ one above the number of $p$. To carry this out on numbers, add one to the input code, encode INC with that register,
and encode the one-element instruction list. These are finite arithmetic operations of the
counter construction in Appendix~\ref{sec:counter-constructions}. A program's number is at
least every register it names: the number of an instruction is above its register, and the code of a
list is above every instruction in it. So $e$ can't name $r$, the register $F(e)$ writes. On input 0 the
run of $F(e)$ ends with 1 in register $r$, and if $e$ stops at all, it ends with 0 there. If $e$ doesn't
stop, then it already differs from $F(e)$, which does. Either way they don't behave the same.
\end{proof}

\paragraph{Proof of the theorem ``The diagonal on the host machine'' (book Theorem~\ref{thm:host-wall}).}
\begin{proof}
(1) The flip, send programs that $d$ calls good to the program lacking $\Pi$ and the others to the
one having it, is carried out by a host program by hypothesis. Theorem~\ref{thm:host-kleene}(2) gives
a program $e$ that computes what its image computes; $\Pi$ holds of $e$ exactly when it holds of the
image, and the image lacks it exactly when $d$ calls $e$ good, so $d$ is wrong at $e$. (2) From
a deciding program $D$ build a host program that computes $D$'s flip on numbers, by the finite-simulation search of Appendix~\ref{sec:counter-constructions}, reading the
Boolean from the first halting run of $D$; then (1) applies. This one is a plain no: the flip (1) needs is built from $D$ itself, so no premise about what can be represented is left over.
\end{proof}

\paragraph{Proof of the theorem ``Rice's theorem for the two-counter machine, on any inputs'' (book Theorem~\ref{thm:tc-rice}).}
\begin{proof}
For every input: given a two-counter instance $H$ starting from counters $a$ and $b$, build a program
that turns the input of $n$ pebbles in bin A into $(6n + 5) \cdot 2^a \cdot 3^b$, runs $H$ inside the
exponents of 2 and 3, divides out every 2 and every 3 when $H$ stops, takes away 5, divides by 6, and
then runs the program that has the property. Because $6n + 5$ shares no factor with 2 or 3, the coding
never mistakes it for an exponent, and the input comes back whole. If $H$ halts, the built program
behaves the same as the one with the property on every input; if not, it never stops on any input,
like the program that loops. As before, if the looping program has the property, put a program
that lacks it at the end instead, and the decider's own answer about the looping program says
which. So a decider for the property would decide two-counter halting.
\end{proof}

\paragraph{Proof of the theorem ``Kleene for the two-counter machine, inputs written as powers of two'' (book Theorem~\ref{thm:tc-kleene}).}
\begin{proof}
(1) The host's three pieces, with two counters in place of registers. First, the universal program.
Load a program's number $e$ and its input $x$ together as $2^x \cdot 3^e$ in bin A, with bin B empty.
The evaluator of Appendix~\ref{sec:counter-constructions}, the one that searches budgets and runs the
program numbered $e$ on $x$, is a finite counter program on finitely many counters, and packing those
counters into two by the prime-exponent construction of Appendix~\ref{sub:prime-packing} gives
$U_2$: started on $2^x \cdot 3^e$ it stops with $2^y$ in bin A exactly when the program numbered $e$,
started on $2^x$, stops with $2^y$. Second, specialising. For a number $c$ and a program $V$, let
$\mathrm{spec}(c, V)$ be $c$ copies of a block that multiplies bin A by three, followed by $V$ with
its jumps moved down. The block is INC and DEC only, so it adds nothing to the ledger and records nothing, and
started on $2^x$ the $c$ copies leave $2^x \cdot 3^c$; so $\mathrm{spec}(c, V)$ on $2^x$ does what $V$
does on $2^x \cdot 3^c$. Third, the fixed point. Let $T$ be a program that carries out $F$: on $2^n$,
$n$ the number of a program $p$, it stops with $2^m$, $m$ the number of $F(p)$. Build $E$ from the
same finite evaluator, with $T$'s number written in, and pack it into two counters as before: on
$2^x \cdot 3^c$ it computes the number of $\mathrm{spec}(c, P_c)$, where $P_c$ is the program numbered
$c$, runs $T$ on that number to get $d$, and runs the program numbered $d$ on $2^x$. Let $c_0$ be the
number of $E$ and $e = \mathrm{spec}(c_0, E)$. On $2^x$, $e$ does what $E$ does on
$2^x \cdot 3^{c_0}$: it builds the number of $\mathrm{spec}(c_0, P_{c_0})$, which is the number of $e$
itself since $P_{c_0}$ is $E$, hands it to $T$, gets the number of $F(e)$, and runs $F(e)$ on $2^x$.
The machine is deterministic, so a packed output is unique, and for every $x$ and $y$, $e$ started
on $2^x$ stops with $2^y$ exactly when $F(e)$ does. (2) is the counterexample that follows.
\end{proof}

\paragraph{Proof of the theorem ``No plain recursion on two counters'' (book Theorem~\ref{thm:no-plain-recursion}).}
\begin{proof}
(2) from (1): let $F(p)$ be a program that multiplies bin A by the constant
$c$ determined by the length of $p$. Multiplication by a fixed constant is the
transfer loop of Appendix~\ref{sec:counter-constructions}. Decoding a finite
program, counting its lines, finding the least prime above a number and writing that loop are
computable operations, so a program carries out $F$ on encoded program numbers.
A fixed point would multiply every plain input by the constant forbidden by (1).

For (1), suppose an $n$-line program did it. Split each state into the two bin counts and a
finite \emph{control}: the line (one of the $n$ lines, or ``stopped'' for any line with no
instruction), the trap, whether the channel is empty, and the table read entry by entry as a
fact's property, its counter, and whether it is \emph{current} (its version is its counter's
version now) or stale. Only properties the program names can enter the table, at most $n$ of
them, so an entry is one of at most $4n$ kinds, and a table of at most sixteen entries is one of
at most $\sum_{k \le 16} (4n)^k \le (4n+1)^{16}$ lists. So there are at most
$Q = 4(n+1)(4n+1)^{16}$ controls. The control and the two counts decide the next control and
counts: a CHECK writes a current fact, a write to a counter makes that counter's facts stale, a
COMMIT passes exactly when a current fact with its property and counter is in the table, a
CERTIFY passes exactly when the channel isn't empty, and the ledger, the flag, the channel's
contents and the versions themselves never decide a branch. Call the control with the two counts
a \emph{configuration}.

Pick a threshold $h \ge 1$ above every number $k$ in a property $\mathsf{Ge}(k)$ the program
names. At $h$ or above, a count answers every question the program can put to it (is it zero, is
it even, is it at least $k$, can a DEC take one off) the same way as any other count of the same
parity at $h$ or above. Call a control together with the parities of both counts a
\emph{pattern}; there are at most $4Q = N$ of them.

\emph{The forced run.} A \emph{stage start} is a configuration with one count below $h$, in the
\emph{exact} bin, and the other at $h$ or above, in the \emph{large} bin, which holds $v$. From
it, run the program while pretending the large bin can never fall below $h$: follow the exact bin
exactly, the large bin only by its parity, and keep a tally of the large bin's net change. This
forced run depends on the control, the exact count and the parity of $v$, and not on $v$ itself.
The real run does exactly what the forced run does, with the large bin at $v$ plus the tally, up
to the first step where $v$ plus the tally drops below $h$, the \emph{crossing}. A move changes a
count by at most one, so at the crossing the large bin holds exactly $h - 1$.

\emph{What a forced run does.} One of three things, fixed by the stage start. (a) It stops. (b)
It comes back to a step with the same control, the same exact count below $h$ and the same
parity of the large bin as an earlier step; it's deterministic in those, so it repeats that
stretch forever, and the exact count stays bounded. (c) Neither. Then it has at most $2Qh$ steps
with the exact count below $h$, one per combination, so from some step on the exact count stays
at $h$ or above for good. Among any $N + 1$ consecutive steps of that tail, two have the same
pattern, at steps $t_1 < t_2 \le t_1 + N$, and between them the exact count changes by some
$\delta$, even because the parities match. If $\delta < 0$, the run from $t_2$ repeats the run
from $t_1$ with the exact count lowered by $|\delta|$ each time round, so after finitely many
rounds it goes below $h$, which the tail never does. So $\delta \ge 0$, and the stretch repeats
forever, lifted by $\delta$ each time round, never dipping below $h$. In every case but (a),
then, the forced run from some step on repeats one stretch forever, and each time round the exact
bin changes by a fixed $g$ and the large bin by a fixed amount; in (b) $g = 0$, and in (c) the
stretch is at most $N$ moves long.

\emph{What the real run does,} once $v$ is large enough for that stage start, meaning the large
bin can't reach the crossing before the forced run stops or its repeating stretch has gone round once in full. If the
forced run stops, the real run stops the same way, with the exact bin at a count fixed by the
stage start and the large one at $v$ plus a fixed amount. If the large bin doesn't lose each time
round, the real run never crosses and never stops. If it loses and $g = 0$, the real run crosses
with both bins below a bound fixed by the stage start. The remaining case is a \emph{transfer}:
each time round the large bin loses $m > 0$ and the exact one gains $g > 0$. That's case (c), so
the stretch has at most $N$ moves, each moving one bin by at most one, and $m, g \le N$; both are
even. Now the stage rule. Start the same stage with $v + m$ in place of $v$. The parity is the
same, so the forced run is the same, and since the crossing comes after the stretch has gone
round once, the bin needs exactly one more time round to get down to $h - 1$. So the real run
crosses one time round later, with the same control, the formerly large bin at $h - 1$, and the
other bin larger by $g$. With $v + km$ it crosses $k$ times round later, the other bin larger by
$kg$. After a transfer, the crossing is a new stage start with the bins' jobs swapped. For
example, a loop that takes two pebbles from the full bin and puts four in the other has $m = 2$
and $g = 4$: two more pebbles going in buy one more trip round the loop, and four more coming
out.

\emph{Choosing the input.} Apart from $v$, a stage start is one of finitely many things: a
control, an exact count below $h$, a parity, and which bin is which. Take one bound $V \ge h$
that's large enough, in the sense above, for every one of them, and above every count that was
called fixed or bounded above: the exact bin of a stage that stops, and both bins at a crossing
with $g = 0$. The configurations with one count below $h$ and the other below $V$ are finitely
many. If the program multiplied every input by $c$, no two different inputs could pass through
the same configuration, because everything after it, the output included, would be the same,
while $cx \ne cx'$. So only finitely many inputs pass through that finite set. Pick an input $x$
that passes through none of it and has $cx$ above $V$. Since $(x, 0)$ is itself a configuration,
$x \ge V$.

\emph{Following that input.} The run from $(x, 0)$ starts a stage: bin B exact at 0, bin A large
at $x$. Every later stage start along the run has its large bin at $V$ or above, or the run would
pass through the finite set. The run stops, since the program multiplies, so no stage runs
forever. No stage crosses with $g = 0$, since that crossing is in the finite set. So each stage
before the last is a transfer, with its own pair $(m_i, g_i)$, and the last one stops. It stops
with its exact bin below $V$ and its large bin at its $v$ plus a fixed amount, and the output $cx$
is above $V$, so the output bin A is the large one. There are finitely many stages, because each
takes at least one step of a run that stops.

\emph{The comparison.} Put $D = 2\prod_i m_i$ and run the input $x + D$. It is $x$ plus a
multiple of $m_1$, so by the stage rule the first stage crosses at the same configuration except
that the other bin is larger by $(D / m_1) g_1 = 2 g_1 \prod_{i \ge 2} m_i$, which is a multiple of
$m_2$. So the second stage crosses the same way, with the next bin larger by
$2 g_1 g_2 \prod_{i \ge 3} m_i$, and so on through every stage. The last stage starts with its
large bin larger by $2\prod_i g_i$, an even amount, so it stops the same way with bin A larger by
exactly that. So the outputs from $x + D$ and from $x$ differ by $2\prod_i g_i$. But
multiplication by $c$ makes them differ by $cD = 2c\prod_i m_i$. Cancel two:
\[
c\prod_i m_i=\prod_i g_i.
\]
Empty products are 1. Every factor on the right is between 1 and $N$, and
$c$ is a prime above $N$, so it divides none of them, and a prime that divides a
product divides one of its factors. So $c$ can't divide the right side. That's the
contradiction.
\end{proof}

\paragraph{One slot, done by hand.}
Do one slot by hand. The property ``even'' has number 1. Checking it on value 4
puts $\langle1,4\rangle=(2\cdot4+1)2^1=18$ in a separate slot. To read that
pair, halve 18 until what's left is odd: one halving gives 9, so the
property number is 1; subtract one from 9 and halve to recover value 4. The host
checks that property of that value. A second register, the mirror, holds 2: the property
number plus one. COMMIT can read this second register while leaving the checked slot
untouched.

When the guest changes the counter that fact concerned, nudge the slot
$18\to19\to18$ and set the mirror to zero. The number is back where it was, but the version
went up twice, so the host's old fact about that slot is stale. Putting a number back
doesn't put an old check back. The mirror is what tells the search to leave that slot alone,
and if the search tried it anyway, the version rule would refuse it.

\begin{center}
\small
\begin{tabular}{@{}lccc@{}}
Act & slot value & slot edition & searchable mirror \\ \hline
Check evenness of encoded 4 & 18 & $v$ & 2 \\
Guest writes its counter: INC slot & 19 & $v+1$ & \\
DEC slot; clear mirror & 18 & $v+2$ & 0 \\ \hline
\end{tabular}
\end{center}

\paragraph{The program, register by register and job by job.}
\textbf{The program.} I'll describe $U$ the way Turing would, one register at a time and then
one job at a time. Registers 0 and 1 hold the guest's $A$ and $B$. Register 2 holds the guest's
whole program as one number: each instruction gets a number (an opcode paired with its operands,
INC as 0, DEC as 1, HALT as 2, CHECK as 3, COMMIT as 4, CERTIFY as 5), and the list of them is
packed with $\langle \cdot, \cdot \rangle$, the empty list as 0. Register 3 is the guest's line.
Registers 4 to 13 are scratch. Registers 14 and 15 count how many CHECKs on $A$ and on $B$ have
passed. Registers 16 to 47 are \emph{mirrors}, sixteen for each guest counter, and registers 48 to
79 are the slots, sixteen for each guest counter. The guest's table holds sixteen facts in
total, across both counters, and all sixteen can be about one counter, so the host gives each
counter room for all of them. The reservation is the host's, and the guest's own capacity stays
sixteen in all. Register
80 is a dead end that's never written.

Each guest step goes like this. At line 1 the program copies the program number and the guest's
line into scratch, putting the originals back as it goes, because copying a counter on this
machine means emptying it into two others and refilling it. If the guest's line is 0, or past
the end of the program, it stops. Otherwise it peels instructions off the front of the copy, one
per line, until it holds the one at the guest's line, and splits that into opcode and operands.
Then a chain of DECs on the opcode sends it to one of six handlers.

\begin{itemize}
\item INC $c$: INC the guest's counter $c$, nudge every slot about $c$ and zero its mirrors, add
one to the guest's line, back to line 1.
\item DEC $c$ $j$: DEC the guest's counter $c$. If it was 0, add one to the guest's line. If it
wasn't, nudge every slot about $c$, zero its mirrors, and set the guest's line to $j$. Back to
line 1.
\item HALT: stop.
\item CHECK $p$ $c$: let $k$ be the count of passed CHECKs on $c$. If $k$ is below 16, build
$\langle$number of $p$, value of $c\rangle$ in slot $k$ of $c$, run CHECK $\mathsf{Slot}$ on that
slot, set its mirror to the number of $p$ plus one, add one to the count and to the guest's line,
and go back to line 1. If $k$ is 16, run CHECK $\mathsf{Slot}$ on the dead end, which holds 0,
which is no pair, so it traps, just as the guest's seventeenth CHECK does.
\item COMMIT $p$ $c$: look through the sixteen mirrors of $c$ for one equal to the number of $p$
plus one. At the first match, run COMMIT $\mathsf{Slot}$ on that slot. With no match, run COMMIT
$\mathsf{Slot}$ on the dead end, which traps.
\item CERTIFY: run CERTIFY, add one to the guest's line, back to line 1.
\end{itemize}

Why the mirrors? COMMIT has to find the slot holding the claim, and it can't read the slot to
find out, because reading would make the fact stale. So beside every slot sits a mirror the
program is allowed to read: the property's number plus one while the slot's fact is live, 0 once
the guest has written the counter. If a guest CHECK fails its property, the host's CHECK on the slot fails too,
since the slot holds exactly the pair the guest checked. If the guest table is
full across both counters, the host's own sixteen-entry table is full as well,
even when this counter's slot count is below sixteen; its CHECK then traps for
lack of room. So both kinds of refusal come through.

\paragraph{The two tally lines for the surcharge.}
Keep two tally lines for a guest whose first three moves each cost one and whose
reading first says yes after the third. The guest's totals are 1, 2, 3. This
compiler pays one PAY on each of the first two moves, then its three-act chain:
the host totals are 1, 2, 5. The excess is two, added once. If the third move
instead costs four, the host runs the chain and one PAY there; both totals then
read 6. These are counts of charged instructions, not counts of all the work needed to decode the guest. Stopping at zero throws away what the chain paid beyond the move's price. A compiler that kept it as credit against the moves after the raise would bring the surcharge back to zero once the guest had paid three from the raising move on, and a guest whose reading starts at yes would be matched exactly from its third unit of cost on. I keep the simpler rule, and every surcharge stated below is this compiler's.


\paragraph{A worked replay, with its table.}
Try CHECK, COMMIT, CERTIFY, CERTIFY on a true claim. The ledger is four, the table holds one
fact, the flag is up and the trap is down, so the replay allowance is two. The replay pays
CHECK, COMMIT, COMMIT, CERTIFY: the same ledger of four and the same summary, in a different order.
The theorem claims that agreement and no more.

\begin{center}
\small
\begin{tabular}{@{}lccccc@{}}
Run & checks & commits & certificates & total paid & final flag \\ \hline
Original & 1 & 1 & 2 & 4 & up \\
Replay & 1 & 2 & 1 & 4 & up \\ \hline
\end{tabular}
\end{center}

Each table has one fact and each channel is nonempty. The counts of the
individual acts differ; the selected final summary agrees. In the replay, the trap latch is up exactly when the fifth number is 1, and the flag
exactly when the fourth is positive.

\paragraph{A worked packing, with its table.}
Try input 2 while the simulated counters hold 1 and 2. The pile is
$17\cdot2\cdot9=306$. Raising the first simulated counter doubles it to 612; lowering
the second divides it by three, to 204. Strip off every factor of two and three and 17
remains; subtract five and divide by six, and the original input 2 is back.

\begin{center}
\small
\begin{tabular}{@{}lccc@{}}
Act & packed pile & simulated counters & input factor \\ \hline
Load input 2; counters 1, 2 & $306=17\cdot2\cdot3^2$ & $(1,2)$ & 17 \\
Increment first counter & $612=17\cdot2^2\cdot3^2$ & $(2,2)$ & 17 \\
Decrement second counter & $204=17\cdot2^2\cdot3$ & $(2,1)$ & 17 \\
Strip factors 2 and 3 & 17 & discarded & 17 \\
Subtract 5; divide by 6 & 2 & & input recovered \\ \hline
\end{tabular}
\end{center}


\section{Reading the instruction rules}

The instruction rules settle some concrete questions about running this machine. Follow
them through a check, a zero decrement and a trapped trace.

\begin{itemize}
\item \textbf{The check price is flat.} CHECK costs 1 whether it's asking ``is this 0'' or ``is this at
least a billion.'' The property language is fixed and each property is one evaluation, so
``pays for its work'' here means that flat price, the same however hard the
question was. A machine that charged by the difficulty of the check would be a different
machine, and I haven't built it.
\item \textbf{DEC on zero isn't a write.} It doesn't change the counter, so it doesn't change the
version. That's the right call for versions (nothing changed), but it means DEC on zero is
a way to move the line number that adds nothing to the ledger and leaves no trace in the record. The idle run in the
subsection on the window, earlier in this chapter, uses exactly that.
\item \textbf{After a trap, a trace still pays.} On programs, a trap stops the run. On arbitrary
traces, instructions after a trap are each charged their cost and none of them changes the
core. That isn't a separate choice. A price belongs to the instruction, the same from every state
(Definition~\ref{def:certsys}), so a CERTIFY on a trapped machine costs what every CERTIFY costs, and
clause (c) of Definition~\ref{def:u34-thiele-complete} has the ledger add it; Theorem~\ref{thm:u34-toll}
has no exceptions for that reason. The choice is one level up, in putting the price on the instruction and
not on the instruction at a state, and what it brings with it is that a trapped
machine is still charged for instructions that change nothing.
\item \textbf{``Checked'' means what the property language can say.} On this machine that's \emph{is
zero}, \emph{is even} and \emph{is at least $n$}, and a live fact means one of those holds of one counter's
current value. It doesn't mean the program is correct, or that anything follows from the
property, or that the property was worth checking. So can a proved machine of this kind check
``sorted,'' the claim the sorter-and-copier table of Section~\ref{sec:established} runs on? It can, through an encoding. A copy of
this machine, instruction for instruction, with its property language left open, has provenance,
checker soundness, no forging and the floor of three proved once for every property language that
comes with an exact equality test and a checker proved equal to its meaning. Give it ``the list
this counter encodes is sorted'' as a property and it certifies a sorted list for three and
refuses an unsorted one forever. The instruction rules are those already given, with the chosen property language in place
of the three original properties. For sorted lists, decode 0 as the empty list and
$2^h(2t+1)$ as the list with head $h$ and tail decoded from $t$, the list code of
Theorem~\ref{thm:u34-tc-instances}.
Every positive code has $t$ smaller than itself, so decoding terminates. The checker
returns yes on an empty list or singleton, and on $a,b,\ldots$ exactly when $a\le b$
and it returns yes on $b,\ldots$. Induction on the list proves that this is exactly
adjacent nondecreasing order. The provenance proof uses only equality of properties, and soundness only that the checker's yes means the property holds, so the same instruction-by-instruction argument applies. What makes that a check of
the list is the proof that the checker means sortedness. Without that proof, ``sorted'' would be a number
some unconstrained code wrote into a counter, and the record would be honest about the counter
and silent about the list.
\item \textbf{The certificate is historical.} Between the COMMIT and the CERTIFY, the counter can
change, and the flag still goes up. The flag says a claim was checked and committed about
an edition. It says nothing about whether the claim is still true.
\item \textbf{One latch, and a channel that moves on.} The flag records that an earned chain
happened, once. After that it stays up, and later CERTIFYs change nothing. The channel is the
part that names a fact: with the flag up from a clean start, it names a fact that some earlier
passing COMMIT committed. But a later commitment replaces it, so at the end
it names the last earned commitment, which needn't be the one the flag rose on. Two runs from
$\mathrm{start}(0, 0)$ show it, worked from the instruction table. First run: CHECK \emph{is zero}
$A$; CHECK \emph{is zero} $B$; COMMIT \emph{is zero} $A$; CERTIFY; COMMIT \emph{is zero} $B$.
Second run: the same five instructions with the last two swapped, so COMMIT \emph{is zero} $B$
comes before CERTIFY. Both end in the same state: both counters 0 at version 0, both facts in the
table, the channel holding (\emph{zero}, $B$, 0), the trap down, a ledger of five and the flag
up. In the first run the flag rose at step four, on the commitment about $A$. In the second it
rose at step five, on the commitment about $B$. So the flag is an ever-certified latch: it says
an earned chain happened at some step, and the end state doesn't say which claim or which step.
A receipt naming the claim the flag first rose on would need that claim kept where nothing can
overwrite it; a receipt with a time on it would also need the step number kept; and several
certifications that each last would need a record richer than one flag. The publication file constructed earlier in this chapter keeps the set of
claims certified at least once. That extension shows exactly what the native
state left out. It still does not record which certification came first or
at what step; those questions would need their own retained answers.
\item \textbf{Sixteen checks per run, ever.} Facts never leave and the table holds sixteen, so a run
gets at most sixteen passing CHECKs in its whole life, and the seventeenth traps. A
certificate needs only one, so this costs nothing for the theorems, but it's a hard wall
for any program that wants to keep checking forever, and it's this machine's small version of
``you can't keep all of it'' from Section~\ref{sec:established}. A CHECK that repeats a claim already in the
table, same property, same counter, same edition, passes again and takes another entry. So the
wall is sixteen passing CHECKs, and sixteen different claims is a different count. Removing
duplicates, evicting old facts or collecting dead ones would each be a different design, and each
would need its own invariants proved.
\item \textbf{The clean start matters.} Provenance, soundness, no forging and the price of three are
all stated from a clean start. Hand the machine a table pre-loaded with facts nobody
checked, and of course none of that follows. The theorems are about what the machine
produces. What you smuggle in at the door is outside them. The door is narrower than clean for some of them: a commitment's provenance and checker soundness use only the table, and hold from any start whose table holds nothing but facts about older editions; the empty channel and the lowered flag come in for the flag's provenance and the price of three. Spell the door out once: the theorems
cover any list of instructions anybody can issue, from a clean start, on a step that runs as
written, with nothing else writing to the machine. They say nothing about whoever sets or
restores the start, about a reset, about a device that misbehaves, or about whoever carries the
flag's value to you.
\item \textbf{The universal programs have stated guests.} $U$ (Appendix~\ref{sec:universal-host}) runs every one of this machine's
programs from every clean start $\mathrm{start}(x, y)$, with halting, answers, flag and ledger matched (the flag over the whole run, the ledger at matching points). $U_P$ (Appendix~\ref{sec:presented-universal}) runs every
computably presented Thiele machine, with its reading decided by a program, its record carried as
a latch, and its ledger matched to within a surcharge of at most two when its reading starts at
no, and of three when it starts at yes. When the guest certifies with a ledger below three, no program on
that host, run from a clean start, matches its ledger exactly. A
machine outside that class, with a reading no program decides, isn't covered by either. The
simpler host with the built-in guest step (Appendix~\ref{sec:simpler-host}) pays one per guest step whatever the step was, so its
ledger is at least the guest's and usually more; $U$'s is exact, and $U_P$'s is exact up to the
surcharge.
\item \textbf{Encoding and running time.} The universality result takes its inputs already
packed into the counters. Getting there from an arbitrary Turing machine goes through
Minsky's exponential packing (Appendix~\ref{sec:counter-constructions}), and that packing is
what it takes to reach the two-counter base this way. Carrying the record to another base has
its own conditions, in Appendix~\ref{sec:every-model}, and anyone who wanted this to run fast would have to
start the analysis over.
\item \textbf{The step charges for the act, not for what the claim is worth.} CHECK \emph{at
least 0} passes on every value, so CHECK, COMMIT, CERTIFY on it raises the flag for three from
any fresh loaded start (the note after Theorem~\ref{thm:u34-tight}), and that's by design. The
step checks order, aboutness and an honest checker. What a claim is worth is the entitlement
count of Section~\ref{sec:entitlement}, a number computed over the rivals a run narrows. A
step sees one state, so it cannot compute that count. A check that cannot fail narrows nothing
and costs three, and the book keeps both numbers, apart.
\item \textbf{The toll on this machine comes from the record layer, with no finiteness premise on the whole machine.}
Section~\ref{sec:a2-not-arbitrary}'s counting of states forces a toll on a finite machine that prices
every merge. This machine is neither. Its counters have no ceiling, as a universal base has to have, and its editions can't have one either (the version rule says why), so it's infinite. And, by design, its unpriced instructions merge states: take a core at line $\ell$ with $c$ at 1 and version
$w$, and a core at line $j - 1$ with $c$ at 0 and version $w + 1$, everything else alike, and
DEC $c$ $j$ sends both to line $j$ with $c$ at 0 and version $w + 1$, at a price of zero. A price of zero
doesn't mean the hardware pays nothing. Any physical heat belongs to the implementation's account. A positive floor
requires the thermal premises and an input distribution with entropy loss;
merely taking a step establishes neither. This ledger does not measure that heat. The
ledger leaves those merges out on purpose. It counts the keeping in the record, not what computing
spends. Computing doesn't have to merge; a reversible base computes by paying in fresh states. On this machine
full merge pricing would charge DEC, an ordinary step, and clause (c) asks for computing to cost
nothing. Charging every move would instead give the clock-style schedule of
Observation~\ref{thm:u34-clock}. Its total counts moves, rather than isolating
record acts. The machine can still retain a record; this total alone does not
say which acts established it. So
merge pricing on this machine applies to the record layer only: a move that sends two different
states with the same line and counters to one state costs at least one. This machine's own prices
meet it. The endless counters would allow a raise that merges nothing, and the flag
doesn't take it: CERTIFY lands a state with the flag down and the same state with the flag up on
one state, a merge of exactly that kind, in plain view. So every price list that meets the premise charges the flag-raising step at least one, with no finiteness premise, and through the trap it charges every CHECK and COMMIT at least one too, so a certificate from a clean start costs at least three on every such list (Theorem~\ref{thm:record-merge-toll}). That's a price I assume for losing distinctions in the
record, and assuming it is all it is. Turning it into heat takes the physical conditions of
Section~\ref{sec:a2-not-arbitrary}; a unit on a ledger isn't heat on its own. The exact price of one, the least the premise allows on each of the three acts, is still a design choice.
\end{itemize}


\section{Running it}

Everything so far is definitions and proofs about them. Does any of it run? This appendix is
me finding out. I built the small machine twice, once straight out of the definitions and once
by hand from the same description, ran both, and compared the whole state after every step.
Then I ran the universal programs on guests small enough to finish, and real programs through
the part of the compiler that executes.

One rule holds for the whole appendix: the runs test the builds, and the proofs carry the theorems. A
run shows what a build did on the inputs it was given, and no theorem in the book leans on one: every theorem in it stands on its argument, on the page.

There are two counts to keep apart as you read. The ledger counts what the price list charges. The
number of interpreter steps counts the work this particular build does to carry the run out.
One guest step can take millions of host steps and add one to the ledger.

\begin{center}
\small
\begin{tabular}{@{}p{.27\linewidth}p{.63\linewidth}@{}}
Stage & What the next stage must preserve \\ \hline
Source orders & Variable values and completed procedure calls \\
Assignments and jumps & The same branches, loops and final answer \\
Counter blocks & Values at block boundaries; spare counters restored \\
Host instructions & The same result, with the opposite DEC convention translated \\
Packed two-counter guest & The result as an exponent in the packed counter \\
Fixed interpreter $U_P$ & Guest counters at matching points; decode the exponent to read the answer \\ \hline
\end{tabular}
\end{center}
This is a path through the compiler, not a claim that one step at each stage
takes the same time. The source compiler inserts no certification chain. Running the executable uses physical hardware and energy. These runs did not measure
heat per host step, and the ledger is not a heat measurement. The times only
tell you which runs are practical on this build. The table shows what ran and what didn't.

\begin{center}
\begin{tabular}{@{}p{0.42\linewidth}p{0.5\linewidth}@{}}
\textbf{What} & \textbf{Status} \\ \hline
The small machine & run: the executable versions agree on every tested step \\
The host machines, with and without PAY & run, compared the same way \\
$U$ and $U_P$ on tiny guests & run, against a separate guest implementation \\
$U$ on a guest that certifies & proved and counted, not run \\ $U_P$ on a guest that certifies & proved, not run \\
Source programs through the compiler & run as far as the host program \\
The two-counter guest and $U_P$ running it & proved, not run \\
The presented machine carried by $U_P$ & proved, not run \\
\end{tabular}
\end{center}

\subsection{Taking the coat off}

The first executable is generated straight from the formal definitions, the way you'd
strip the coat off a proof and keep the program underneath. The second was written
separately, by hand, from the same description. Both implement the small machine, the host with a register for every whole number, the host with PAY and the loaders, and both run the universal programs; their instruction lists the second build reads as data printed from the definitions, so it checks them by running them, not by writing them again. The compiler below has an executable runner too.

Whole numbers matter here. A machine word in one build is 63 bits wide, and the guest
codes $U$ reads go up past $2^{100}$, so ordinary fixed-width arithmetic would wrap around
and lie. The mathematical machine never overflows, so its executable uses integers of any
size, with the arithmetic chosen to match the definitions: subtraction stops at zero and
division by zero gives zero. What sits between the definitions and the running code, the
mappings, the compilers and the routines that read input and write output, is code I had to
trust.

\subsection{A second machine, written by hand, to check the first}

Comparing two builds step by step catches one of them disagreeing with the description, on the
runs I tried. They're compared on the \emph{whole} state after \emph{every} step: the program line, every counter
and register and its version, the fact table in order, the channel, the trap latch, the ledger and
the flag. The comparison covers exhaustive sets of short programs, random programs and both hosts, and
they agree on every tested step. $U$ and $U_P$ on tiny guests are compared at fifteen checkpoints two hundred steps apart, over the first three thousand host steps of each run, instead of after every step, and the longer runs below against a separate guest
implementation, once per pass of the loop and at the stop. I also broke the machine on purpose, giving it a fact table one too small, a
CHECK, COMMIT and CERTIFY that cost nothing, a commitment that ignores versions, and a write that doesn't bump the version, and each break
gets caught. A second set pins the boundaries this book draws in words, and those runs come out as the words say: two runs that end in the same state with different first certifications, a check that
can't fail certifying from every start, a CHECK and six commitments costing seven without a flag, a write
after the commitment and a flag over a claim that's now false, sixteen copies of one CHECK filling
the table, a CERTIFY costing one at a ledger already at two, and a merge between two yes-states
that leaves the entering state alone.

A mistake both builds share gets past the comparison, and so does a mistake in the definitions
themselves, since both builds follow them.

\subsection{\texorpdfstring{$U$}{U} on a real guest}

Now the universal program, which is where I expected trouble. The guest I wanted to run first was
the obvious one, a guest that goes the whole way: CHECK, COMMIT, CERTIFY, flag up, run inside $U$ on
the built machine. That isn't what ran, and the reason is the way the guest is stored. $U$ keeps the whole guest program as one number,
a number that grows like two raised to the sum of the codes for the guest's instructions (for a
guest of one instruction it's exactly a power of two), and it reads that number
the only way a counter can, one unit at a time. So its speed depends on the code, and the code grows
fast: a guest with one CHECK in it has code $2^{24}$, and a guest with the whole chain has a code
near $2^{106}$.

So I came down to what could finish. The table below is what ran, on the executable host, compared against a separate
guest implementation at every pass of the loop and at the stop: the host stopped exactly when the guest
did, and at the stop showed the guest's counters, trap latch, ledger and flag.

\begin{center}
\begin{tabular}{@{}llrr@{}}
guest & host & host steps & time \\ \hline
count A down from 2 & $U$ & 25,953,317 & \\
one passing CHECK & $U$ & 402,654,059 & 618 s \\
one failing CHECK & $U$ & 201,327,155 & 305 s \\
one passing CHECK & $U_P$ & 402,654,059 & 617 s \\
one failing CHECK & $U_P$ & 201,327,155 & 313 s \\
\end{tabular}
\end{center}

The first row's count is pinned exactly: if the build changed so that it moved by one step, the
comparison would fail. The times are from one ordinary computer and are only there to give you the
scale. The counts work out to about twelve host steps for every unit of the guest's code each time $U$ goes round
its loop, which it does once for every guest step and once more to see that the guest has
stopped: the failing CHECK traps on its first pass, $12 \times 2^{24}$ and a little, and the
passing one takes two passes, $24 \times 2^{24}$ and a little. The growth is exponential in the size of
the guest. A guest with one CERTIFY in it would have code $2^{32}$, about fifty billion host
steps, and at the measured rate of about 650,000 steps a second it would take something like
twenty-two hours, nearly a day. The guest I wanted, with the whole chain, has a code near
$2^{106}$, which at the same twelve steps a unit is about $10^{33}$ host steps for each pass; at the
measured rate, with this encoding on that computer, a run of its three steps comes to around
$10^{20}$ years. So $U$ on a guest that certifies is proved, in Appendix~\ref{sec:universal-host},
and has never been run. Runs this long only fit in memory at all because the build can repack the
host's register store in the middle of a run; repacking changes how the registers are stored, not
what they hold.

What a run that long can be is counted. Four loop shapes do most of $U$'s work.
Count from the DEC instruction that tests the loop counter; any instructions before
that test are still counted separately. You can derive each shortcut from the arrows:
\begin{itemize}
\item In the emptying loop, a positive counter $v$ takes one DEC back to its own test
and becomes $v-1$. At zero, one DEC exits. Thus $E(0)=1$ and
$E(v+1)=1+E(v)$, giving $E(v)=v+1$.
\item In the copying loop, a positive DEC jumps back one line to INC on a distinct
destination, which returns to the test. Each unit moved therefore takes two steps,
and the zero test takes one: $C(0)=1$, $C(v+1)=2+C(v)$, so $C(v)=2v+1$.
\item With two destination INCs before the test, each positive pass takes three steps.
Thus $D(0)=1$, $D(v+1)=3+D(v)$ and $D(v)=3v+1$. If the two destinations are the same
register, it receives two increments per pass; the step count stays three.
\item In the halving loop, two successful DECs on the source and one INC on the
quotient return to the same test. After $k=\lfloor v/2\rfloor$ such passes,
$3k$ steps have removed $2k$ from the source and added $k$ to the quotient.
The remaining source is zero or one. Follow that final branch instruction by instruction.
\end{itemize}

Each successful decrement or increment also raises its register's version once.
The zero test does not write. So the same counts determine the version changes as well
as the counter values and next line. These loops neither charge nor change the table,
channel, trap or flag. Replacing one completed loop stretch by its derived effect and
count therefore preserves the whole endpoint and the number of steps. Repeating that
replacement throughout a run preserves both by addition.

The numerical tally below applies those four rules and steps the other instructions
individually. It agrees with the three runs above that were also stepped all the way.
For the guest CHECK, COMMIT, CERTIFY, using the zero predicate on $A$, the passing start
$(0,0)$ ends with the flag up and ledger three. From $(5,0)$ the first check fails and
the guest stops with ledger one. Appendix~\ref{sec:universal-blocks} supplies the program
blocks and their guest correspondence; the displayed totals count that particular encoding.

\begin{center}
\small
\begin{tabular}{@{}lr@{}}
start & host steps \\ \hline
$(0,0)$ & 3,894,222,629,847,358,066,688,933,266,198,944 \\
$(5,0)$ & 973,555,646,581,507,138,873,939,791,446,643 \\
\end{tabular}
\end{center}

What that table shows is a count, not a run. Nothing stepped those $10^{33}$ instructions, and the row in the table
at the top of this section says so.

Those big counts come from this encoding and from reading it one pebble at a time. They
measure the work this interpreter does, and that's all they measure. Whether some other
interpreter, or some other base, could do the same job fast is a separate construction with its
own running-time analysis. Nothing measured here puts a floor under every way of doing it.

\subsection{From source programs to the host}

Running guests one CHECK at a time doesn't tell you whether the built machine can do anything you'd
recognise as a program. You'd want to hand it gcd, or a sort, written the way people write them, and
get the answer out, and writing gcd by hand in INC and DEC is no way to find out. So the last piece
of the build is a compiler. The source is a small language you could write by hand: whole-number
variables, addition, subtraction that stops at zero, comparisons, \emph{if}, \emph{while}, and
procedures that return several results, with the rule that a procedure can only call ones defined
before it, so there's no recursion. Each translation stage has a proof that it preserves the program's result: the procedures are
inlined, the result is flattened to assignments and jumps, that is compiled to Minsky's counter
machine through the counter-machine compiler, that to a program for the host machine, that to a
two-counter guest, and the guest to a run of $U_P$.

\begin{theorem}[The compiler is correct]
For every source program with no recursion, and every input: the program computes $y$ exactly when
the runner (the bounded evaluator of the compiled host program) halts with $y$ for some amount of fuel; exactly when the compiled host program
halts with $y$; exactly when the two-counter guest halts with $y$; and exactly when $U_P$ running that
guest halts with $y$, with its first counter at 0, its ledger at 0 and its flag down. So a terminating compiled program delivers its answer and certifies nothing. To certify a
claim about that answer, a target program must add a CHECK, a COMMIT and a CERTIFY;
this compiler does not insert that chain. The runner also
takes exactly as many steps as the host program does, with the rest of its fuel left over, and a run
that doesn't halt has used all its fuel.
\end{theorem}

\begin{proof}
Six arrows and a runner, each argued below, and the statement is their composition. Four of the
arrows are between machines that take steps, and for those one argument does both directions, so
I give it first.

\emph{The simulation argument.} Take two machines $S$ and $T$, each deterministic, where $S$ has a
next state from every state inside its program, and a relation $\sim$ between their states such
that: a state of $S$ inside its program is related only to states of $T$ inside its program, and
a state of $S$ outside it only to states of $T$ outside it; and from any related pair with $S$
inside its program, one step of $S$ is matched by one or more steps of $T$ that land related to
$S$'s next state. Then from related starts, $S$ halts exactly when $T$ does, and they halt
related. Forward: start with the related pair. Follow the stretch of $T$ that
matches $S$'s first step; its endpoint is the related start needed for the
second stretch. Keep joining stretches until $S$ stops. There are only
finitely many steps to match, so this gives a finite run of $T$ whose final
state is related to $S$'s final state. The latter is outside its program,
so the former is outside its own. Backward: suppose $T$ halts after $k$ steps and $S$ doesn't halt. Then $S$ takes at least
$k + 1$ steps without leaving its program, the matched stretches give $T$ a run of at least
$k + 1$ steps all of whose stretch ends are inside its program, and $T$, being deterministic, has
only one run from its start, which halted at step $k$. That's a contradiction, so $S$ halts, and
by the forward direction it halts related to where $T$ did.

\emph{Unfolding the calls.} Imagine a program written on cards. A call says
``carry out the orders on that other card, then come back.'' Here a card may
call only a lower-numbered card. Card zero therefore calls nothing. Replace
every call on card one by the orders from card zero; then do card two, using
the already expanded lower cards, and continue. There are finitely many cards,
and every replacement has already been built. The process ends with calls
replaced by actual orders.

The variables need their own places too. Give each active procedure a block
of numbered boxes, its \emph{frame}. If the caller's block starts at box $b$
and uses $n$ boxes, put the called procedure's block immediately after it,
starting at $b+n$. A call becomes four pieces: copy the arguments into that
new block; set its other boxes to zero; carry out the expanded orders using
those box numbers; copy the answers back into the caller's named boxes.
A local variable numbered $x$ in the called procedure is thus the global
box numbered $b+n+x$. No assignment reaches below the procedure's own block.

Why does this preserve a terminating run? Begin with the lowest-numbered
procedure. It has no calls, so its expanded orders are its original orders
with the boxes renamed. Assignments change the same values, tests choose
the same branches, and loops repeat the same bodies. Now take a procedure
whose lower-numbered calls have already been settled. Each replacement call
builds exactly the arguments and zeros that the original call would receive.
The expanded lower procedure returns the same answers, and the copy-back
puts them in the same caller boxes. The caller's surrounding orders then
continue from the same values. Working upward through the numbered procedures
settles the main program.

The comparison works backward as well. Whenever an expanded call finishes,
its four pieces have run in their written order. The middle piece is the
expanded lower procedure, whose terminating runs I have already matched to
the original lower procedure. Its preparation and copy-back are just the
original call's input and output operations. Replace that finished four-piece
execution by the call it implements. The surrounding assignments, branches
and completed loop repetitions then give the same original execution.
A finite execution in either description therefore gives a finite execution
in the other, with the same answer. Neither can terminate when the other does
not.

In symbols, matching a source arrangement $e$ to the global boxes $G$ in a
frame starting at $b$ means $G(b+x)=e(x)$ for each local box $x$. The comparison
above preserves that equality at the beginning and end of every translated
statement.

\emph{Flattening.} The flat code of a statement placed at address $a$: an assignment is one
assignment; a sequence is the first statement's code followed by the second's; \emph{if} $b$ is a
jump on false to the else part, the then part, a jump past the else part, and the else part; and
\emph{while} $b$ with a body of $l$ instructions is a jump on false to $a + l + 2$, the body at
$a + 1$, and a jump back to $a$. Every jump in the code of a statement lands inside that code or
at its end, the address just past it. Now follow a completed structured execution.
An assignment performs the same change; a sequence joins its two runs; a branch
follows the same test answer. For each loop iteration, the test falls through,
the body's code takes the run to $a+l+1$, and the jump returns it to $a$.
Repeat for the remaining iterations; the final false test jumps to the end.
Thus an execution from variables $e$ to $e_1$ gives a flat run from address $a$
to the end of the code with those same final variables.

Read a completed flat run in the other direction. Because no jump escapes past
the code's end, a run that leaves it must visit that end first. Assignments and
branches recover their structured operations immediately; a sequence splits at
the second statement's first address. For a
loop, the run's first step is the test. If it is false, the run jumps to the end, and the loop's
derivation is the one with a false test. If it is true, the run enters the body's code, which it
can leave only through its end at $a + l + 1$ (the body's own backward statement), where the jump
returns it to $a$ with a strictly shorter run left. Read the next iteration the same way.
This cannot continue beyond the finite run I started with: its final test must
be false, giving the last piece of the structured loop execution. So the
flat program at address 1 halts exactly when the structured statement has a derivation, with the
same final variables.

\emph{Counter machine.} Here a DEC jumps when its register is zero and goes on otherwise. Variable
$x$ lives in register $x + 1$, register 0 is kept at 0, so DEC 0 $j$ is a plain jump to $j$, and
registers above the variables are working space. The relation: register 0 is 0, registers 1 to
$n$ hold the flat variables, and every register above is 0. Each flat instruction becomes a block built from small loops. Follow one
pass and then count how many pebbles remain to be moved:
Emptying a register: each pass takes one off, and the loop exits at zero. Moving $s$ into $d$:
each pass takes one off $s$ and adds one to $d$, so $s + d$ never changes, and at exit $s$ is 0
and $d$ holds the old sum. Copying: move $s$ into two registers at once, then move one of them
back. Cut-off subtraction $t - u$: take one off each while both are positive; at exit one of them
is 0, and the other holds the difference, or 0. Comparison: take one off each in step; whichever
reaches zero first, or both together, decides $<$ and $=$, and the block empties both. Build a larger expression from these smaller jobs. A constant is a
fixed number of increments into an empty working register. A variable is
copied there with its original restored. For a sum or cut-off subtraction,
compute the two smaller expressions in separate working registers, combine
them with the loops just described, and empty the temporary registers.
Starting with the innermost pieces and working outward, the final value
lands in register $t$, the registers above $t$ return to zero, and those
below $t$ are left alone. A test computes its comparison in the same spare
registers, chooses the true or false exit, and clears the spare registers
on either exit. Combining tests follows the selected exits of their smaller
tests, each of which has already restored the registers. Thus a test leaves
the registers as it found them and changes only which address comes next. So an
assignment block computes the value into the first working register, empties register $x + 1$,
moves the value in, and lands related; a jump on false runs the test's code with the true exit at
the next block and the false exit at the block of $j$; a jump is DEC 0 $j$. The linker places the
blocks one after another and sends every jump to the start of its target's block. That's the
one-step matching. The flat machine has a next state from every address in its program, and an
instruction that names a variable outside the frame just advances, so the step is total. The
counter machine is deterministic. The simulation argument gives both directions: the counter
program from the loaded inputs leaves its code with register $\mathit{out} + 1$ equal to $y$
exactly when the flat program halts with $y$ in $\mathit{out}$.

\emph{Host.} The host's DEC is the other way round: it jumps when its register is positive, taking
one off, and goes on at zero. An INC $x$ stays INC $x$. A DEC $x$ $j$ at host address $a$ becomes
three instructions: DEC $x$ $(a + 3)$, then INC 0, then DEC 0 $j'$, where $j'$ is the host address
of $j$'s block. If register $x$ is positive, the first takes one off and jumps to $a + 3$, the
start of the next block, which is what the counter machine does. If it is 0, the run goes on to
INC 0, which raises register 0 to one more than it held, and DEC 0 $j'$ finds it positive, takes
the one back off and jumps to $j'$, which is the counter machine's jump at zero, with register 0
as it was. The relation: the same registers, at corresponding addresses. Each counter-machine step
is matched by one or three host steps. The host program has no CHECK, COMMIT, CERTIFY or PAY, and
INC and DEC change only a register and the line, at no cost, so every host state in the run has
the start's facts, channel, trap latch, ledger and flag. The counter machine's step is total, the
host is deterministic, and the simulation argument gives: the host program halts with register
$\mathit{out} + 1$ equal to $y$ exactly when the counter program leaves its code with that
register equal to $y$.

\emph{Guest.} The host program's registers $r_0, \ldots, r_{k-1}$ are packed into the guest's
counter $B$ as $\prod_i p_i^{r_i}$, one prime $p_i$ per register, with counter $A$ the spare. That
is the packing of Appendix~\ref{sub:prime-packing}, with the two counters' jobs swapped. The
relation: the guest is at the start of the block for the matching instruction, $B$ holds the code
of the registers, and $A$ is 0. An INC on register $i$ is the stretch that multiplies $B$ by
$p_i$; a DEC on register $i$ is the stretch that divides $B$ by $p_i$ with remainder and takes the
positive branch when the remainder is 0, the zero branch otherwise, restoring $B$. The appendix
argues each stretch: it ends, because each of its loops empties the counter it counts down; it
lands with $A$ at 0; and, by unique factorisation, it changes the exponent of $p_i$ alone, by the
right amount. The host's step is total on these programs and the guest is deterministic, so the
simulation argument gives: the guest halts with $B$ coding registers whose register
$\mathit{out} + 1$ is $y$, that is, with $y$ the exponent of that register's prime in $B$, exactly
when the host program halts with that register at $y$. The guest has no record instruction, so its
ledger stays 0 and its flag down.

\emph{$U_P$.} $U_P$ keeps pace with its guest the way $U$ does, with the PAY case added (the proof of Theorem~\ref{thm:u34-presented}). This guest has no PAY and no record instruction, so on it $U_P$ runs only branches $U$ has, and the argument of Theorems~\ref{thm:u34-U-sim} and~\ref{thm:u34-U-halting} goes through: from a freshly loaded host $U_P$ halts exactly when the guest does, with its registers 0 and 1 holding the guest's $A$ and $B$, and its
ledger and flag equal to the guest's. So it halts with register 0 at 0, register 1 the code with
$y$ in it, ledger 0 and flag down.

\emph{The runner.} The runner keeps the host program and the registers in binary tries, a tree
indexed by the bits of a number. Read the address one bit at a time: a zero chooses one branch and a one
the other. Looking up an address follows that route to its stored value,
or returns the default zero if no value was stored there. Storing at an
address follows the same route and changes its final value, creating missing
branches as needed. Every other address either branches off before that end
or names a different node, and keeps its value. Thus the tree performs the
same lookup and single-address update as the register map it represents. So each round of the runner's loop does
exactly what one host step does to the same registers, and the runner stops when the line leaves
the program or the fuel is spent. If it halts after $k$ rounds, the host run of $k$ steps is
halted at the end of the program with the runner's registers, and the fuel left is the fuel given
less $k$. If the source program has a derivation, the arrows above give a halting host run, of
some length $k$, and the host, being deterministic, has only that run; any fuel of at least $k$
makes the runner take those $k$ steps and stop. If the runner hasn't halted, it used every unit of
fuel on host steps that stayed inside the program, so the host is still running after that many
steps, and a derivation, which would give a halting host run of some length $k$, would need $k$
larger than the fuel.

Composing: the program computes $y$ exactly when the unfolded program does, exactly when the flat
program halts with $y$, exactly when the counter program does, exactly when the host program does,
exactly when the guest halts with $y$ coded in $B$, exactly when $U_P$ halts as stated, and exactly
when the runner halts with $y$ on some fuel.
\end{proof}

The theorem says nothing about how many steps the compiled program takes next to the source; it's
about answers.

The runner and the compiler are implemented, and the build has been run on a set of programs I picked
to be annoying: gcd, factorial, a primality program written with procedures, an Ackermann tower, a sorting
network on a list packed into one number, and a tiny interpreter for another machine written in the
source language. Then come 220 random programs that halt, 160 random programs with loops that might not,
programs that never halt run for a fixed budget, and programs mutated on purpose to make sure the
comparison notices. Answers are compared with direct arithmetic and separate interpreters of the source and
compiled programs. Ackermann's function at $(3, 5)$ is 253; the
source program takes 64,768 operations to get there, and the compiled host program takes about 119
million steps. The executable runner does a few million host steps a second, so that's under a minute.

That works out to about 1,800 host steps per source operation for Ackermann, and up to about a hundred thousand
for multiplication by repeated addition, because every addition copies a running total one unit at a
time. On this build the work grows with the size of the numbers being handled, because every
operation is one pebble at a time. How it grows with the size of the program I haven't
measured. The last two stages, the two-counter guest and $U_P$ running it, are proved and have
never been run, and $U_P$ on a compiled guest would be slow for the same reason $U$ on a
certifying guest is. Still, I like that 253. It came out of nothing but INC and DEC, and a proof
that they mean what the source meant.

\subsection{What wasn't built}

There's no chip. The model is not a circuit, and the lack of one says nothing about it. The example was never cut into metal: there's no circuit for the small machine, no
simulation of one at the level of gates, and no board ever ran any of this. The small machine exists
as mathematical definitions and executable implementations that
agree on every tested step. That's the whole of what ``it runs'' means in this book.

I'd love to see it in metal some day. But builds that can be held against each other on every step are worth a lot to
me.
A board with nothing to hold it up against isn't, and that's the trade I made.

%===========================================================================


\section{Constructing the quoting operation}

The proof in Section~\ref{sec:diagonal} uses a program that takes the code of a term
and returns the code of that code. Here's how to build it.

Draw the instruction ``add one'' on a card. Now make a second card describing the
first card's marks, rather than carrying out the addition. The count has not risen:
you have described an instruction. Describe the second card in turn and you have
a description of a description. Quoting does that second job mechanically. In the
language below, take the variable of index 0. Its code is a tree with three lambda
nodes, then an application of variable 2 to the numeral for 0. Quoting that code
builds a description of those three lambdas and that application; it does not
look up variable 0 or execute the term it describes.

Nothing in the operation lets a program inspect itself by magic: its input is an
ordinary tree encoded as a closed lambda term.

The construction follows Forster and Smolka (2017), whose construction for L has the same
three syntax builders on Scott-encoded children, an auxiliary procedure for numerals, and a quote
procedure $Q$ with the same specification, defined by the same recursion on codes; the
self-application that uses $Q$ in Section~\ref{sec:diagonal} is their fixed-point lemma, after
Barendregt's (1984) proof for the full lambda calculus. What follows is my rebuilding of it, with
the reductions written out.

\subsection{Numerals and syntax as data}

For readability, the formulas below give names to bound variables. Renaming a bound
variable changes nothing; the underlying syntax uses the de Bruijn indices of
Definition~\ref{def:u12-L}. Applications associate to the left. Write $X\Downarrow Y$
here for a finite sequence of the weak call-by-value steps ending in the value $Y$.

Scott numerals are the following closed values:
\[
\overline{0}=\lambda z.\lambda s.z,
\qquad
\overline{n+1}=\lambda z.\lambda s.s\,\overline n.
\]
So a numeral selects either the supplied zero case or the supplied successor case,
passing its predecessor to the latter. The cases below are values, so choosing one
does not prematurely run the other.

The three syntax constructors, operating on already encoded children, are
\[
\begin{aligned}
\mathrm{mkvar}&=\lambda n.\lambda v.\lambda a.\lambda l.(v\,n),\\
\mathrm{mkapp}&=\lambda p.\lambda q.\lambda v.\lambda a.\lambda l.a\,p\,q,\\
\mathrm{mklam}&=\lambda p.\lambda v.\lambda a.\lambda l.l\,p.
\end{aligned}
\]
Each builder is a term of the language $L$ of Definition~\ref{def:u12-L}.
For a variable of index $n$, an application and an abstraction, respectively, define
their codes by the following resulting values:
\[
\begin{aligned}
\ulcorner\#n\urcorner&=\lambda v.\lambda a.\lambda l.(v\,\overline n),\\
\ulcorner p\,q\urcorner&=\lambda v.\lambda a.\lambda l.a\,\ulcorner p\urcorner\,\ulcorner q\urcorner,\\
\ulcorner\lambda p\urcorner&=\lambda v.\lambda a.\lambda l.l\,\ulcorner p\urcorner.
\end{aligned}
\]
In particular, $\mathrm{mkvar}\,\overline n$, $\mathrm{mkapp}\,\ulcorner p\urcorner\,\ulcorner q\urcorner$ and
$\mathrm{mklam}\,\ulcorner p\urcorner$ reduce to these three codes. A code is closed even when
the term it describes has a free variable: the index of that variable is stored
as a numeral. It is data inside the code, and the code itself has no free variable.

\subsection{Recursion without running under a lambda}

The quoting procedure will have to call itself on the smaller pieces of a
description. There is a catch: this language works out an argument before
handing it to a function. Hand over the recursive calculation itself and it
starts calculating at once, before the receiving function has decided whether
it needs it. That can send you straight back into the same calculation forever.

Instead, hand over a recipe waiting for an ingredient. In this language a
lambda is such a recipe, and it counts as a finished value even though its
body has not run. The receiving function can keep the recipe and call it when
it reaches a smaller piece. In the expression below, the little recipe
$\lambda a.x\,x\,f\,a$ is the part doing that job. The letter $a$ names the
argument it is waiting for; $f$ names the work to do once the recipe and the
current input have been handed over.
\[
W=\lambda x.\lambda f.\lambda v.\ f\,(\lambda a.x\,x\,f\,a)\,v,
\qquad
\operatorname{rec}(F)=\lambda v.W\,W\,F\,v.
\]
When $F$ and $u$ are closed values, beta reduction gives
\[
\operatorname{rec}(F)\,u
\longrightarrow W\,W\,F\,u
\longrightarrow^{*} F\,(\lambda a.W\,W\,F\,a)\,u
=F\,(\operatorname{rec}(F))\,u.
\]
The recursive argument is a lambda, so it's already a value. It starts work only when
$F$ applies it to an argument. That is the whole reason this construction survives
the call-by-value rule.

\subsection{Quoting numerals}

Let $E_j=\ulcorner\#j\urcorner$ for $j=0,1,2$, and let
$Z=\ulcorner\overline0\urcorner$, obtained directly from the finite syntax of
$\overline0=\lambda\lambda.1$. Write $\mathrm{mklam}^k(X)$ for $k$ nested applications of the
syntax builder $\mathrm{mklam}$ to $X$; it is an abbreviation for constructing $k$ lambda nodes, and the
language gains no new operation from it.

Define
\[
F_N=\lambda r.\lambda n.\ n\,Z\,
       (\lambda k.\mathrm{mklam}^2(\mathrm{mkapp}\,E_0\,(r\,k))),
\qquad Q_N=\operatorname{rec}(F_N).
\]
Follow the two branches. Zero chooses the already built description $Z$.
A number one larger than $n$ chooses the other branch and hands it the
predecessor $n$. That branch first asks the same procedure to describe the
predecessor. Then it puts the description inside an application of variable
zero and wraps two lambdas around it. Those are precisely the marks in the
numeral for $n+1$.

Zero therefore works. If the procedure describes $n$ correctly, the wrapping
just described makes it describe $n+1$ correctly. Starting at zero and taking
one more step proves the claim for every natural number. Every call on a
positive number asks about a smaller one, so it reaches zero and finishes.
In the notation for termination and returned value, the result is
$Q_N\,\overline n\Downarrow\ulcorner\overline n\urcorner$.

Do the smallest case by hand. Call the successor case $G$. Then $Q_N\,\overline1$ steps to
$F_N\,Q_N\,\overline1$ and on to $\overline1\,Z\,G$; the numeral hands its predecessor
$\overline0$ to $G$, the recursive call $Q_N\,\overline0$ comes back as $Z$, and the builders
return $\mathrm{mklam}(\mathrm{mklam}(\mathrm{mkapp}\,E_0\,Z))$. That is the code of
$\lambda\lambda.(\#0\,\overline0)$, which is $\overline1$.

\subsection{Quoting terms}

For a proposed recursive call $r$, define three handlers:
\[
\begin{aligned}
H_V&=\lambda n.\mathrm{mklam}^3(\mathrm{mkapp}\,E_2\,(Q_N\,n)),\\
H_A(r)&=\lambda p.\lambda q.\mathrm{mklam}^3(\mathrm{mkapp}\,(\mathrm{mkapp}\,E_1\,(r\,p))\,(r\,q)),\\
H_L(r)&=\lambda p.\mathrm{mklam}^3(\mathrm{mkapp}\,E_0\,(r\,p)).
\end{aligned}
\]
Now set
\[
F_Q=\lambda r.\lambda t.t\,H_V\,(H_A(r))\,(H_L(r)),
\qquad Q=\operatorname{rec}(F_Q).
\]
The handlers are lambdas and so are values. The Scott code of the input chooses
one of them and hands it the encoded children. Start at the leaves and build
outward. At each node the handler builds the syntax of that node's encoding:
$Q\,\ulcorner t\urcorner\Downarrow\ulcorner\ulcorner t\urcorner\urcorner$. There are
three kinds of node to consider:

\begin{itemize}
\item For $t=\#n$, the first handler constructs three lambda nodes over an application
of variable $2$ to the quoted numeral. That is the syntax of
$\lambda\lambda\lambda.(\#2\,\overline n)$, the code of $\#n$.
\item For $t=p\,q$, quote the two children first. The second handler
constructs three lambdas over variable $1$ applied to those two codes. That is
the syntax of the application code.
\item For $t=\lambda p$, the last handler constructs three lambdas over variable $0$
applied to the quoted child code. That is the syntax of the abstraction code.
\end{itemize}

Each recursive call works on a proper child of the finite input tree, or on a
smaller numeral. No chain of calls can keep descending through a finite tree
or decreasing a natural number forever. Each builder takes finitely many beta
steps, so the construction finishes and returns the code just described.

The $\mathrm{mkapp}$ used in Section~\ref{sec:diagonal} is this builder. With $Q$ and that
builder written out, every piece of the self-referential term there is on the page.

%===========================================================================

\section{Computing with counters}

Put a few beans in one cup and leave a second empty. Write a short list of orders:
add one, take one away if there is one, go to this line, stop. A finger on the page
keeps your place. Carry out the orders literally. The question in this appendix
is how to give the same job to another person with different equipment (a strip
of paper, stacks of cards, or a different set of cups) and get the same answer
whenever the first person stops.

Agree on how the input is laid out and how to read the output before either starts.
One person must not quietly solve the problem while preparing the other's cups.
The reductions in this appendix concern functions on natural numbers. To compute a
relation $R(\vec x,y)$ means to halt with output $y$ exactly when $R(\vec x,y)$ holds,
and to produce only numerical outputs on halting runs. The machines below are
deterministic, so any relation they compute is the graph of a partial function.
A translation must preserve both the output and whether the run terminates.
Its input and output encodings are part of the translation.

\subsection{Finite programs made from two instructions}

A counter holds a natural number. INC adds one and advances to the next line.
DEC subtracts one and jumps to its named line when the counter is positive; at zero,
it leaves the counter alone and advances. A finite program uses finitely many counters.
An unconditional jump to $j$ is INC $r$ followed by DEC $r$ $j$: it restores $r$'s
value and always jumps. A zero test is DEC $r$ to a positive branch that first restores
$r$; its fall-through is the zero branch. These operations preserve values, not versions.
They are used here for computation, where there are no live facts to preserve.

Try copying a cup with three beans. Put two empty cups beside it and keep a supply
of spare beans. Each time you remove a bean from the first cup, put one in each of
the other two. When the first is empty, move the third cup's beans back into it.
Now the first and second each have three, and the third is empty. Try zero beans,
then five. At every pass one more of the original beans has been counted into
both other cups; that is why it works for any starting count, not just the three
you tried. These are the moves behind the copy $C(x,y,t)$ in the next appendix.

Emptying a counter is a DEC loop back to its own line. Moving a pile decrements
its source and increments its destination until the source is empty. Copying
uses two such destinations and then restores the source from one of them.
Addition, multiplication by a fixed natural number, and bounded loops are finite
combinations of these moves. Division by a fixed positive number $b$ needs only
$b$ control states: decrement the input, cycle through remainders $0,\ldots,b-1$,
and increment a quotient counter whenever the cycle returns to zero. On exit the
quotient is in that counter and the remainder is in the control state.

\subsection{A tape in finitely many counters}

Cut a strip into numbered squares and put a finger on one square. Lift the squares
on its left into a stack, nearest on top, and do the same on its right. Moving your
finger one place right is now two acts: put the square you were reading on top of
the left stack, then take the top of the right stack as the new square. Write a
new mark before moving and the stack keeps the changed square. Move left and the
two stacks exchange jobs. You can walk the strip either way without losing its
order. An exhausted stack supplies a fresh blank square.

Now replace each stack by a count. Use digit 1 for a blank and digit 2 for a mark,
with base 3. Push a mark onto an empty stack: its code is 2. Push a blank on top:
its code is $3\cdot2+1=7$. Count seven beans in groups of three. The one left over
says ``blank on top''; the two groups say the remaining stack has code 2. Pop that
and the marked square comes back. Multiplication pushes and division pops.

In general, number a finite tape alphabet by $1,\ldots,q$, including its blank,
and put $b=q+1$.
Encode a finite stack whose top symbol is $d$ and whose remaining stack has code $n$
by $bn+d$; the empty stack is 0. A push is multiplication by $b$ followed by $d$
increments. A pop is division by $b$: the remainder is the top symbol and the quotient
is the remaining stack. The digits are nonzero, so no nonempty stack has code 0.
Popping an empty stack gives remainder 0, which is no digit; the control reads it as the blank and leaves the stack empty.

A tape configuration is a control state, the scanned symbol, and two stacks: the
symbols to the left of the head, nearest first, and the symbols to its right,
nearest first. The control state and scanned symbol have finitely many possibilities,
so they fit in the program's line number. To write $d$ and move right, push $d$ onto
the left stack and pop the right stack for the new scanned symbol. A left move does
the reverse. The arithmetic just described implements each transition in a finite
terminating stretch of counter instructions. Extra stored blanks do not change the
represented tape. A finite number of tapes uses two stacks per tape and finitely
many scratch counters. Halting states translate to halting lines. Induction over tape
steps proves the simulation; termination of every arithmetic stretch proves its converse.

In the other direction a tape can hold the finitely many counter values as unary
strings separated by delimiters. INC appends a mark in the selected string; DEC erases
one when present and otherwise detects the empty string. The current instruction is
finite control. Each scan and rewrite terminates, so halting and the output are preserved.
Multiple tapes may be represented on one tape by a list of tracks with marked head
positions: a complete scan reads the marked symbols, a second scan writes their changes
and shifts their markers. Extend the represented interval by a blank cell when needed.
A fixed-length binary code for the finite alphabet implements each scan with binary
symbols and finitely many intermediate control states. So this argument also applies
to the single-tape binary halting question.

\subsection{Packing finitely many counters into two}


Try holding two counts in one pile of eighteen beans. Count how many times you can
halve eighteen exactly: once. Count how many times you can divide it exactly by
three: twice. Reassemble the eighteen beans after each reading. Those are the two stored counts,
1 and 2. Double the pile to 36 and
only the first stored count rises. Divide 36 by three to get 12 and only the second
falls. Divide by three again: four beans, storing 2 and 0.

Try taking one more off that second count. Four beans do not divide evenly into
threes. The attempted division must restore the four and take the zero branch;
keeping only the quotient would corrupt the first count. This is why the procedure
below remembers the remainder. Two and three share no factor, so changing the
number of one kind of factor leaves the other alone. Assign another prime to
each extra count. An all-zero collection is represented by one bean: the empty
product is 1.

Let $p_0,\ldots,p_{k-1}$ be distinct primes. At the boundary between simulated
instructions, keep
\[
A=\prod_{i=0}^{k-1}p_i^{r_i},\qquad B=0,
\]
where $r_i$ are the simulated counter values. The simulated line is finite control.
The representation is unique by prime factorization. There are enough primes for
any finite program: a prime divisor of $1+\prod_i p_i$ is not among the primes already
chosen, so the list can always be extended.

To simulate INC on $r_i$, empty $A$ into $B$, adding $p_i$ to $B$ per decrement,
then move $B$ back to $A$. The result is $p_iA_{\rm old}$ and $B=0$.
To simulate DEC, divide $A$ by $p_i$ using the finite remainder cycle described above.
On exit $A=0$, $B=q$ and the control holds $d$, where $A_{\rm old}=p_iq+d$ and
$0\le d<p_i$. If $d=0$, the original number was divisible by $p_i$, so $r_i>0$.
Move $B$ to $A$ and take the positive branch. Its exponent of $p_i$ has decreased by
one and every other exponent is unchanged. If $d>0$, then $r_i=0$. Restore $A$ by
moving $B$ back with $p_i$ increments per unit, add $d$, and take the zero branch.
Both branches restore $B=0$. All loops terminate because their source counter decreases.
Induction gives the represented run at every simulated instruction boundary, and
no simulated instruction can conceal an infinite internal run.

This is universality with an encoding. With input $n$ in the first simulated counter
and the others zero, the packed input is $2^n$ when $p_0=2$. It is supplied to the
two-counter machine as its input; the construction does not claim that two counters
starting at $(n,0)$ can first manufacture $2^n$. That is precisely the distinction
Schroeppel's theorem in Section~\ref{sec:smallest} turns on, and the undecidability used by the
one-counter theorem in Appendix~\ref{sec:every-model} ranges over packed starts too.
For a tuple of inputs, supply its product of prime powers. If a numerical output is
needed in ordinary notation, the external decoding reads the designated prime's
exponent. A host with finitely many additional counters can perform this decoding itself
by repeated fixed-prime division.

\subsection{Why halting has no decider}

Write two responses on a card: ``If the predictor says STOP, keep going forever.
If it says KEEP GOING, stop.'' Suppose somebody offers a predictor that always
finishes and always gets a program's eventual behaviour right. Give it this very
card, with these instructions for using its answer. Work through both possible
answers. STOP makes the card keep going; KEEP GOING makes it stop. Waiting longer
for the predictor does not create a third correct answer: its promise was to
finish with one of the two.

The card can be a program because each response is a finite instruction, and
waiting for and branching on the predictor's answer are computable operations.
Here is how to include the input explicitly. Finite tape programs can be written
as finite binary words, and a program can interpret such a word by looking up
the next transition and maintaining the simulated tape.
Invalid descriptions may be assigned one fixed halting program. Suppose a terminating
algorithm $H(p,x)$ decided whether the program described by $p$ halts on $x$.
Build a program $D$ which, on input $p$, computes $H(p,p)$, loops forever if the answer
is yes, and halts if it is no. Let $d$ describe $D$. Either answer to $H(d,d)$ contradicts
what $D(d)$ does. So no such deciding algorithm exists.

Translate a tape program and its input by the stack construction and then prime packing.
Both translations are finite algorithms, and the translated two-counter run halts exactly
when the tape run halts. A decider for two-counter halting would then give $H$.
The same reduction, followed by Theorem~\ref{thm:u34-halting}, applies to the small machine.
Halting instances can be enumerated by simulating every description and input for successively
larger finite budgets. If the nonhalting instances could also be enumerated, interleaving
those two enumerations would decide each instance. This proves the impossibility underlying
the reduction formulation of Theorem~\ref{thm:u34-undecidable}.

\subsection{Seven descriptions of the same computation}
\label{sub:seven-constructions}

Here are the translations behind Theorem~\ref{thm:lift-models}. Each one keeps the
partial function, including where it runs forever. The constructions
are the familiar ones behind the results of Minsky, Shepherdson and Sturgis, and
Forster, Larchey-Wendling and their coauthors; the encodings and the reason each
translation works are given here.

\emph{Binary stacks.} A binary stack is a finite word with push and pop operations,
with a distinct empty case. Encode the empty word as 1 and pushing $d\in\{0,1\}$
as $n\mapsto2n+d$. Division by two pops it, with 1 tested before division.
Finite counters implement these operations. Conversely two stacks hold a tape as
above; encode every symbol of a larger alphabet by a fixed-length binary block.
A push or pop of that block takes a fixed finite number of stack operations, and the
control stores a partially read block. A tape simulates counters, so the translations
connect binary stack machines and counter machines in both directions.

\emph{The two decrement conventions.} Some counter machines jump on zero and
fall through after a positive decrement; the convention of this book does the reverse.
To translate a conditional decrement, send its two outcomes to separate short blocks
that unconditionally jump to the intended destinations. In this book's convention an unconditional jump is the restoring INC/DEC pair above. Under jump-on-zero that pair falls through, because the DEC finds the counter positive; there the jump is a DEC on one extra counter that nothing ever increments, so it is always zero. So either convention implements the other's instruction, with the same final values on the original counters and one finite stretch per step, and the translation into jump-on-zero spends one counter more. Relabel each target
to the beginning of its translated block. This proves the equivalence of the two Minsky
instruction conventions without identifying their individual steps.

\emph{FRACTRAN.} A program is an ordered finite list of positive rational numbers.
At a positive integer $N$, multiply by the first fraction whose product with $N$ is
an integer; halt if none applies. Counter machines simulate this by fixed-denominator
division, multiplication and an ordered scan of the finite list.
For the reverse translation, assign distinct primes $p_r$ to counters and $q_i$ to
control states. Represent a counter configuration by
$q_i\prod_r p_r^{n_r}$. Every translated instruction uses a fresh intermediate state
prime, so its state marker never cancels out of a fraction.
For INC $r$ going to $j$, use
$t_i p_r/q_i$, followed by the transfer $q_j/t_i$.
Try it with a counter represented by factors of 2, a starting line marked by 5,
a destination marked by 7 and an intermediate marker 11. Ten beans, $5\cdot2$,
mean line 5 with counter value 1. The fraction $22/5$ takes ten to 44; $7/11$
takes 44 to 28. Divide the result by its line marker 7: four remains, two factors
of 2. The line has changed and the counter has risen by one. The intermediate 11
keeps the two parts of the instruction separate.
For DEC $r$ with positive target $j$ and zero target $k$, use, in this priority order,
\[
\frac{t_i}{q_i p_r},\qquad \frac{u_i}{q_i},
\]
with transfers $q_j/t_i$ and $q_k/u_i$. All state and intermediate primes are distinct
from the counter primes. Exactly one state marker is present. Fractions belonging
to other states cannot apply. At state $i$, the first decrement fraction applies exactly
when $n_r>0$; otherwise the second applies. The next fraction transfers the intermediate
marker to the destination marker. A halting state has no outgoing fraction. Each
counter step then takes two fraction steps, and the prime exponents preserve all
counter values. Read the output as the exponent of the output counter's prime.

\emph{Partial recursive functions.} Compare two orders. ``Try the first ten tickets''
gives you a stopping point before you begin. ``Try tickets in order until one
passes'' does not. If each test finishes, the second order can still keep you
working forever because every ticket fails. If one test itself never finishes,
you never reach the next ticket. Those are the differences between a bounded
loop and the search used here.
Start with zero, successor and projections;
close under composition, primitive recursion and minimization. In minimization,
search $0,1,2,\ldots$ for the first argument where the tested function returns zero;
every earlier test must terminate with a positive result. A counter program implements
the initial functions directly. Disjoint scratch counters preserve arguments across
composition. For primitive recursion it keeps the recursion index and accumulated result,
increments the index, and runs the step program until it reaches the supplied bound.
For minimization it runs the tests successively; an undefined test leaves the run there
forever, as the definition requires. Induction on these constructions supplies a finite
counter program for every partial recursive function.

Conversely encode a counter configuration by nested pairs
$\langle m,n\rangle=(2n+1)2^m$, with a separate finite code for the line.
Addition, multiplication, powers, bounded division and bounded search are primitive
recursive: their loops have a supplied natural bound and can be written as recursion
on that bound. Pair decoding is bounded by the encoded positive number. The next
configuration and the halt test are then primitive recursive for any fixed finite
program. The configuration after $t$ steps is primitive recursion on $t$, keeping halted
configurations fixed. Minimize the total halt test over $t$, then read the output counter.
This is the program's partial function, and it is undefined exactly when the program
never halts.

\emph{The language $L$.} Fold two instruction cards shut. One says ``put three beans
in the output cup and stop''; the other says ``keep adding beans forever.'' A zero
input is to choose the first card. If your rule makes you carry out both cards
before choosing, you never deliver the three beans. You have to keep the
instructions folded until the choice is made. An abstraction does that job in
weak call-by-value: its body waits until the abstraction is applied.
The Scott numeral definitions and the delayed recursion
construction in Appendix~\ref{sec:quoting-construction} give zero, successor, predecessor,
case distinction and guarded recursion as terms. A case on a numeral $n$ can be delayed
explicitly as
\[
n\,(\lambda d.u)\,(\lambda p.\lambda d.(v\,p))\,(\lambda a.a):
\]
at zero it evaluates $u$, and at a successor it evaluates $v$ on the predecessor.
Both branches are values until selected, so call-by-value does not evaluate the unused
branch. Composition is application; primitive recursion repeats a step on the predecessor;
minimization is a guarded search that evaluates each test before moving to its successor.
Build the term in the same order as the function: begin with its initial
functions, then use the corresponding term operation each time composition,
recursion or minimization is used. Each operation consumes the results of the
already built pieces exactly as the function definition requires. On numerical
inputs, a computation that finishes therefore returns the Scott numeral of its
answer. If a required subcomputation never finishes, evaluation waits for it;
if minimization never finds a successful test, its search keeps going. In either
case there is no resulting value.

In the other direction, encode a finite term as a finite string or as the nested-pair
code of its syntax tree. Reading the outer constructor, finding a variable index,
renaming bound indices and substituting a finite term are terminating traversals of
finite trees. A tape algorithm performs them with an explicit stack of unvisited children.
For weak call-by-value evaluation it evaluates an application's function and argument,
then substitutes the argument value in the abstraction's body; it never reduces inside
an unapplied abstraction. Each operation is a finite tape computation. Induction on
a terminating evaluation gives a terminating tape run with its value; reading a terminating
run backwards gives the evaluation tree. Scott numerals are decoded by counting successor
constructors. So an $L$ computation of a numerical relation is a tape computation of it.

The tape and counter translations, the two counter conventions, binary stacks, FRACTRAN,
partial recursive functions and $L$ now lie in one connected family of translations.
Composing the translations preserves their input/output relation. None of them keeps time,
space, or a ledger of what was paid just because it keeps a partial function. Where I need a ledger
kept, I build that in on purpose, as the universal program of
Appendix~\ref{sec:universal-blocks} does; time and space I don't claim at all.

\subsection{Searching for a finite computation}

Give yourself one token and run one instruction of the program you want to test.
If it has not stopped, reset its starting arrangement, give yourself two tokens
and try again. Then three. A program that stops after seven instructions is found
on the seven-token attempt. Seeing six attempts fail says only that it has not
stopped within six instructions; it does not tell you it runs forever.

The evaluators used in Appendix~\ref{sec:small-diagonal} work by that search. A program
and a natural budget determine a finite simulation. Enumerate budgets
$0,1,2,\ldots$, and stop at the first one that exhibits a halting run, returning its output.
If no such run exists, the search continues forever. Encoding a finite program, decoding
an instruction, and performing one step use the finite operations above, so this search
is itself a partial recursive function and has a finite counter program.
The host can execute that program instruction for instruction, using as many registers
as the construction requires. Its computation uses only INC, DEC and HALT, so it adds nothing to the ledger and leaves the fact table, channel and flag alone. This supplies the evaluators and
program transformations used in the host's s-m-n and fixed-point arguments.


\section{The universal program in blocks}

Put the guest's instructions on one sheet and the host's on another, with a finger
marking the current line on each. The guest takes one instruction. The host takes
as many of its smaller instructions as it needs, then returns to the head of its
sheet. Compare the two arrangements there: counts, live facts, commitment, flag
and ledger. The host is allowed to do more work between comparisons. It is not
allowed to lose a fact, revive a dead one, or charge another unit for its bookkeeping.
The lists below let you carry out that comparison instruction by instruction.

Here is the finite list of instructions that is $U$. A block is a list, not an extra
instruction: replace every block name by its list, place the lists consecutively, and
resolve each label to its instruction number. The first instruction is at address 1.
Within a block, $o$ is its first address and $o+k$ its offset $k$. Unless a branch says
otherwise, execution continues into the next block. The only final instructions are
INC, DEC, HALT, CHECK, COMMIT and CERTIFY, with the meanings and prices of
Definition~\ref{def:u34-step} on the host's registers (Definition~\ref{def:u34-multi}).

\subsection{Arithmetic blocks}

Write $J(a,p)$ for the two-instruction list INC $a$; DEC $a$ $p$.
It jumps to $p$ and restores the value of $a$.
The following lists specify the other arithmetic blocks. Semicolons concatenate lists;
a named block in a list is placed at the next available address. A target such as
$o+3$ refers to the enclosing block's start, not the nested block's start.

\begingroup\small
\begin{longtable}{@{}>{\raggedright\arraybackslash}p{0.19\linewidth}>{\raggedright\arraybackslash}p{0.72\linewidth}@{}}
\toprule
Block & Expansion at address $o$ \\
\midrule\endhead
$Z(x)$ & DEC $x$ $o$ \\
$J_0(x,p)$ & DEC $x$ $(o+3)$; $J(x,p)$; INC $x$ \\
$M(x,y)$ & $J(x,o+3)$; INC $y$; DEC $x$ $(o+2)$ \\
$M_2(x,y,z)$ & $J(x,o+4)$; INC $y$; INC $z$; DEC $x$ $(o+2)$ \\
$C(x,y,t)$ & $Z(y)$; $Z(t)$; $M_2(x,y,t)$; $M(t,x)$ \\
$H(a,x,p)$ & $Z(a)$; $M(x,a)$; $J(a,o+8)$; INC $x$; DEC $a$ $(o+11)$; $J(a,p)$; DEC $a$ $(o+7)$ \\
$P(a,x,y)$ & $Z(a)$; $M_2(x,a,a)$; INC $a$; $M(a,x)$; $J(a,o+22)$; $M_2(x,a,a)$; $M(a,x)$; DEC $y$ $(o+13)$ \\
$N(a,x,y)$ & $J_0(x,0)$; $Z(y)$; $J(a,o+8)$; INC $y$; $H(a,x,o+7)$ \\
\bottomrule
\end{longtable}
\endgroup

Their lengths, in the order $Z,J_0,M,M_2,C,H,P,N$, are
$1,4,4,5,11,12,23,20$. The source and destination of a move are distinct;
$M_2$ may have equal destinations, as in $P$. In $C$ the three registers are distinct;
in $H$, $a\ne x$; in $P$ and $N$, all three are distinct.

Each claim about these blocks follows from a decreasing counter. $Z$ decrements to
zero. $M$ adds the old source to the destination and empties the source; $M_2$ does it
to both destinations. So $C$ leaves $x$ unchanged, writes its old value to
$y$, and leaves $t=0$. $J_0$ tests zero without changing the value: a positive
counter is decremented and then restored, while zero takes the jump.

In $H$, first move the old $x$ into $a$. Each complete pass removes two from $a$ and
adds one to $x$. So it leaves $a=0$ and $x=\lfloor x_{\rm old}/2\rfloor$;
an even old value jumps to $p$, an odd one falls through. In $P$, the first four
blocks make $x=2x_{\rm old}+1$ and $a=0$. Each subsequent pass doubles $x$ while
decrementing $y$, so it ends at $x=(2x_{\rm old}+1)2^{y_{\rm old}}$, with $a=y=0$.
In $N$, zero jumps to address 0, which holds no instruction, so the run stops. A positive number is repeatedly halved
until odd, counting the even divisions in $y$. So
$N$ sends $x=\langle m,n\rangle$ to $x=n$, $y=m$, $a=0$.
These invariants prove both the returned values and termination on the stated inputs.

\subsection{Selection, comparison and fetching}

The selector $D(x,t;p_0,\ldots,p_{k-1})$ has length $3k$.
For $i=0,\ldots,k-1$, place at offset $3i$ the three instructions
DEC $x$ $(o+3i+3)$; $J(t,p_i)$.
If the entering value is $i<k$, it takes target $p_i$ with $x=0$.
If it is at least $k$, it falls through with $x$ reduced by $k$.
This follows by counting the decrements before the first zero.

The equality block $E(x,y;t_1,t_2,u;p)$ first makes copies
$C(x,t_1,u)$ and $C(y,t_2,u)$, occupying offsets 0 to 21. Its remaining list is
\[
\begin{array}{r l}
22:&\mathrm{DEC}\ t_1\ (o+29),\\
23:&J_0(t_2,p),\\
27:&J(u,o+32),\\
29:&\mathrm{DEC}\ t_2\ (o+22),\\
30:&J(u,o+32),\\
32:&Z(t_1);\ Z(t_2).
\end{array}
\]
The listed offsets count expanded instructions, so the block has length 34.
The temporary registers are distinct from each other and from $x,y$.
The two copies count down together. If they reach zero together, the block jumps to
$p$; if one reaches zero first, it clears both and falls through. So it preserves
$x,y$, leaves its temporary counters zero at either exit, and branches exactly on equality.

For fetching, define $F(q,l;w,k,h,a,t;p_0,p_e)$ by the list
\[
\begin{array}{r l}
0:&C(q,w,t),\\
11:&C(l,k,t),\\
22:&J_0(k,p_0),\\
26:&\mathrm{DEC}\ k\ (o+27),\\
27:&J_0(w,p_e),\\
31:&N(a,w,h),\\
51:&J_0(k,o+56),\\
55:&\mathrm{DEC}\ k\ (o+27).
\end{array}
\]
Its length is 56. Register $q$ holds a list encoded by 0 for empty and
$\langle\text{head},\text{tail}\rangle$ for a nonempty list, and $l$ is a line number.
The scratch registers are distinct and separate from $q,l$. A zero line takes $p_0$.
Otherwise decrementing the line copy once makes $k=l-1$. Each pass removes the next
list head into $h$ and its tail into $w$. If $k=0$, that head is the selected
instruction and the block exits; otherwise it decrements $k$ and repeats.
An exhausted list takes $p_e$. The original program and line values are restored by
the copies. Every call terminates: each complete pass decreases $k$, or exits.

\subsection{Registers and instruction codes}

Put $T_i=i+4$ for $0\le i<10$. For bank $c\in\{0,1\}$ and slot $0\le k<16$, put
\[
R_c=c,\quad Q_c=14+c,\quad V_{c,k}=16+16c+k,\quad S_{c,k}=48+16c+k.
\]
Register 2 is the program, register 3 its line, and register 80 is the dead end.
$V$ are the mirrors, $S$ the slots, and $Q$ the counts of successful checks in each bank.
Property codes are $0$ for zero, $1$ for even, and $n+2$ for at least $n$.
The code of an instruction is the pair of its opcode and operand:
INC has opcode 0 and operand $c$; DEC has opcode 1 and operand $\langle c,j\rangle$;
HALT has opcode 2 and operand 0; CHECK has opcode 3 and operand $\langle c,p\rangle$;
COMMIT has opcode 4 and that same operand form; CERTIFY has opcode 5 and operand 0.
A program is their list code. These rules specify the loaded number in register 2.

A bump of slot $S$ and mirror $V$ is INC $S$; DEC $S$ to the following instruction;
$Z(V)$. It has length 3. It leaves the slot's value alone, increases its version by
two, and empties the mirror. Let $B_c$ concatenate the sixteen bumps of bank $c$,
in increasing slot order; its length is 48.

\subsection{The fifteen blocks of \texorpdfstring{$U$}{U}}

The following table fixes the order, start address and expanded length of every major
block. Labels in the definitions below mean these addresses; $H$ without arguments denotes the head at 1
and $K$ the halt block at 98.

\begingroup\small
\begin{longtable}{@{}>{\raggedright\arraybackslash}p{0.46\linewidth}rr@{}}
\toprule
Block & Start & Length \\
\midrule\endhead
Head & 1 & 97 \\
Halt & 98 & 11 \\
INC bank selection & 109 & 8 \\
INC bank 0 & 117 & 52 \\
INC bank 1 & 169 & 52 \\
DEC bank selection & 221 & 28 \\
DEC bank 0 & 249 & 60 \\
DEC bank 1 & 309 & 60 \\
CHECK bank selection & 369 & 28 \\
CHECK bank 0 & 397 & 1022 \\
CHECK bank 1 & 1419 & 1022 \\
COMMIT bank selection & 2441 & 28 \\
COMMIT bank 0 & 2469 & 628 \\
COMMIT bank 1 & 3097 & 628 \\
CERTIFY & 3725 & 4 \\
\bottomrule
\end{longtable}
\endgroup

The head is
$F(2,3;T_0,T_1,T_2,T_3,T_4;K,K)$;
$Z(T_0)$; $N(T_3,T_2,T_5)$;
$D(T_5,T_4;109,\allowbreak 221,\allowbreak 98,\allowbreak 369,\allowbreak 2441,\allowbreak 3725)$; $J(T_4,K)$.
Fetching supplies an instruction code in $T_2$. Unpacking leaves its operand in $T_2$
and its opcode in $T_5$. The six-way selector chooses its handler. The final jump
handles an invalid opcode, though a loaded program has none.
The halt block is $Z(T_0);\cdots;Z(T_9)$ followed by HALT.

INC's bank selector is $D(T_2,T_4;117,169);J(T_4,K)$.
Each other bank selector first runs $N(T_3,T_2,T_6)$, leaving the property code or
jump target in $T_2$ and the bank in $T_6$, then selects its two bank blocks with
$D(T_6,T_4;\text{bank 0},\text{bank 1});J(T_4,K)$.

The INC block for bank $c$ is INC $R_c$; $B_c$; INC 3; $J(T_4,H)$.
The DEC block, starting at $o$, is
DEC $R_c$ $(o+5)$; INC 3; $Z(T_2)$; $J(T_4,H)$;
$B_c$; $Z(3)$; $M(T_2,3)$; $J(T_4,H)$.
A zero counter follows the first return, advancing the guest line without bumping
slots. A positive counter jumps over it, bumps the bank and installs the target line.

A CHECK slot block for $(c,k)$ is the following 60 instructions after expansion:
\[
\begin{aligned}
&C(R_c,T_8,T_4);\ C(T_2,T_9,T_4);\ P(T_3,T_8,T_9);\ M(T_8,S_{c,k});\\
&\mathrm{CHECK}\ \mathsf{Slot}\ S_{c,k};\ Z(V_{c,k});\ M(T_2,V_{c,k});\\
&\mathrm{INC}\ V_{c,k};\ \mathrm{INC}\ Q_c;\ \mathrm{INC}\ 3;\ J(T_4,H).
\end{aligned}
\]
The CHECK bank block at $o$ is $C(Q_c,T_7,T_4)$, then a selector on $T_7$ whose
sixteen targets are $o+62+60k$, then CHECK $\mathsf{Slot}$ 80 and $J(T_4,K)$,
then the sixteen slot blocks. Its length is
$11+48+1+2+16\cdot60=1022$.
The selected unused slot starts at zero, so the move into it installs exactly the
pair of property code and guest value. The host's own table has capacity sixteen:
if it is full, its CHECK traps even when this bank's count is below sixteen.
So the capacity is shared between the banks, just as it is in the guest. A count of sixteen in one bank already fills the table, so the dead-end check and the full-table trap are the same case.

A COMMIT slot block is COMMIT $\mathsf{Slot}$ $S_{c,k}$;
$Z(T_2)$; INC 3; $J(T_4,H)$, of length 5.
The COMMIT bank block at $o$ first increments $T_2$, making the wanted mirror value
one more than the property code. It then runs, for $k=0,\ldots,15$,
$E(V_{c,k},T_2;T_7,T_8,T_4;o+548+5k)$.
After those comparisons come COMMIT $\mathsf{Slot}$ 80 and $J(T_4,K)$, followed
by the sixteen slot blocks. Its length is $1+16\cdot34+1+2+16\cdot5=628$.
The comparison reads the mirror, leaving the slot untouched. Its first match
commits that slot's current fact; absence of a match reaches the dead end and traps.
The CERTIFY block is CERTIFY; INC 3; $J(T_4,H)$, of length 4.

These expansions determine every instruction and target. Adding the fifteen lengths
gives 3728. There are sixteen ordinary CHECK sites and one dead-end CHECK in each
bank, sixteen ordinary COMMIT sites and one dead-end COMMIT in each bank, and one
CERTIFY site: $2\cdot17+2\cdot17+1=69$ paid sites. Every other instruction costs nothing.

\subsection{Why the blocks simulate a guest}

Run a case that catches more than a wrong sum. Start the guest with four beans
in $A$ and run CHECK Even $A$. Put its fact on a card saying ``even, version 0.''
The host copies the value into an unused slot as the pair (Even, 4), whose code
is $(2\cdot4+1)2^1=18$, and checks that slot. Put one unit on each ledger.

Now let the guest add one bean and take it away. Its pile is four again, but its
version is 2. The old fact still says version 0. On the host, each of those guest
writes bumps the slot by one and back: 18 to 19 to 18. The slot's value ends where
it began, while its version rises and its mirror is cleared. Try COMMIT Even $A$
using the old check. Both machines must refuse, even though four is even. Both
pay one more unit for that failed attempt. A translation that merely compares
the piles at the end would miss the failure this experiment is designed to catch.

Start again and leave out the two bean moves. CHECK, COMMIT, CERTIFY must succeed
on both sheets, raising both flags for exactly three units. The longer list of host instructions adds nothing to the ledger for copying or looking up a line. These
two runs test the reason for keeping both values and versions. The argument for
all runs has to keep the following correspondence at every return to the head.

At the head, registers 0, 1 and 3 give the guest counters and line. Each successful
guest CHECK has one host slot fact holding its property and checked value. Its mirror
is the property code plus one while the guest fact is current, and zero after the
guest counter has been written. The matching host slot fact is current in exactly
the first case. Count registers say how many slots have been used in their bank;
unused slots and mirrors are zero. Tables have corresponding entries in the same
order, and channels name corresponding entries. The ledgers and flags agree.
Loading a guest establishes every part of this relation.

Fetching and dispatching terminate and preserve that relation's stored values.
For INC or a positive DEC, the guest writes its counter; the selected bank's bumps
invalidate exactly its former slot facts and clear their mirrors. A zero DEC writes
nothing and bumps nothing. These cases add nothing to either ledger.
For CHECK, an unused slot receives the exact pair tested by the guest. Both checkers
answer the same, and the tables have the same length. A passing check appends the
corresponding fact and updates its mirror and count. A false check or a full table
traps both machines, after the same single paid act. For COMMIT, a nonzero mirror
matches exactly a current guest fact with the requested property and counter. Its
host slot has not been read or changed, so the corresponding COMMIT succeeds.
If none matches, both machines trap. Again exactly one unit is paid. For CERTIFY,
channels are either both empty or both hold corresponding commitments, so the
instruction either traps both or raises both flags, for one unit.

All arithmetic between these paid acts terminates by the block invariants. A normal
case returns to the head; a halting or trapped case stops. No arithmetic block has a
paid instruction. This proves the one-step correspondence and bounds every
intermediate host ledger between the guest's old and new ledgers. Induction proves
the run correspondence of Theorems~\ref{thm:u34-U-sim} and~\ref{thm:u34-U-flag}.
It also proves the reverse halting implication: an infinite guest run entails
infinitely many nonempty host stretches, and a running guest instruction cannot
produce a stopped host except in its own halting or trapping case.

\subsection{Adding PAY}

For $U_P$, use opcode 6 with operand 0 for PAY, append a PAY handler to the six handlers,
and append its block after CERTIFY. The block is PAY; INC 3; $J(T_4,H)$, of length 4.
The head's selector now has seven entries and is three instructions longer. Every
block after the head then starts three addresses later. Resolve targets to those
new addresses; the head itself remains at address 1.
The new PAY block starts at 3732. The total length is $3728+3+4=3735$ and there
are $69+1=70$ paid sites. The interpretation of $\mathsf{Slot}$ uses the priced
host's property language from Appendix~\ref{sec:presented-universal}; the arithmetic
blocks and their invariant do not depend on the particular meaning of a property code.
A PAY guest step is one PAY host step surrounded by arithmetic that costs nothing, so the
same ledger argument applies.


\section{Every assumption in one place}

This appendix gathers the conditions used throughout the book and shows what each supplies.
Use it to follow the connections between the results: which count needs finiteness, which
record needs permanence, and which charge follows from which price rule. Each theorem also
states its own hypotheses where the argument is made.

The labels A2 to A6 in this appendix are shorthand for the conditions listed below.
They are not additional claims: each label names the rule written beside it. There is no A1: A2
is the toll's name in the mathematical specification, and the others follow it.

\subsection{Systems and their conditions}

\begin{definition}[Thiele machine with a two-place record]\label{def:app-certsys}
A Thiele machine (Definition~\ref{def:thiele}) has a record in any order. The results of
this appendix are stated on the two-place record of Definition~\ref{def:certsys}: a set of
states $S$, a set of moves $I$, a step function
$\mathrm{step}: S\times I\to S$, a cost $\mathrm{cost}:I\to\mathbb{N}$, and a reading
$\mathrm{cert}:S\to\{\mathrm{no},\mathrm{yes}\}$, satisfying the toll: every move that
turns the reading from no to yes costs at least one. This is the condition of
Section~\ref{sec:toll}, called A2 in the table below. Every result stated ``for any
Thiele machine'' in this appendix assumes it.
\end{definition}

\begin{definition}[Quantitative Thiele machine]\label{def:app-qcs}
A quantitative Thiele machine is a Thiele machine with a \emph{witness},
a function $W$ from states to $\mathbb{N}$, and a threshold $N \in \mathbb{N}$, satisfying:
(A3) $W(s) + \mathrm{cost}(i) \ge W(\mathrm{step}(s, i))$, the witness grows by at most
what the move cost; and (A5) a state that reads yes has $W(s) \ge N$.
An optional additional condition is (A4) $W(s) \le W(\mathrm{step}(s, i))$.
\end{definition}

\begin{theorem}[The floor raised to a threshold]\label{thm:app-quantitative}
In a quantitative Thiele machine, every run from a state with witness $w$ to a state that reads yes has total cost at least $N \mathbin{\dot{-}} w$; from a state with witness $0$ (this starting condition is called A6) that is at least $N$.
\end{theorem}

\begin{proof}
Put the witness readings in a row: $w_0=w,w_1,\ldots,w_n$.
Beside the move from $w_j$ to $w_{j+1}$ write its charge $c_j$.
A3 says that move can add no more than $c_j$ to the witness:
$w_{j+1}-w_j\le c_j$. Add these inequalities. Each intermediate
reading appears once with a plus and once with a minus, so only the last
reading minus the first survives:
\[
 w_n-w\le c_0+\cdots+c_{n-1}.
\]
At the end the reading says yes, so A5 puts $w_n$ at least at $N$.
The bill is therefore at least $N-w$. If the witness started above $N$,
that subtraction is negative and supplies no positive floor; the bill is
still at least zero. That is what the truncated subtraction
$N\mathbin{\dot{-}}w$ says. At $w=0$ the floor is $N$.
For an empty run the same argument reads $w_n=w$ and bill zero.
Nothing here required the witness to rise on every move (A4), or used the toll.
\end{proof}

A3 and A4 bound opposite
directions, and neither implies the other: witness 1, cost 0, witness 0 satisfies A3 and
breaks A4, and witness 0, cost 0, witness 1 does the reverse. These conditions describe a general setting.
The small machine's own floor of three is proved directly (Observation~\ref{thm:u34-min-price})
and never goes through it.

The table below puts the labels side by side, so none of them is left without a meaning.

\begingroup\small
\begin{longtable}{@{}>{\raggedright\arraybackslash}p{0.2\linewidth}>{\raggedright\arraybackslash}p{0.75\linewidth}@{}}
\toprule
\textbf{Label} & \textbf{Condition} \\
\midrule
\endhead
\bottomrule
\endlastfoot
A2 & the toll \\
A3 & the witness grows by at most the cost \\
A4 & the witness never drops \\
A5 & yes needs a witness at the threshold \\
A6 & the run starts with witness zero \\
\end{longtable}
\endgroup

\subsection{The premises the derivations take}

These are definitions used as hypotheses. Nothing proves them in general; each theorem that
uses one says so in its statement.

\begingroup\small
\begin{longtable}{@{}>{\raggedright\arraybackslash}p{0.22\linewidth}>{\raggedright\arraybackslash}p{0.35\linewidth}>{\raggedright\arraybackslash}p{0.35\linewidth}@{}}
\toprule
\textbf{Premise} & \textbf{What it says} & \textbf{Where it's used} \\
\midrule
\endhead
\bottomrule
\endlastfoot
One next state & every state has exactly one next state, and states can be told apart & the step law, Theorem~\ref{thm:step-law}, and merge or a fresh state, Proposition~\ref{prop:heat-or-room}(1); no finiteness \\
Finite & a list with no repeats contains every state & the permanent-flip theorem, Theorem~\ref{thm:perm-flip-merges}, where only the states reading yes need to be finitely many, and what's built on it: merge or revoke, the squeeze bound, the entropy road and its heat floor, the certainty bound and the direction (Theorems~\ref{thm:merge-or-revoke}, \ref{thm:squeeze}, \ref{thm:u04-entropy-toll}, \ref{thm:certainty-bill} and~\ref{thm:dir-merge}); the noisy and unitary steps; the exact count of states, Proposition~\ref{prop:heat-or-room}(2), for bounded memories; recurrence, Theorem~\ref{thm:reversible-keeps-nothing} \\
Permanent & no move turns the reading from yes to no & the permanent-flip theorem; the toll derived, Observation~\ref{thm:toll-derived}; the squeeze bound and the entropy road; a raise with no merge, Proposition~\ref{prop:raise-needs-room}; merge or a fresh state; the latch, as (R4); the direction, as growth \\
Merge pricing & every move that merges costs at least one & the toll derived, Section~\ref{sec:a2-not-arbitrary}; forced pricing, Theorem~\ref{thm:forced-iff-merge}; the direction, Theorem~\ref{thm:dir-merge}(3) \\
Halving price & each unit of cost pays for at most one halving of distinct states & the squeeze bound; narrowing along a run, Section~\ref{sec:narrowing} \\
Entropy pricing & on a finite machine, each move costs at least the bits it removes from any spread & the third road to the toll; the certainty bound, Theorem~\ref{thm:certainty-bill} \\
Landauer's principle & for a stated physical memory whose logical states are alike in their own entropy, removing $\Delta H$ bits gives off mean heat $\ge k_B T \ln 2 \cdot \Delta H$ & the heat floor; with $k_B T \ge 0$ (or $> 0$) as a separate hypothesis on each theorem \\
Recursion (generic) & every transformation the substrate can represent has a fixed point & the diagonal, Section~\ref{sec:diagonal} \\
Commitment contract & a commitment bit is set only when the claim holds, and whenever it holds & the commitment escape, Section~\ref{sec:windows} \\
Computably presented & four programs compute the driver, step, price and reading on valid codes; the reading program also halts and returns false on junk codes & the universal host for presented machines, Appendix~\ref{sec:presented-universal} \\
Clean start & an empty fact table, an empty channel and the flag down & every provenance theorem about the small machine, Section~\ref{sec:smallest}; the floor of three needs less, and that section says exactly what \\
Record-layer merge pricing & a move that sends two different states with the same base part to one state costs at least one & the toll on the example machine, Theorem~\ref{thm:record-merge-toll} \\
Structural dissipation pricing & every collapsing instruction dissipates at least one unit, and dissipation is at most cost & Section~\ref{sec:landauer}, Premise~\ref{prem:u24-landauer-pair} \\
Covering tree, paid steps & a decision tree whose leaves cover the narrowing, and at least as many steps charged at least one as the tree is deep & entitlement, Theorems~\ref{thm:u08-count}, \ref{thm:ent-cs} and~\ref{thm:ent-machine} \\
Share kept by a test & each passed test keeps at least a $2^{-k}$ share of the candidates, for a $k$ fixed before the run; every run meets it for $k$ its own biggest test & entitlement made by the run, Theorems~\ref{thm:ent-run} and~\ref{thm:ent-run-machine} \\
Lift interface & a claim language with an exact equality test, a checker proved to say yes exactly when the claim holds, a claim true at one loaded state and false at another, room for one fact & the lift, Theorem~\ref{thm:lift} \\

Representable flip & the substrate can represent the flip of the candidate decider & the diagonal, Theorem~\ref{thm:u12-no-inside-decider} \\
Quorum intersection; charge predicate & as transcribed in Section~\ref{sec:models} (the setting also carries at most one parent, which no proof uses) & Theorems~\ref{thm:u13-accountable-safety} and~\ref{thm:u13-gas-exactness} \\
\end{longtable}
\endgroup

Four remarks on the table. First, the recursion premise holds by construction for the nat
family and by proof for L (Section~\ref{sec:diagonal}) and for the programs of the small machine's
host (Appendix~\ref{sec:small-diagonal}).
Second, the substrate used for the diagonal also carries a different rule from the toll: the counter
never goes down on a step. Neither implies the other, and the diagonal uses neither. Third, merge pricing is an abstract lower bound on every merging instruction.
It is the pairwise restriction of worst-case entropy pricing, which is equivalent
to the halving price. Landauer's premise instead bounds heat from the uncertainty
actually removed under its thermal conditions. A known starting state supplies no
positive floor from that premise (Theorem~\ref{thm:known-state}). Preparing or erasing an uncertain record of that start would be a separate
physical process with its own premises; knowing the start alone proves no
payment. Under the fixed-protocol entropy-balance assumptions in
Section~\ref{sec:bits-zero}, a protocol merging $K$ inputs has a worst input
with at least $k_B\ln K$ of entropy production. This compares inputs to
one protocol rather than charging every known-state execution. Riechers, Boyd, Wimsatt and Crutchfield (2020) derive additional work bounds under
time-symmetric control of metastable memories. Their low-error bound depends on the
initial distribution and the memory's time-reversal behaviour; it is not a universal
heat charge on each logical move.

Fourth, finiteness, because it's the premise people most often think the book leans on, and it
leans on it less than it looks. The step law and the first half of merge or a fresh state need none: on any machine run by one rule, finite or endless, a run that raises a reading that stays up either merges two of its
states somewhere or takes a fresh state at every step. Finiteness is what takes the second option away. With infinitely many
states reading yes, a price list that charges only merges lets a machine certify at a cost of zero,
the machine that keeps every history, the first way out (Section~\ref{sec:a2-not-arbitrary}), by taking a fresh state at every step. That
isn't a way around what keeping takes; it's a price list that charges nothing for the fresh states. And the finite set the count needs can be very small. When the record steers nothing and every move adds its own price to the ledger, the count runs over one state of the rest of the machine, on the record's own readings, and for a flag that's a set of two. That's how the machine of Part~III, which has no end, still pays for its raise in a merge (Theorem~\ref{thm:record-merge-toll}). Applying a finite-state count to a device requires a justified finite set of distinguishable states. The abstract theorem does not supply that physical premise.




\section{The claim ledger}

This section lists everything this book claims, one line each, with what kind of claim it is and
where it lives. ``Proved'' means proved under the premises printed in that section, and nothing more.
If a line here ever says more than its section does, the section is right and the line is
wrong. ``Directly from the rules'' means the claim can be read off the stated
operations by checking their cases. ``Directly from the interface'' means its
requirements already include the claimed behavior. These claims are true, but
carry almost nothing beyond what those definitions already say. If you're counting results, count those
differently. I mark them so they don't pad the list, and where they're stated in the book they
carry the label \emph{observation}, with three exceptions: Theorem~\ref{thm:u34-toll}, part (3)
of Theorem~\ref{thm:u34-presented-surcharge}, and the parts of Theorem~\ref{thm:slice} its line names.

Ledger is a good word for it, and it's older than bookkeeping. Around 1400 a \emph{ledger} was a
book that lay permanently in one place, a big church breviary, say, from the old verb
\emph{leggen}, to lay; the account-book sense turns up in the 1500s. Staying put is
the job here too. The home sections are long and they wander, and this is the one place where
every claim lies still, so you can come back and find it where you left it.

\subsection{Claims about the model}

The model comes first. Every line in this table is about any system that has the pieces the line
names: a state and a step, plus whatever the result asks for, a reading and the toll, a
counter and a price list, a window, programs. None of them is about my machine, except the lines that run these results on it,
which say so. The table
runs over several pages, in reading order, except that a line running a result on the small
machine sits beside the result it runs.

There's a way to use it that takes five minutes. Pick the sentence in this book you trust
least. Find its line. Read the Kind column before anything else. Then go to the home section
and check that the line says no more than the section does. If it says more, you've found my
mistake, and the end of this section says what I'd like done about it.

\begingroup\small
\begin{longtable}{@{}>{\raggedright\arraybackslash}p{0.242\linewidth}>{\raggedright\arraybackslash}p{0.344\linewidth}>{\raggedright\arraybackslash}p{0.344\linewidth}@{}}
\toprule
\textbf{Claim} & \textbf{Kind} & \textbf{Home} \\
\midrule
\endhead
\bottomrule
\endlastfoot
Every run of a machine with one next state for every state takes fresh states until its first return, which lands on its start or is a merge, and after it only repeats steps already taken, so it merges at most once; the only steps that spend neither a fresh state nor a merge are on a loop, which makes nothing new & proved & Section~\ref{sec:what-a-chair-is}, Theorem~\ref{thm:step-law} \\
A step that may have no next state can raise a permanent reading with no merge; with each halted state stepping to itself, any state that steps onto a halted state shares that landing with it, a merge & proved & Section~\ref{sec:what-a-chair-is}, Theorem~\ref{thm:halting-hole} \\
Writing a 0 on a blank cell, with the write defined on written cells too, is a merge of blank with zero; a remembering state added to the cell still merges; the endless tape and the bilateral shift avoid the merge by taking a fresh state every step; one Thiele machine runs computably presented Thiele machines as guests, paying at least one unit for each rise of a guest's record, and no level codes all its own readings & proved & Section~\ref{sec:what-a-chair-is}, Theorems~\ref{thm:write-merge} and~\ref{thm:no-self-containment}; the guests, Theorems~\ref{thm:u34-presented} and~\ref{thm:parts-nesting} \\
A reversible step on finitely many states returns every run to its start and keeps nothing & proved & Section~\ref{sec:what-a-chair-is}, Theorem~\ref{thm:reversible-keeps-nothing} \\
A noisy step that preserves the even spread cannot raise a permanent reading; one that can raise such a reading strictly concentrates the even spread, raising its collision weight and leaving its Shannon entropy below the maximum & proved & Section~\ref{sec:what-a-chair-is}, Theorem~\ref{thm:noisy-chairs} \\
A finite-dimensional unitary that keeps a subspace keeps its complement, so no unitary step raises a permanent reading in finite dimension & proved; finite dimension is used once, for ``onto'' & Section~\ref{sec:what-a-chair-is}, Theorem~\ref{thm:quantum-chairs} \\
A biased walk hands its surroundings strictly positive entropy per tick at any nonzero speed, and heat times time is at least the distance squared & proved; Schnakenberg's formula taken as the model & Section~\ref{sec:what-a-chair-is}, Theorem~\ref{thm:time-bill} \\
The measured floor, the measured time-heat rate, the measured recurrence and the collapse-heat searches & cited, not reproved; the bead's erasure law reproduced by simulation & Section~\ref{sec:what-a-chair-is} \\
Every scale pays the same way: when a coarse description is exact on a set of fine states and its rule merges, the fine rule merges there, steps out of that set, or the set is endless & proved & Section~\ref{sec:what-a-chair-is}, Theorem~\ref{thm:every-scale} \\
Thiele machine: states, moves, step, cost, reading & definition & Section~\ref{sec:toll}, Definition~\ref{def:certsys} \\
The toll & choice & Section~\ref{sec:toll}, Definition~\ref{def:certsys} \\
The universal floor: no to yes costs at least one & proved (premise: the toll), directly from the rules & Section~\ref{sec:toll}, Observation~\ref{thm:universal-floor} \\
One unit per crossing: a run costs at least its number of steps from no to yes, the converse gives the toll, and the count is reached & proved (premise: the toll), directly from the rules & Section~\ref{sec:toll}, Observation~\ref{thm:per-crossing} \\
With a price per move name, no system whose flag rises keeps every run at most its number of raises and still obeys the toll & proved & Section~\ref{sec:toll} \\
The floor speaks about the world, or computing power & not claimed & Section~\ref{sec:toll} \\
A computation is in scope just by being one & not claimed; someone has to supply the five parts and prove the toll & Section~\ref{sec:toll} \\
On a finite machine, a permanent flip merges & definition; proved (no price premise) & Section~\ref{sec:a2-not-arbitrary}, Definition~\ref{def:fin-perm-inj}, Theorem~\ref{thm:perm-flip-merges} \\
Merge pricing & assumed (premise) & Section~\ref{sec:a2-not-arbitrary}, Premise~\ref{prem:merge-pricing} \\
Finite, permanent and merge pricing give the toll & proved (premise: merge pricing), directly from the rules & Section~\ref{sec:a2-not-arbitrary}, Observation~\ref{thm:toll-derived} \\
On any machine, a raise by a move that never turns the reading off and merges nothing sends the yes-states one-to-one into a part of themselves that leaves one out (Hilbert's hotel), so it needs infinitely many yes-states; with finitely many, that raise merges & proved (premise: the move never turns the reading off; nothing else beyond telling states apart) & Section~\ref{sec:a2-not-arbitrary}, Proposition~\ref{prop:raise-needs-room} \\
Anything kept is paid for in a merge or in fresh states (finite-time dissipation bounds for specified thermal stochastic models are cited physics, proved for the walker): any run of one rule that raises a record and keeps it either merges two of its own states or is in a new state at every step, so on $M$ states a non-merging run that raises it has fewer than $M$ steps, and a bounded memory only delays the merge & proved; Bennett's trade-off cited, not reproved & Section~\ref{sec:a2-not-arbitrary}, Proposition~\ref{prop:heat-or-room} \\
Retention bound $M \ge 1 + (T+1)(N-1)$, exact, and $M \ge N + T(N-L)$ with $L$ fixed places & proved; the fibre form by the same count, on the page & Section~\ref{sec:a2-not-arbitrary}, Proposition~\ref{prop:heat-or-room} \\
On a finite machine, a flip merges or revokes & proved & Section~\ref{sec:a2-not-arbitrary}, Theorem~\ref{thm:merge-or-revoke} \\
Forced to be priced exactly when it merges & definition; proved & Section~\ref{sec:a2-not-arbitrary}, Definition~\ref{def:forced-priced}, Theorem~\ref{thm:forced-iff-merge} \\
On a finite machine, permanent records are forced; the converse fails & proved (the three cards) & Section~\ref{sec:a2-not-arbitrary}, Corollary~\ref{cor:permanent-forced} \\
Each premise of the toll has a machine that walks past it & proved (the strip, which pays by taking states it has never used; the coin; the stamp); certifying at zero cost on the ledger is proved too, and on the strip it means the price list charges nothing for fresh states & Section~\ref{sec:a2-not-arbitrary} \\
Landauer's principle is true of the world & not claimed; a named premise, with its critics named & Section~\ref{sec:a2-not-arbitrary} \\
Halving price & assumed (premise) & Section~\ref{sec:a2-not-arbitrary}, Premise~\ref{prem:halving} \\
The squeeze bound & proved (premise: halving price) & Section~\ref{sec:a2-not-arbitrary}, Theorem~\ref{thm:squeeze} \\
The halving price holds exactly when no pile-up exceeds $2^{\mathrm{cost}}$; the eight-state price list is the least that meets it; the small machine is not halving-priced & proved (finite machines; Landauer's price premise stays named) & Section~\ref{sec:a2-not-arbitrary} \\
\end{longtable}
\endgroup

\begingroup\small
\begin{longtable}{@{}>{\raggedright\arraybackslash}p{0.242\linewidth}>{\raggedright\arraybackslash}p{0.344\linewidth}>{\raggedright\arraybackslash}p{0.344\linewidth}@{}}
\toprule
\textbf{Claim} & \textbf{Kind} & \textbf{Home} \\
\midrule
\endhead
\bottomrule
\endlastfoot
Forgetting nothing removes no bits & proved & Section~\ref{sec:a2-not-arbitrary}, Theorem~\ref{thm:u04-injective-entropy} \\
A permanent flip from the even spread removes at least $\log_2 \frac{m+k}{m}$ bits and at least $\frac{2k}{m+k}$ bits & proved & Section~\ref{sec:a2-not-arbitrary}, Theorem~\ref{thm:entropy-ceiling} \\
Landauer's principle & assumed (premise) & Section~\ref{sec:a2-not-arbitrary}, Premise~\ref{prem:landauer} \\
Heat floor of a permanent flip & proved (premises: Landauer, $k_B T \ge 0$, the even spread) & Section~\ref{sec:a2-not-arbitrary}, Corollary~\ref{cor:heat-floor} \\
In the worst case over spreads a permanent flip removes at least a full bit, whatever the number of yes-states; from the even spread it removes at least $\log_2 \frac{m+k}{m}$ and at least $\frac{2k}{m+k}$ bits & proved & Section~\ref{sec:a2-not-arbitrary} \\
Positive heat for any full spread & proved (premises: Landauer, $k_B T > 0$) & Section~\ref{sec:a2-not-arbitrary}, Corollary~\ref{cor:u04-full-support-heat} \\
Entropy pricing gives the toll & definition; proved (premise: entropy pricing) & Section~\ref{sec:a2-not-arbitrary}, Definition~\ref{def:u04-entropy-priced}, Theorem~\ref{thm:u04-entropy-toll} \\
Entropy pricing is equivalent to the halving price; merge pricing is their weaker restriction to pairs. The largest entropy loss is $\log_2$ of the largest pile-up & proved & Section~\ref{sec:a2-not-arbitrary}, Theorem~\ref{thm:one-premise} \\
On a known state the premise sets no positive heat floor; the distinction was kept by whoever knows the state & proved; the second half is a reading & Section~\ref{sec:a2-not-arbitrary}, Theorem~\ref{thm:known-state} \\
The bits removed, one term per state & proved & Section~\ref{sec:a2-not-arbitrary}, Theorem~\ref{thm:u04-point-sum} \\
Any machine worth building is of this kind & not claimed & Section~\ref{sec:a2-not-arbitrary} \\
An eight-state machine meets every premise, so the toll holds on it outright & proved & Section~\ref{sec:a2-not-arbitrary} \\
The small machine runs the eight-state machine through a window, move for move, paying exactly its charges & proved (from live states; the setup's claim can't fail) & Section~\ref{sec:part3-on-machine}, Theorem~\ref{thm:frag-small} \\
The small machine as a whole is not merge-priced & proved & Section~\ref{sec:part3-on-machine}, Theorem~\ref{thm:frag-merge-split} \\
On the small machine CERTIFY, COMMIT and CHECK merge, every merge-priced cost obeys the toll, and record-layer merge pricing holds for its own prices and gives the toll, with no finiteness premise & definition; proved (premise: record-layer merge pricing, for the toll) & Section~\ref{sec:a2-not-arbitrary}, Premise~\ref{prem:record-merge}, Theorem~\ref{thm:record-merge-toll} \\
Permanence is the job to pick & choice & Section~\ref{sec:a2-not-arbitrary} \\
The machine's own narrowing is priced & proved (premise: halving price) & Section~\ref{sec:narrowing}, Theorem~\ref{thm:machine-narrowing} \\
An observer can learn with nothing merged: on the four-state demon, a run of cost zero under a price list meeting the halving price narrows two candidates to one, and WIPE costs at least one & definition; proved (one price list that allows it) & Section~\ref{sec:narrowing}, Definition~\ref{def:knowledge}, Theorem~\ref{thm:observer-free} \\
Learning borrows: the copy uses up a blank page and can be handed back (a second MEASURE clears it), and the page is paid for when it's wiped & proved on the four-state demon; the wipe merges & Section~\ref{sec:narrowing} \\
Learning costs nothing on every price list, or takes no fresh state & not claimed; a price list may charge it, and it uses up a blank page & Section~\ref{sec:narrowing} \\
\end{longtable}
\endgroup

\begingroup\small
\begin{longtable}{@{}>{\raggedright\arraybackslash}p{0.242\linewidth}>{\raggedright\arraybackslash}p{0.344\linewidth}>{\raggedright\arraybackslash}p{0.344\linewidth}@{}}
\toprule
\textbf{Claim} & \textbf{Kind} & \textbf{Home} \\
\midrule
\endhead
\bottomrule
\endlastfoot
Three states are the minimum for teaching at zero cost after the first look & proved & Section~\ref{sec:narrowing}, Theorem~\ref{thm:three-states} \\
The certainty bound: the uncertainty about a start splits into the part the run destroyed, at most its cost in bits, and the part the end state holds; whatever the end state can't tell you about what an observer saw was paid for in units of cost & proved (premise: entropy pricing) & Section~\ref{sec:narrowing}, Theorem~\ref{thm:certainty-bill} \\
``No Free Insight,'' the toll's first name, names a distinction made to stay & choice (the name kept, the scope stated) & Section~\ref{sec:narrowing} \\
Watching narrows like a quantum measurement & not claimed & Section~\ref{sec:narrowing} \\
Structure: relations, kept under chosen changes, carried by any stuff, relative to a frame & choice (my overlay of the sources; they never proved it together) & Section~\ref{sec:structure} \\
Turing's four parts have no slot for ``these belong together'' & not claimed; I looked and didn't find one, which is weaker & Section~\ref{sec:structure} \\
Equivalence relation, partition & definition & Section~\ref{sec:structure}, Definition~\ref{def:equiv-partition} \\
Sameness and sorting are one object & proved (textbook, Halmos) & Section~\ref{sec:structure}, Theorem~\ref{thm:equiv-partition} \\
A partition is the start, smallest by the count of names & choice & Section~\ref{sec:structure} \\
Nobody prices vouching for a sorting & not claimed; only that I found no one who does & Section~\ref{sec:structure} \\
A counter kept by a price list & definition & Section~\ref{sec:counter}, Definition~\ref{def:u06-counter} \\
Conservation & proved  & Section~\ref{sec:counter}, Theorem~\ref{thm:conservation} \\
Only one honest count & proved  & Section~\ref{sec:counter}, Theorem~\ref{thm:initiality} \\
The fold on lists is unique; it reaches states only under a condition & proved  & Section~\ref{sec:counter} \\
Initiality picks the prices, or makes units of cost physical & not claimed & Section~\ref{sec:counter} \\
The floor and a ceiling (a run never costs more than its raises) leave exactly one cost function, one unit per raise and nothing else, when the price may depend on where the step starts; the small machine is not priced that way, by design & proved & Section~\ref{sec:counter} \\
The toll is the potential method & definition; proved & Section~\ref{sec:counter}, Definition~\ref{def:potential}, Theorem~\ref{thm:potential} \\
The potential method is new here & not claimed (it's Tarjan's) & Section~\ref{sec:counter} \\
The counter is the cost-semantics frameworks' writer & proved for the writer: a run in the writer over whole numbers returns the final state and the total cost & Section~\ref{sec:counter} \\
The counter embeds the cost-semantics frameworks' type systems and their before-the-run checking of grades & not claimed; by design the grade is written by the step during the run & Section~\ref{sec:counter} \\
A window that forgets the counter can't recover it on every state; on reached states it takes two runs with one view and two counters & proved (premise: the counter can be changed on its own); the premise is needed  & Section~\ref{sec:counter} \\
A decoder needs constant fibres & proved (any types) & Section~\ref{sec:windows}, Theorem~\ref{thm:fibre} \\
No exact window price across a collision & proved (premise: a collision), directly from the rules & Section~\ref{sec:windows}, Observation~\ref{thm:shadow-price} \\
\end{longtable}
\endgroup

\begingroup\small
\begin{longtable}{@{}>{\raggedright\arraybackslash}p{0.242\linewidth}>{\raggedright\arraybackslash}p{0.344\linewidth}>{\raggedright\arraybackslash}p{0.344\linewidth}@{}}
\toprule
\textbf{Claim} & \textbf{Kind} & \textbf{Home} \\
\midrule
\endhead
\bottomrule
\endlastfoot
On the toy, no verifier on bare transcripts is sound and complete, and none that reads only the bare part of a richer transcript is either, for any pair like it & proved, over arbitrary types & Section~\ref{sec:windows} \\
Full-state, commitment and response verifiers are sound and complete for their explicit explanation relations; commitment is sound and complete exactly when binding and honesty hold & proved; verification cost is stipulated, and no proof reads it & Section~\ref{sec:windows}, Definition~\ref{def:u21-escape-relations}, Theorem~\ref{thm:u21-escapes} \\
On every certified Thiele-complete machine, no verifier on bare transcripts is sound and complete for the record, or for a ledger rise of 3 from a clean loaded state & proved & Section~\ref{sec:windows}, Theorem~\ref{thm:ver-complete} \\
The one-round ``ask'' models a cheating prover & not claimed & Section~\ref{sec:windows} \\
The window fact is mysticism, or Rice & not claimed & Section~\ref{sec:windows} \\
Entitlement: object, checked relation, paid event & definition (the book's word) & Section~\ref{sec:entitlement}, Definition~\ref{def:u08-entitlement} \\
A decision tree has at most two to its depth leaves & definition; proved & Section~\ref{sec:entitlement}, Definition~\ref{def:u08-dtree}, Lemma~\ref{lem:u08-tree-width} \\
Paid narrowing, counted & proved with the total of the charges, which equals the counter's rise by conservation; both are proved (premises: a covering tree, paid steps) & Section~\ref{sec:entitlement}, Theorem~\ref{thm:u08-count} \\
Structural entitlement over any Thiele machine: a strictly stronger test, a paid crossing, a cost that covers the index bits & proved (premises: covering tree, paid steps) & Section~\ref{sec:entitlement}, Theorem~\ref{thm:ent-cs} \\
A reading that a grouping was checked and committed sees only its partition and entitles, uniformly over keys, exactly the groupings that refine it & proved & Section~\ref{sec:entitlement}, Theorem~\ref{thm:partition-reading} \\
Structural entitlement for every certified Thiele-complete machine: stronger test, earned chain, index bits at most the record moves, ledger rise at least 3 & proved (premises: a clean start, and M1 to M6) & Section~\ref{sec:part3-on-machine}, Theorem~\ref{thm:ent-machine} \\
A floor with no supplied tree, from the run's own CHECKs & proved (answers taken at the rivals as they start) & Section~\ref{sec:part3-on-machine}, Theorem~\ref{thm:ent-questions} \\
Structural entitlement with the narrowing made by the run, over any Thiele machine whose test steps are named: every drop made by a test the run passed, the exact product count, index bits at most $k$ times the cost & proved (premises: test steps charged at least one, a prior containing the start, a $2^{-k}$ share kept by each test) & Section~\ref{sec:entitlement}, Theorem~\ref{thm:ent-run} \\
The same for every certified Thiele-complete machine: index bits at most $k$ times the passing CHECKs, at most $k$ times the ledger's rise; the rounded sum with no share assumed; the share can't be dropped & proved (premises: a clean start, the record up at the end, a prior containing the start, the share for part (4)) & Section~\ref{sec:part3-on-machine}, Theorem~\ref{thm:ent-run-machine}, Proposition~\ref{prop:ent-share-needed} \\
What the record adds to the price bound from a clean start: the earned chain and a floor of 3 & proved & Section~\ref{sec:part3-on-machine}, Theorem~\ref{thm:ent-observed} \\
The weighted count and the partition form & proved; the weights are whole numbers, and no probabilities enter & Section~\ref{sec:part3-on-machine} \\
The eight-start member, and a search through $n$ bits with up to 16 questions & proved; the 16 is the machine's fact table & Section~\ref{sec:part3-on-machine}, Proposition~\ref{prop:ent-small}, Theorem~\ref{thm:ent-search} \\
The covering hypothesis can't be dropped & proved, by a counterexample on the small machine & Section~\ref{sec:part3-on-machine} \\
The time tax: no program of zero-cost two-counter moves decides, in fewer than $m$ moves, $A \ge m$ from its final line and $B$, one paid move decides it, and three paid moves certify it & proved, for programs of INC and DEC & Section~\ref{sec:part3-on-machine}, Theorem~\ref{thm:time-tax} \\
The lists, tree and fibres are read off the run & not claimed, except in the worked members & Section~\ref{sec:part3-on-machine} \\
Entitlement is truth, or permission & not claimed & Section~\ref{sec:entitlement} \\
The count prices why candidates were dropped & not claimed; it prices the size of the drop & Section~\ref{sec:entitlement} \\
Raise, non-raising move & definition & Section~\ref{sec:insight}, Definition~\ref{def:u09-events} \\
The toll prices only the crossing, flat in depth & fact read off the definitions & Section~\ref{sec:insight}, Remark~\ref{rem:u09-three} \\
Sorting moves by whether the reading moves & choice & Section~\ref{sec:insight} \\
Nothing makes the crossing wait for a check & fact read off the definitions & Section~\ref{sec:insight}, Remark~\ref{rem:u09-no-wait} \\
A paid flag proves its claim true & not claimed & Section~\ref{sec:insight} \\
\end{longtable}
\endgroup

\begingroup\small
\begin{longtable}{@{}>{\raggedright\arraybackslash}p{0.242\linewidth}>{\raggedright\arraybackslash}p{0.344\linewidth}>{\raggedright\arraybackslash}p{0.344\linewidth}@{}}
\toprule
\textbf{Claim} & \textbf{Kind} & \textbf{Home} \\
\midrule
\endhead
\bottomrule
\endlastfoot
Base, cover, base-driven record, latch & definition & Section~\ref{sec:latch}, Definition~\ref{def:u11-base}, Definition~\ref{def:u11-honest} \\
Over any base, a base-driven record is a latch & proved, directly from the rules & Section~\ref{sec:latch}, Observation~\ref{thm:u11-latch} \\
Every reachable latch is a base-driven record & proved & Section~\ref{sec:latch}, Theorem~\ref{thm:u11-converse} \\
A finite reversible machine cannot raise a permanent reading & proved; an unbounded-history extension can, by taking fresh states & Section~\ref{sec:latch}, Theorem~\ref{thm:u11-reversible} \\
Two records are two latches & proved & Section~\ref{sec:latch}, Lemma~\ref{lem:u11-pair} \\
The toggle: driven, can switch off, no latch & proved & Section~\ref{sec:latch}, Lemma~\ref{lem:u11-toggle} \\
The timestamp, a record switched on by a hidden clock: permanent, undriven & proved & Section~\ref{sec:latch}, Lemma~\ref{lem:u11-clock} \\
The timestamp is undriven on the states its runs reach, while the timestamp started in sync is a latch there and is rejected only by (R1) over every state & proved & Section~\ref{sec:latch}, Lemma~\ref{lem:u11-clock-sync} \\
Base-driven records agree up to the schedule & proved  & Section~\ref{sec:latch} \\
The event is an unforced choice & proved, the converse applied twice  & Section~\ref{sec:latch}, Corollary~\ref{cor:u11-event-free} \\
The schedule doesn't fix the probabilities & proved, with the branching latch refuted & Section~\ref{sec:latch}, Theorem~\ref{thm:u11-probabilistic} \\
One state space up to renaming & false: an extension may carry any hidden field & Section~\ref{sec:latch} \\
The latch on a Turing machine, a register machine, L & definition; proved (and the relation is an equivalence) & Section~\ref{sec:latch}, Definition~\ref{def:u11-weak-equiv}, Theorem~\ref{thm:u11-four-bases} \\
Weak base equivalence connects those bases & not claimed & Section~\ref{sec:latch} \\
A growing record is a family of latches & definition; proved, and one latch isn't enough & Section~\ref{sec:latch}, Definition~\ref{def:u11-growing}, Theorem~\ref{thm:u11-growing} \\
Thiele-complete means a universal base whose moves cost nothing and can't touch the record, every other move costing exactly one, a record that stays up, blank starts, and some run whose record ends up from one start and down from another; nothing in it says what the record is about & definition & Section~\ref{sec:complete}, Definition~\ref{def:thiele-complete} \\
Every Thiele-complete machine obeys the toll, raises its record only by a mark costing exactly one, counts its marks on its ledger, pays at least one on any run that raises the record, and runs every two-counter program & proved, for every machine and every interface meeting the definition & Section~\ref{sec:complete}, Theorem~\ref{thm:tc-has} \\
A clock is not Thiele-complete, nor is a machine whose record never rises, and each of the four clauses is needed & proved & Section~\ref{sec:complete}, Observation~\ref{obs:tc-clock}, Theorem~\ref{thm:tc-clauses} \\
Every Thiele-complete machine hides its record and its ledger from its window, with a gap of one unit, and one is reached & proved & Section~\ref{sec:complete}, Theorem~\ref{thm:tc-window} \\
Every universal base becomes Thiele-complete by adding marks, with no claims and no checker & proved & Section~\ref{sec:complete}, Theorem~\ref{thm:tc-lift} \\
Certified Thiele-complete is exactly Thiele-complete with clean starts, the earned chain and a claim that can fail; the inclusion is strict & proved & Section~\ref{sec:certified}, Theorem~\ref{thm:tc-decomposition} \\
In a partial order the cups fix the place; in a preorder they need not & proved & Section~\ref{sec:direction}, Theorem~\ref{thm:dir-lines} \\
The toll over any direction: every step that turns a cup over costs at least one exactly when every run that ends somewhere not at or below its start does, and exactly when every cup obeys the toll on its own; a run then costs at least its number of such steps & proved & Section~\ref{sec:direction}, Theorem~\ref{thm:dir-toll} \\
A chain of $n+1$ stages needs $n$ write-once Boolean cups & proved; ordinary rewritable binary storage is a different question & Section~\ref{sec:direction}, Theorem~\ref{thm:dir-bits} \\
On a finite machine, a step that turns a cup over merges or revokes; with a growing record it merges, and merge pricing gives the toll & proved & Section~\ref{sec:direction}, Theorem~\ref{thm:dir-merge} \\
Finiteness, growth and merge pricing each do work for the direction toll & proved (three machines) & Section~\ref{sec:direction}, Theorem~\ref{thm:dir-escapes} \\
A window can price rises exactly without determining the record value & proved, under the collision and decidability conditions stated there & Section~\ref{sec:shadow-exact}, Theorem~\ref{thm:dir-window} \\
With an injective shadow step, merged states lie in one shadow fibre & proved; the shadow-injectivity premise is needed & Section~\ref{sec:shadow-exact}, Theorem~\ref{thm:dir-fibre} \\
Halving pricing bounds weighted conditional entropy; an infinite landing pile admits no finite halving price & proved; average loss still depends on the distribution & Section~\ref{sec:shadow-exact}, Theorems~\ref{thm:dir-bits-loss} and~\ref{thm:dir-infinite} \\
\end{longtable}
\endgroup

\begingroup\small
\begin{longtable}{@{}>{\raggedright\arraybackslash}p{0.242\linewidth}>{\raggedright\arraybackslash}p{0.344\linewidth}>{\raggedright\arraybackslash}p{0.344\linewidth}@{}}
\toprule
\textbf{Claim} & \textbf{Kind} & \textbf{Home} \\
\midrule
\endhead
\bottomrule
\endlastfoot
Substrate; shortcut predicate & definition & Section~\ref{sec:diagonal}, Definition~\ref{def:u12-substrate}, Definition~\ref{def:u12-shortcut-predicate} \\
No inside decider with a representable flip & proved (premise: the substrate's recursion theorem) & Section~\ref{sec:diagonal}, Theorem~\ref{thm:u12-no-inside-decider} \\
The nat family: each candidate decider misses program two in the instance built from it & proved & Section~\ref{sec:diagonal}, Theorem~\ref{thm:u12-nat-family} \\
An outside decider gets the nat family right & proved  & Section~\ref{sec:diagonal} \\
Recursion, Rice and halting for L & definition; proved & Section~\ref{sec:diagonal}, Definition~\ref{def:u12-L}, Theorem~\ref{thm:u12-L} \\
Undecidable means unknowable; the wall separates Turing degrees & not claimed & Section~\ref{sec:diagonal} \\
Casper setting; one storage slot; a gas schedule & definition (transcribed; the gas schedule's charge premise stipulated) & Section~\ref{sec:models}, Definition~\ref{def:u13-casper}, Definition~\ref{def:u13-slot}, Definition~\ref{def:u13-gas-schedule} \\
Accountable safety & proved (premise: quorum intersection; the one parent the setting carries is never used) & Section~\ref{sec:models}, Theorem~\ref{thm:u13-accountable-safety} \\
The quorum premise can hold; on a one-validator chain one vote finalizes genesis and nobody is slashed & proved & Section~\ref{sec:models}, Theorem~\ref{thm:u13-casper-instances} \\
A persistent write to a fresh slot nets at least 22100, and to a warm one at least 20000 & proved, on the transcription & Section~\ref{sec:models}, Theorem~\ref{thm:u13-sstore} \\
Gas is exact exactly when commits pay one & proved (premise: the stipulated charge predicate) & Section~\ref{sec:models}, Theorem~\ref{thm:u13-gas-exactness} \\
A permanent write on a finite machine merges, and the merge lands on the record's yes side & proved & Section~\ref{sec:models}, Theorem~\ref{thm:u13-logical-payment} \\
Platform and quote; inclusion check; bounds fragment & definition & Section~\ref{sec:models}, Definition~\ref{def:u13-quote}, Definition~\ref{def:u13-inclusion}, Definition~\ref{def:u13-pcc} \\
The quote can't attest what it never measured & proved & Section~\ref{sec:models}, Theorem~\ref{thm:u13-quote} \\
\end{longtable}
\endgroup

\begingroup\small
\begin{longtable}{@{}>{\raggedright\arraybackslash}p{0.242\linewidth}>{\raggedright\arraybackslash}p{0.344\linewidth}>{\raggedright\arraybackslash}p{0.344\linewidth}@{}}
\toprule
\textbf{Claim} & \textbf{Kind} & \textbf{Home} \\
\midrule
\endhead
\bottomrule
\endlastfoot
The RFC's inclusion examples pass & proved (symbolic hash) & Section~\ref{sec:models}, Theorem~\ref{thm:u13-ct-examples} \\
The proof-carrying checker is the verification condition & proved & Section~\ref{sec:models}, Theorem~\ref{thm:u13-pcc} \\
A transcription is the real system & not claimed & Section~\ref{sec:models} \\
A toll is a better design than a wall & not claimed & Section~\ref{sec:models} \\
Ecosystem, proliferation, pointer; copied and relied on, independent carriers; observation game & definition & Section~\ref{sec:pointer}, Definition~\ref{def:u14-ecosystem}, Definition~\ref{def:u14-reliance}, Definition~\ref{def:u14-game} \\
Independent copies on every state, with two observers or more: the event is equivalent at any two states & proved & Section~\ref{sec:pointer}, Theorem~\ref{thm:u14-independent-constant} \\
In the specified replicated-header toy, gas used and finality are copied, and the chosen action maps rely only on finality & proved for those maps & Section~\ref{sec:pointer}, Theorem~\ref{thm:u14-ledger} \\
Certification is the event whose records are copied and relied on by parties who didn't produce them & thesis & Section~\ref{sec:pointer}, Thesis~\ref{conj:u14-pointer} \\
In the toy, ``certified'' is the unique pointer & proved & Section~\ref{sec:pointer}, Theorem~\ref{thm:u14-toy} \\
The survey; the swap shows the verdict follows the observer maps & proved & Section~\ref{sec:pointer}, Theorem~\ref{thm:u14-survey} \\
One blind observer stops proliferation & proved & Section~\ref{sec:pointer}, Theorem~\ref{thm:u14-blind} \\
Agreement doesn't buy permanence; durable views that report the event do & proved & Section~\ref{sec:pointer}, Theorem~\ref{thm:u14-toggle} \\
The formal objects prove the thesis & not claimed; ``independent'' is defined, ``deployed'' and ``faithfully model'' are defined nowhere & Section~\ref{sec:pointer} \\
The early engine derived physics & not claimed; it ran a simulation & Section~\ref{sec:origin} \\
The floor raised to a threshold & proved (premises: the witness starts at zero, grows by at most the cost, and meets the threshold) & Appendix~\ref{sec:assumptions}, Theorem~\ref{thm:app-quantitative} \\
\end{longtable}
\endgroup



\subsection{Claims about the machine}

Part III built one small machine, with the hosts that run it in an appendix, and the ledger keeps
them apart from the model. The machine is the small one of Section~\ref{sec:smallest}, two counters and six instruction kinds, where commitments are earned
and universality is proved, with the definition of certified Thiele-complete that every machine can be
held against, the universal program $U$ on a host made of the same instruction kinds, and the
priced program $U_P$ that runs every computably presented Thiele machine. Every line in the
next tables is about those machines, or about every machine meeting that definition, or about
the physics section that closes Part III and the game check in Appendix~\ref{sec:chsh}, or about the hosts, bases, models, parts and
tools the appendices build around them, and nothing else, under the premises printed in its section.

\begingroup\small
\begin{longtable}{@{}>{\raggedright\arraybackslash}p{0.242\linewidth}>{\raggedright\arraybackslash}p{0.344\linewidth}>{\raggedright\arraybackslash}p{0.344\linewidth}@{}}
\toprule
\textbf{Claim} & \textbf{Kind} & \textbf{Home} \\
\midrule
\endhead
\bottomrule
\endlastfoot
Checker and arithmetic meaning agree & proved on the small machine, under the stated premises & Section~\ref{sec:smallest}, Lemma~\ref{lem:u34-eval} \\
Two-counter simulation and halting correspondence & proved on the small machine, under the stated premises & Section~\ref{sec:smallest}, Theorem~\ref{thm:u34-simulation} (step), Theorem~\ref{thm:u34-halting} (halting) \\
Halting is undecidable by the two-counter reduction & proved on the small machine, under the stated premises & Section~\ref{sec:smallest}, Theorem~\ref{thm:u34-undecidable} \\
Exact ledger, permanent certification and toll & proved on the small machine, under the stated premises, directly from the rules & Section~\ref{sec:smallest}, Theorem~\ref{thm:u34-toll} \\
A Thiele machine & proved on the small machine, under the stated premises & Section~\ref{sec:smallest}, Theorem~\ref{thm:u34-instance} \\
Certification factors as a latch over the core base & proved on the small machine, under the stated premises & Section~\ref{sec:smallest}, Theorem~\ref{thm:u34-latch} \\
A passing commitment has a prior matching check at the same version & proved on the small machine, under the stated premises & Section~\ref{sec:smallest}, Theorem~\ref{thm:u34-commit-earned} \\
Certification has a check, commitment and certification in that order & proved on the small machine, under the stated premises & Section~\ref{sec:smallest}, Theorem~\ref{thm:u34-cert-earned} \\
Unearned commitment and uncommitted certification trap & proved on the small machine, under the stated premises & Section~\ref{sec:smallest}, Theorem~\ref{thm:u34-failures} \\
Current-version checked claims are true & proved on the small machine, under the stated premises & Section~\ref{sec:smallest}, Theorem~\ref{thm:u34-soundness} \\
No stored fact is forged on a clean-start trace & proved on the small machine, under the stated premises & Section~\ref{sec:smallest}, Theorem~\ref{thm:u34-no-forging} \\
Bounded table keeps facts and traps when full & proved on the small machine, under the stated premises & Section~\ref{sec:smallest}, Theorem~\ref{thm:u34-bounded} \\
Equal computational windows can hide different records & proved on the small machine, under the stated premises & Section~\ref{sec:smallest}, Theorem~\ref{thm:u34-separation} \\
The window cannot recover the ledger, flag or commitment permission & proved on the small machine, under the stated premises & Section~\ref{sec:smallest}, Theorem~\ref{thm:u34-no-oracle} \\
From a clean start, a run ending with the flag up adds at least three to the ledger & proved on the small machine, under the stated premises, directly from the rules & Section~\ref{sec:smallest}, Observation~\ref{thm:u34-min-price} \\
That minimum of three is reached by a run whose check could have failed, and when that check fails the program is refused forever & proved on the small machine, under the stated premises & Section~\ref{sec:smallest}, Theorem~\ref{thm:u34-tight} \\
The step refuses trivial checks & not claimed; ``is at least 0'' passes every value and certifies for three too & Section~\ref{sec:smallest} \\
From a clean start, provenance, checker soundness, no forging and the floor of three hold for any property language with an exact equality test and a checker proved to say yes exactly when the claim holds, a sorted list read out of a counter included & proved on the generic machine  & Section~\ref{sec:smallest} \\
With the flag up from a clean start, the channel names a fact an earlier passing COMMIT committed, and a later commit can replace it & proved on the small machine  & Section~\ref{sec:smallest} \\
The guarantees hold against whoever sets the start, a reset, or a second writer that doesn't bump versions & not claimed; they are about programs from a clean start, one step at a time & Section~\ref{sec:smallest}, Section~\ref{sec:scope} \\
Certified Thiele-complete means a universal base that costs nothing and can't touch the record, a record that stays up once it's up, an earned record with a checker proved to mean its claims, exact prices of one per check, commit and certificate, and a claim whose check can fail & definition & Section~\ref{sec:smallest}, Definition~\ref{def:u34-thiele-complete} \\
A machine certified Thiele-complete over a claim language fixed first is certified Thiele-complete, and from a clean state its raised record stands on a claim whose meaning held when it was checked and when it was committed & proved & Section~\ref{sec:smallest}, Theorem~\ref{thm:u34-tc-over} \\
Every certified Thiele-complete machine obeys the toll, runs every two-counter program, counts its record moves on its ledger, pays at least three on any run from a clean state that ends with the record up, has a check that fails and a run that certifies, and from a clean state raises its record only by a CERTIFY & proved, for every machine and every interface meeting the definition & Section~\ref{sec:smallest}, Theorem~\ref{thm:u34-tc-consequences} \\
A clock (every move charged, any reading or any latch as the record) obeys the toll with a two-counter machine running inside it and is not certified Thiele-complete; nor is any machine whose every move costs something & proved, directly from the rules & Section~\ref{sec:smallest}, Observation~\ref{thm:u34-clock} \\
Each clause does work: a paid latch meets the universal base and exact prices clauses and isn't certified Thiele-complete, and a silent machine meets three clauses and fails non-vacuity & proved & Section~\ref{sec:smallest}, Theorem~\ref{thm:u34-tc-clauses} \\
The four clauses of certified Thiele-complete are independent: for each, a Thiele machine with a two-counter machine inside meets the other three, fails it, and is certified Thiele-complete through no interface & proved & Section~\ref{sec:smallest}, Theorem~\ref{thm:tc-independent} \\
Certified Thiele-complete at any unit: with record moves at one fixed price $c \ge 1$ the class is the unit class with its prices multiplied by $c$, and the clock fails at every unit & proved & Section~\ref{sec:smallest}, Theorem~\ref{thm:u34-tc-unit} \\
The small machine, the same machine over any property language with an exact equality test, a checker proved to say yes exactly when the property holds and some property true of one number and false of another, and the sorted-list machine are certified Thiele-complete & proved on the small machine and on the generic machine & Section~\ref{sec:smallest}, Theorem~\ref{thm:u34-tc-instances} \\
Every certified Thiele-complete machine hides its record and its ledger from its base window & proved, for every machine and every interface meeting the definition, directly from the rules & Section~\ref{sec:smallest}, Observation~\ref{thm:u34-tc-window} \\
A counter beside a Turing machine changes nothing & proved & Section~\ref{sec:machine-vs-program}, Theorem~\ref{thm:u19-tm-lift} \\
A Turing machine can't price a reading & not claimed & Section~\ref{sec:machine-vs-program} \\
The calibration is a unit conversion & definition; proved & Section~\ref{sec:landauer}, Definition~\ref{def:u24-scale}, Theorem~\ref{thm:u24-calibration} \\
Structural dissipation pricing (a whole-unit floor on collapsing instructions, separate from Landauer's principle); dissipation at most cost & assumed (premise) & Section~\ref{sec:landauer}, Premise~\ref{prem:u24-landauer-pair} \\
The premises fit when every counted step costs at least one, and fit a one-bit reset & proved & Section~\ref{sec:landauer}, Theorem~\ref{thm:u24-consistency} \\
A collapse needs the merge clause: INC $A$ raises ``$A \ge 1$'' one way, one-to-one and costing nothing; on a finite machine a one-way step always merges & proved  & Section~\ref{sec:landauer}, Premise~\ref{prem:u24-landauer-pair} \\
The small machine's whole instruction set fails the premise pair & proved; its unpriced steps merge  & Section~\ref{sec:landauer} \\
Keeping history makes every step injective; the history takes fresh states, one more entry every step & proved & Section~\ref{sec:landauer}, Theorem~\ref{thm:u24-history} \\
The heat is fixed by the chosen gap & definition; proved & Section~\ref{sec:landauer}, Definition~\ref{def:u24-calorimeter}, Theorem~\ref{thm:u24-calorimeter} \\
The canonical reset's rates aren't in detailed balance at any gap or temperature; with rates that are, a fixed-gap relaxation moves toward $1/(1 + e^{\Delta/k_B T})$, which is $1/5$ at the gap $2 k_B T \ln 2$ & proved  & Section~\ref{sec:landauer} \\
A driven reset with rates in detailed balance, its gap raised from 0 to $D$ in $N$ equal steps, pays between $kT \ln 2 - kT e^{-D/kT}$ and $kT \ln 2 + D/(2N)$ in work & definition; proved & Section~\ref{sec:landauer}, Definition~\ref{def:u24-driven}, Theorem~\ref{thm:u24-driven} \\
The counter is physical dissipation, in joules & not claimed & Section~\ref{sec:landauer} \\
Permanence predicts more than Landauer & not claimed & Section~\ref{sec:landauer} \\
The theorems prove physical reality & not claimed & Section~\ref{sec:landauer} \\
No fixed host program can keep the guest's facts on an exact copy of a guest counter & proved, for any host property language and any translation into a finite list of host properties & Appendix~\ref{sec:universal-host}, Theorem~\ref{thm:u34-nocopy} \\
A universal Thiele machine for a class: one host and one program, any guest loaded into a start state, halting matched, the guest's states decoded in order, the host's flag the guest's reading latched, the host's ledger the guest's plus a surcharge, never negative and bounded for the class, and the flag raised only by the host's own moves under its own toll & definition (surcharge 0 for $U$; for $U_P$ at most 2 when the guest's reading starts at no and exactly 3 when it starts at yes, so 3 over the class) & Appendix~\ref{sec:universal-host}, Definition~\ref{def:universal-thiele} \\
One fixed program $U$, in the small machine's own instruction kinds and with no guest step built in, runs every small-machine program from every clean start $\mathrm{start}(x, y)$, step for step, with halting and answers matched & proved on the host machine & Appendix~\ref{sec:universal-host}, Theorem~\ref{thm:u34-U-sim}, Theorem~\ref{thm:u34-U-halting} \\
$U$'s flag rises at some step if and only if the guest's does, and its ledger equals the guest's at every matching point & proved on the host machine & Appendix~\ref{sec:universal-host}, Theorem~\ref{thm:u34-U-flag} \\
$U$'s raised flag stands on $U$'s own check, commit and certificate on a slot, and on the guest's own earned chain for the same claim & proved on the host machine & Appendix~\ref{sec:universal-host}, Theorem~\ref{thm:u34-U-earned} \\
$U$'s host is certified Thiele-complete, and whether $U$ halts on a loaded guest is undecidable & proved on the host machine; undecidable in the reduction sense of Theorem~\ref{thm:u34-undecidable} & Appendix~\ref{sec:universal-host}, Theorem~\ref{thm:u34-U-complete} \\
One fixed program $U_P$, $U$ with one paid move added, runs every computably presented Thiele machine: halting matched, the state decoded at every matching point, the host's flag up exactly when the guest's reading has been yes, and the host's ledger the guest's plus a surcharge & proved on the priced host & Appendix~\ref{sec:presented-universal}, Theorem~\ref{thm:u34-presented} \\
$U_P$'s raised flag stands on its own check, commit and certificate on a slot whose claim is the guest's own reading of a state where it said yes & proved on the priced host & Appendix~\ref{sec:presented-universal}, Theorem~\ref{thm:u34-presented} \\
The surcharge is at most two when the reading starts at no, and zero when the first raising move costs at least three & proved & Appendix~\ref{sec:presented-universal}, Theorem~\ref{thm:u34-presented-surcharge} \\
No program on the host machine, run from a clean start, exactly matches a guest certified with its whole ledger below three & proved, for every program from a clean start, directly from the rules & Appendix~\ref{sec:presented-universal}, Theorem~\ref{thm:u34-presented-surcharge} \\
The priced host is certified Thiele-complete, with the paid move read as a check that never passes & proved & Appendix~\ref{sec:presented-universal} \\
A demonstration machine (one move of price three, reading ``at least three'') runs on $U_P$, with the host's ledger exactly $3n$ at the host time matching move $n$ & proved; never run & Appendix~\ref{sec:presented-universal} \\
A machine whose reading no program decides, or whose next move no program computes, runs on $U_P$ & not claimed; outside the class & Appendix~\ref{sec:presented-universal} \\
A guest's reading that goes back to no is carried as no & not claimed; the host's flag is the latch of the reading & Appendix~\ref{sec:presented-universal} \\
\end{longtable}
\endgroup

\begingroup\small
\begin{longtable}{@{}>{\raggedright\arraybackslash}p{0.242\linewidth}>{\raggedright\arraybackslash}p{0.344\linewidth}>{\raggedright\arraybackslash}p{0.344\linewidth}@{}}
\toprule
\textbf{Claim} & \textbf{Kind} & \textbf{Home} \\
\midrule
\endhead
\bottomrule
\endlastfoot
The simpler host runs any small-machine program as a guest, and halts with any two-counter program & proved on the small machine, by construction: its guest step runs the guest's own step as a primitive; the guest program is held as a list & Appendix~\ref{sec:simpler-host}, Theorem~\ref{thm:u34-universal} \\
From a loaded host, or any state where the two agree, the simpler host's mirror always equals the guest's flag & proved on the small machine, directly from the rules & Appendix~\ref{sec:simpler-host}, Observation~\ref{thm:u34-mirror} \\
From a host that agrees, the simpler host's step collects the guest's toll; on every trace its ledger covers the guest's & proved on the small machine & Appendix~\ref{sec:simpler-host}, Theorem~\ref{thm:u34-host-toll} \\
The simpler host's mirror rises only in a guest step that passed CERTIFY & proved on the small machine, for every sequence of host moves & Appendix~\ref{sec:simpler-host}, Theorem~\ref{thm:u34-no-free-host} \\
The simpler host's own instructions never touch the guest & proved on the small machine & Appendix~\ref{sec:simpler-host}, Theorem~\ref{thm:u34-no-free-host} \\
From a loaded host, the simpler host's raised mirror means the guest's own run checked, committed and certified the same claim & proved on the small machine  & Appendix~\ref{sec:simpler-host} \\
The simpler host is a Thiele machine & proved on the small machine & Appendix~\ref{sec:simpler-host}, Theorem~\ref{thm:u34-host-instance} \\
The simpler host's ledger equals the guest's & not claimed; a guest step with budget $b$ costs $b$ whatever the guest instruction cost, so its ledger only covers the guest's; $U$'s is exact & Appendix~\ref{sec:simpler-host}, Theorem~\ref{thm:u34-host-toll} \\
Rice's theorem on the host machine, for properties that depend on everything a program ends with (registers, table, channel, trap, ledger, flag) & proved; undecidable in the reduction sense of Theorem~\ref{thm:u34-undecidable} & Appendix~\ref{sec:small-diagonal}, Theorem~\ref{thm:host-rice} \\
Rice's theorem for two-counter programs, from the clean start & proved; the fact table isn't compared & Appendix~\ref{sec:small-diagonal}, Theorem~\ref{thm:guest-rice} \\
s-m-n, and a universal host program for host programs & proved, with no premise on register 2 & Appendix~\ref{sec:small-diagonal}, Theorem~\ref{thm:host-smn} \\
Kleene's recursion theorem on the host machine & proved; the fixed point is of the function & Appendix~\ref{sec:small-diagonal}, Theorem~\ref{thm:host-kleene} \\
Kleene's recursion theorem up to the record on the host machine: stopping, register 0, trap latch, ledger exactly, flag, number of facts, whether the channel is empty & proved; the host program simulates the program $T$ that carries out the transformation, marks and all, and pays for the whole record in one go once it has computed & Appendix~\ref{sec:small-diagonal}, Theorem~\ref{thm:host-kleene-record} \\
A fixed point of full behaviour (every register, what the facts name) on the host machine & refuted for registers, and for the fact table, the channel and the versions  & Appendix~\ref{sec:small-diagonal}, Theorem~\ref{thm:host-no-exact} \\
No inside decider with a representable flip, and no host decider, for functional properties of host programs & proved (the recursion premise is a theorem there) & Appendix~\ref{sec:small-diagonal}, Theorem~\ref{thm:host-wall} \\
Rice's theorem for two-counter programs on every input (stop, counters, trap latch, ledger, flag, and what each fact and the channel are about) & proved & Appendix~\ref{sec:small-diagonal}, Theorem~\ref{thm:tc-rice} \\
Kleene's recursion theorem for two-counter programs, inputs written as powers of two, for the function & proved; refuted once the recorded facts count & Appendix~\ref{sec:small-diagonal}, Theorem~\ref{thm:tc-kleene} \\
Kleene's recursion theorem for two-counter programs on plain inputs & refuted: a carried-out transformation has no fixed point & Appendix~\ref{sec:small-diagonal}, Theorem~\ref{thm:no-plain-recursion} \\
Every universal base lifts with an exact checker, a nontrivial claim and room for a fact; finite simulations lift after grouping & proved under the stated interface conditions & Appendix~\ref{sec:every-model}, Theorems~\ref{thm:lift} and~\ref{thm:lift-macro} \\
Finite branching, state-independent moves and an indiscriminate record order give the stated obstructions to lifting & proved & Appendix~\ref{sec:every-model}, Theorem~\ref{thm:lift-fails} \\
Every box with a universal base is a slice of its lift with the window claims and room for a fact: forgetting the panel runs as the box runs while the trap is down, reaches exactly what the box reaches, isn't moved by CHECK, COMMIT or CERTIFY, and gives neither the flag nor the ledger, the chain's ledger exactly 3 above the empty run's; with no room in the table no interface certifies the lift; a Thiele-complete machine's computing moves are its moves that cost nothing, through every interface, and its window is a slice onto two counters; a machine run as one rule has no universal base, no base by simulation and no universal-base slice, and is a Thiele machine only as a clock or with a reading that never rises & proved; that the slice runs as the box runs and that the record moves leave it alone, directly from the rules, and the converse's universal base and step-for-step clauses directly from the interface & Appendix~\ref{sec:every-model}, Definition~\ref{def:slice}, Theorem~\ref{thm:slice} \\
With the program in the state, a reading is kept by every operation exactly when it never comes down from any state; a cell some program clears is a simulated bit & proved, directly from the rules & Appendix~\ref{sec:every-model}, Observation~\ref{obs:program-in-state} \\
The seven models compute the same numerical relations under their encodings; one counter has decidable halting & proved & Appendix~\ref{sec:every-model}, Theorems~\ref{thm:lift-models} and~\ref{thm:one-counter} \\
Products and sequential hand-offs preserve the stated structure; shared writers and unsuitable record readings can break it & proved under the stated interfaces and loading conditions & Appendix~\ref{sec:parts}, Theorems~\ref{thm:parts-product}--\ref{thm:parts-breaks} \\
A following host run pays for the guest's record rises; presented nesting surcharges add & proved; every host run is presented on its states with finitely many live registers, so towers go to any depth & Appendix~\ref{sec:parts}, Theorem~\ref{thm:parts-nesting} \\
The source compiler preserves halting and numerical output through the host, two-counter guest and universal program & proved; final two stages not executed & Appendix~\ref{sec:running}, Theorem~\ref{thm:running-compiler} \\
Executable implementations agree on the tested whole-state transitions & tested, not a guarantee about untested inputs or physical hardware & Appendix~\ref{sec:running} \\
\end{longtable}
\endgroup

\begingroup\small
\begin{longtable}{@{}>{\raggedright\arraybackslash}p{0.242\linewidth}>{\raggedright\arraybackslash}p{0.344\linewidth}>{\raggedright\arraybackslash}p{0.344\linewidth}@{}}
\toprule
\textbf{Claim} & \textbf{Kind} & \textbf{Home} \\
\midrule
\endhead
\bottomrule
\endlastfoot
The integer check makes the pinned matrix PSD & definition; proved & Appendix~\ref{sec:chsh}, Definition~\ref{def:u23-pinned}, Theorem~\ref{thm:u23-notrap-psd} \\
PSD forces $S^2 \le 8$ & proved (algebra only) & Appendix~\ref{sec:chsh}, Theorem~\ref{thm:u23-psd-tsirelson} \\
A passed check means $S^2 < 8$, and no smaller bound holds & proved: $S^2 < 8$, with passing tallies as close to 8 as desired & Appendix~\ref{sec:chsh}, Corollary~\ref{cor:u23-notrap-bound} \\
A PSD pinned matrix has a PSD completion & proved & Appendix~\ref{sec:chsh}, Remark~\ref{rem:u23-sound-not-complete} \\
Every deterministic plan fails the check & proved; no such plan's pinned matrix is PSD  & Appendix~\ref{sec:chsh}, Remark~\ref{rem:u23-sound-not-complete} \\
The level-1 test is the size-five case & definition; proved & Appendix~\ref{sec:chsh}, Definition~\ref{def:u23-genrealizable}, Corollary~\ref{cor:u23-general} \\
The 1+AB matrix ties its four-body cells, the constraint $\langle \{A_0, A_1\} \otimes \{B_0, B_1\} \rangle = 0$; the level-1 pinnings don't imply it, and a classical strategy breaks it & proved  & Appendix~\ref{sec:chsh}, Appendix~\ref{sec:chsh-derivation} \\
A tally above two, with every pair asked, isn't one plan replayed & proved & Appendix~\ref{sec:chsh-derivation}, Theorem~\ref{thm:chsh-tally} \\
A tally above two proves nonlocality; a pass proves the data quantum & not claimed & Appendix~\ref{sec:chsh} \\
A finite local model meets the bound & proved; Fine's theorem in correlation form is proved as an iff, and the converse of the single-bound statement is false  & Appendix~\ref{sec:chsh-derivation}, Theorem~\ref{thm:chsh-fine}, Section~\ref{sub:finite-local-construction} \\
The elliptope gate is sound; the PR box fails it & proved & Appendix~\ref{sec:chsh-derivation}, Theorem~\ref{thm:chsh-gate-sound} \\
The arc-sine boundary: a table has a PSD completion, and so is quantum, exactly when its arcsines pass the four inequalities {\footnotesize(if and only if)} & proved & Appendix~\ref{sec:chsh-derivation} \\
Every PSD completion comes from a quantum strategy, and the completion set is exactly the set of quantum correlator tables in finite dimension, real or complex & proved & Appendix~\ref{sec:chsh-derivation} \\
A two-qubit strategy reaches $2\sqrt{2}$; every strategy in finite dimension, with real or complex amplitudes, has $S^2 \le 8$ and a PSD level-1 moment matrix; every deterministic plan is such a strategy & proved & Appendix~\ref{sec:chsh-derivation} \\
The integer route to the bound is new mathematics & not claimed; the route is NPA's & Appendix~\ref{sec:chsh} \\
The integer check passes exactly when every pair was sampled and the pinned matrix is PSD & proved & Appendix~\ref{sec:chsh}, Theorem~\ref{thm:chsh-exact} \\
Honest play at Tsirelson's point fails the exact check about 85 percent of the time & simulated, with a first-order limit of $(2 + \sqrt 2)/4$; no theorem & Appendix~\ref{sec:chsh} \\
The check with a tolerance $\varepsilon$ is the exact check at $\varepsilon = 0$; a pass means operator norm at most $1 + \varepsilon$ and $S^2 \le 8(1 + \varepsilon)^2$ & proved, including the earned CHECK  & Appendix~\ref{sec:chsh} \\
A flag raised on the host machine from a clean start stands on a passing CHECK of the packed counts, a COMMIT and a CERTIFY; the chain run from a fresh start certifies for 3 exactly when the check passes & proved & Appendix~\ref{sec:chsh}, Theorem~\ref{thm:chsh-earned} \\
Five tallies: 12/5 and 14/5 certify; 16/5, the PR box and the all-ones plan are refused forever & proved & Appendix~\ref{sec:chsh}, Proposition~\ref{prop:chsh-demos} \\
The tally was played, or was honest, or puts a bound on nature & not claimed; it is loaded into a counter at no charge on the ledger & Appendix~\ref{sec:chsh} \\
\end{longtable}
\endgroup

\subsection{What I propose}

What ties all of those results together is a proposal, and it's mine: put the relations a
computation carries, and what keeping them takes, into the state and the step. The items below
spell that out, along with one narrower bet inside it, the pointer thesis about
certification. Every one of them reaches past the theorems, and none of them adds to the
physics, which stays in the premises the book names and the results it cites.

\begin{itemize}
\item That the axis belongs in the step (Section~\ref{sec:scope}, Section~\ref{sec:conclusion}).
\item That the usual account of a computation is a shadow of the step, for machines whose control never reads their record (Section~\ref{sec:shadow-exact}). For a computer with a move for each two-counter instruction, under the record layer of Appendix~\ref{sec:every-model}, part of that is proved: the computer is a slice of a certified Thiele-complete machine, what's left with the layer forgotten, and it can't see the record or the ledger (Theorem~\ref{thm:slice}). That layer's control reads its trap, so the slice follows the computer only while the trap is down. The proposal is the rest: that this is how to see computation in general. The small machine's control does read its own record layer: a commitment that finds no fact traps. So for it the claim is about its ordinary computer, not the whole machine.
\item Within the record construction, that certification is the event selected by the reliance criterion of the pointer thesis (Section~\ref{sec:pointer}).
\item That this is a structural account of computation, with a cost model as one part of it (Section~\ref{sec:scope}).
\item That the toll is the right thing to price. The alternatives are a wall or a penalty for contradiction (Section~\ref{sec:models}, Section~\ref{sec:conclusion}).
\end{itemize}

And there are three things this book doesn't settle, each with what would settle it. Whether Landauer's critics are
right, which would leave every theorem standing and one premise unowned by the world: that takes
measurements of real devices (Section~\ref{sec:a2-not-arbitrary}). Whether certification really is the
pointer event: that takes a way to tie the definitions to the systems the thesis is about, and Section~\ref{sec:pointer}
sets out the three comparisons needed. And whether there's a smaller start than a partition: that takes a ruler
someone can state and a smaller structure measured by it (Section~\ref{sec:structure}).

That's the whole ledger. If you find a claim in this book that isn't on it, or a line that
says more than its home does, that's my mistake, and I'd want to hear about it.

%===========================================================================


\section{The CHSH derivation, in full}

Appendix~\ref{sec:chsh} tells the story; this appendix does the algebra. Each item gives
the statement with its hypotheses and the steps you need to check it. Nothing here says
that a matrix is a physical state. Where quantum strategies come in, I say exactly how they
stand to the matrices.

\subsection{Setting and notation}


Start with four piles of tally sheets, one for each pair of settings Alice
and Bob can be given. On each round their answers either agree or disagree.
Give agreement a score of $+1$ and disagreement a score of $-1$. In each pile,
add the scores and divide by the number of rounds. That pile's average is its
\emph{correlator}. It says how strongly the two answers agree for that pair of
settings.

For example, three agreements and one disagreement give a total of two over
four rounds, hence an average of one half. The four averages make a little
square table: Alice's setting chooses the row, Bob's chooses the column.
The score I will study adds three cells of that table and subtracts the
bottom-right one. Here are the names for those quantities.

\begin{definition}[Settings, outcomes, correlators]
Alice receives a setting $x \in \{0,1\}$ and Bob a setting $y \in \{0,1\}$. Each answers a
bit, recoded as an outcome $a, b \in \{+1, -1\}$ (bit 0 becomes $+1$, bit 1 becomes $-1$).
The correlator $E_{xy}$ is the average of the product $ab$ over the rounds with settings
$(x, y)$. The CHSH quantity is
\[
S = E_{00} + E_{01} + E_{10} - E_{11}.
\]
I collect the four correlators into the $2 \times 2$ matrix
\[
E = \begin{pmatrix} E_{00} & E_{01} \\ E_{10} & E_{11} \end{pmatrix},
\]
whose rows are indexed by Alice's setting and whose columns by Bob's. The rows of $E$ are
$(E_{00}, E_{01})$ and $(E_{10}, E_{11})$; the columns are $(E_{00}, E_{10})$ and
$(E_{01}, E_{11})$.
\end{definition}

\begin{definition}[Counters]
There are eight counts: for each setting pair, $\mathrm{same}_{xy}$ counts rounds
whose two answer bits matched and $\mathrm{diff}_{xy}$ counts rounds whose bits differed. A
round adds one to exactly one of the eight. Write
\[
N_{xy} = \mathrm{same}_{xy} + \mathrm{diff}_{xy}, \qquad
D_{xy} = \mathrm{same}_{xy} - \mathrm{diff}_{xy}, \qquad
E_{xy} = \frac{D_{xy}}{N_{xy}}.
\]
\end{definition}

A matched round contributes $ab = +1$ and a differing round $ab = -1$, so $D_{xy}$ is the sum
of the products and $D_{xy}/N_{xy}$ is their average. Since $|D_{xy}| \le N_{xy}$, every
correlator from counters lies in $[-1,1]$. When a bucket is empty, its correlator is defined
as 0. The score is $E_{00}+E_{01}+E_{10}-E_{11}$.

\subsection{The classical bound}


Write a fixed plan on a card: Alice's answer to setting zero, her answer to
setting one, and Bob's two answers. Each answer is a sign, plus or minus.
Before multiplying anything, look at Bob's pair. Either the signs agree or
they differ. If they agree, adding them gives plus or minus two and subtracting
them gives zero. If they differ, the sum is zero and the difference is plus
or minus two. One of those two quantities vanishes; the other cannot exceed
two in size. The next calculation uses exactly that split.

\begin{lemma}[Deterministic strategies]\label{lem:chsh-deterministic}
Let $a_0, a_1, b_0, b_1 \in \{+1, -1\}$. Then $a_0 b_0 + a_0 b_1 + a_1 b_0 - a_1 b_1$ is $2$ or $-2$.
\end{lemma}

\begin{proof}
Factor: $a_0 b_0 + a_0 b_1 + a_1 b_0 - a_1 b_1 = a_0(b_0 + b_1) + a_1(b_0 - b_1)$. If
$b_0 = b_1$ then $b_0 + b_1 = \pm 2$ and $b_0 - b_1 = 0$; if $b_0 \ne b_1$ the other way
round. In both cases exactly one of the two brackets is $\pm 2$ and the other is $0$, so the
expression is $\pm 2$ times a sign, which is $+2$ or $-2$. There are sixteen sign choices,
and every one of them gives exactly $\pm 2$. None gives $0$ or $\pm 4$.
\end{proof}

You can also just list the sixteen sign choices and add each one up. Nothing in the argument needs real
numbers: the signs are whole numbers, so the bound holds the same way whether you read the
correlators as rationals (what a counter gives) or as reals.

A fixed card is only the first case. You can shuffle among several cards,
perhaps using some more often than others. Every card still tells each player
an answer using that player's setting alone. The shuffle changes how often
each card contributes; it does not change the score that a particular card
would produce. Since each card's score lies between minus two and two, their
average must stay between those same ends. The following definition writes
out that weighted shuffle.

\begin{definition}[Finite factorizable model]
Four real correlators $E_{00}, E_{01}, E_{10}, E_{11}$ are factorizable when there are a
natural number $\lambda_{\max}$, weights $p(\lambda)$ and response signs $A(x, \lambda)$,
$B(y, \lambda)$ for $\lambda = 0, \ldots, \lambda_{\max}$ such that:
$0 \le p(\lambda) \le 1$;
$\sum_{\lambda = 0}^{\lambda_{\max}} p(\lambda) = 1$; every $A(x, \lambda)$ and every
$B(y, \lambda)$ is $-1$ or $+1$; and for all $x, y$,
$E_{xy} = \sum_{\lambda=0}^{\lambda_{\max}} p(\lambda) A(x, \lambda) B(y, \lambda)$.
In words: a hidden card $\lambda$ is drawn with probability $p(\lambda)$, and each player reads
an answer off the card using only their own setting. Only the correlators are modelled.
\end{definition}

\begin{theorem}[The classical bound for a finite local model]
If $E$ is factorizable, then $-2 \le E_{00} + E_{01} + E_{10} - E_{11} \le 2$.
\end{theorem}

\begin{proof}
Substitute the factorization: $S = \sum_\lambda p(\lambda) S_\lambda$, where $S_\lambda$ is
the deterministic expression of Lemma~\ref{lem:chsh-deterministic} at the signs
$A(\cdot, \lambda), B(\cdot, \lambda)$. Each $S_\lambda \in [-2, 2]$, the weights are
nonnegative and sum to 1, so $-2 = \sum_\lambda p(\lambda)(-2) \le S \le
\sum_\lambda p(\lambda) \cdot 2 = 2$.
\end{proof}

\begin{remark}[Scope of the classical bound]\label{rem:chsh-fine-scope}
Fine's 1982 theorem is an equivalence between several descriptions of local models and the
Bell inequalities. The preceding theorem proves one direction: a finite factorizable model
obeys the bound. The converse of this single inequality is false: $(1, 1, -1, 1)$ has $S = 0$,
but it breaks another of the sign variants ($E_{00} + E_{01} - E_{10} + E_{11} = 4$), so no
local model gives it. The correlation-form equivalence needs each correlator between $-1$ and $1$
and all eight inequalities; with those, Section~\ref{sub:finite-local-construction} builds the
local model. The preceding theorem is stated for finitely many hidden cards ($\lambda\le\lambda_{\max}$), and that loses nothing. Whatever the hidden variable is and however it is drawn, the correlators see it only through its four answer signs, so lumping together the cards that carry the same four signs gives at most sixteen cards and the same four correlators. All the lumping needs is a probability for each of the sixteen sign patterns. \end{remark}

The finite-tally statement concerns one fixed strategy replayed on every recorded round.

\begin{definition}[A tally consistent with one local plan]
For bits $a_0, a_1, b_0, b_1 \in \{0, 1\}$ and a counter record $w$, $w$ is locally consistent with
$(a_0, a_1, b_0, b_1)$ when, for each setting pair $(x, y)$: if $a_x = b_y$ then
$\mathrm{diff}_{xy} = 0$, otherwise $\mathrm{same}_{xy} = 0$; and every bucket is sampled,
$N_{00}, N_{01}, N_{10}, N_{11} > 0$.
\end{definition}

\begin{theorem}[A tally above 2 is not one replayed strategy]
Let $w$ be a counter record whose rational score $S(w)$ is greater than 2.
Then for all bits $a_0, a_1, b_0, b_1$, $w$ is not locally consistent with
$(a_0, a_1, b_0, b_1)$. Every locally consistent record has $S(w) = 2$ or $S(w) = -2$; the first statement is that fact turned around.
\end{theorem}

\begin{proof}
Under local consistency each sampled bucket is pure: if $a_x = b_y$ every round in the
bucket matched, so $E_{xy} = N_{xy}/N_{xy} = +1$; otherwise every round differed and
$E_{xy} = -1$. In signs, $E_{xy} = \alpha_x \beta_y$ with $\alpha_x = (-1)^{a_x}$ and
$\beta_y = (-1)^{b_y}$. The score is then the deterministic expression of
Lemma~\ref{lem:chsh-deterministic}, which is 2 or $-2$. A score above 2 contradicts this.
\end{proof}

The theorem rules out one explanation, a single four-bit plan used on every round. It says
nothing about a local model that redraws its hidden state every round; Appendix~\ref{sec:chsh} gives a four-round coin-flip tally with score 4. The clause that every bucket is sampled does work here. An empty bucket reads as 0, and the plan that answers 0 to everything, never asked the pair 11, would score $1 + 1 + 1 - 0 = 3$.

\subsection{Tsirelson's bound from algebra}


Now put the cards aside and look just at four numbers in the square. I want
to bound their three-plus-one-minus score from their squared sizes. The useful
trick is to write the proposed allowance minus the squared score as a pile
of squares. A square cannot be negative. If the difference really is that
pile, the score cannot outrun the allowance.

The identity below is the receipt for that trick. The letters $a,b,c,d$ are
the four entries read across the first row, then across the second. It holds
for any four real numbers. After proving it, I will put in the two bounds
on row length that this part of the checker supplies.

\begin{lemma}[The six-square identity]
For all real $a, b, c, d$,
\[
\begin{aligned}
&4(a^2 + b^2 + c^2 + d^2) - (a + b + c - d)^2 \\
&\quad = (a-b)^2 + (a-c)^2 + (a+d)^2 + (b-c)^2 + (b+d)^2 + (c+d)^2.
\end{aligned}
\]

\end{lemma}

\begin{proof}
Expand the square on the left:
$(a+b+c-d)^2 = a^2 + b^2 + c^2 + d^2 + 2ab + 2ac - 2ad + 2bc - 2bd - 2cd$. So the left side is
\[
3(a^2 + b^2 + c^2 + d^2) - 2ab - 2ac + 2ad - 2bc + 2bd + 2cd.
\]
On the right, each variable appears in exactly three of the six brackets ($a$ in
$a-b$, $a-c$, $a+d$; $b$ in $a-b$, $b-c$, $b+d$; $c$ in $a-c$, $b-c$, $c+d$; $d$ in $a+d$,
$b+d$, $c+d$), so the squared terms add to $3(a^2+b^2+c^2+d^2)$. The cross terms are
$-2ab$ from $(a-b)^2$, $-2ac$ from $(a-c)^2$, $+2ad$ from $(a+d)^2$, $-2bc$ from $(b-c)^2$,
$+2bd$ from $(b+d)^2$ and $+2cd$ from $(c+d)^2$. Both sides agree term by term.
\end{proof}

Since the right side is a sum of squares, $S^2 \le 4 \sum_{xy} E_{xy}^2$ for any four reals.

\begin{theorem}[Row-bound form of Tsirelson's bound]
For all real $E_{00}, E_{01}, E_{10}, E_{11}$ with $E_{00}^2 + E_{01}^2 \le 1$ and
$E_{10}^2 + E_{11}^2 \le 1$, $S^2 \le 8$. The absolute form $|S| \le \sqrt 8 = 2\sqrt 2$ follows.
\end{theorem}

\begin{proof}
Adding the two row bounds gives $\sum_{xy} E_{xy}^2 \le 2$. By Lemma~\ref{lem:chsh-sos},
$S^2 \le 4 \sum_{xy} E_{xy}^2 \le 8$.
\end{proof}

\begin{theorem}[Tightness]
The point $E_{00} = E_{01} = E_{10} = 1/\sqrt 2$, $E_{11} = -1/\sqrt 2$ satisfies both row
bounds and has $S = \sqrt 8$.
\end{theorem}

\begin{proof}
Each row has squared length $1/2 + 1/2 = 1$, and $S = 4/\sqrt 2 = 2\sqrt 2 = \sqrt 8$.
\end{proof}

\begin{remark}[Where equality can happen]
Equality $S^2 = 8$ under the row bounds forces
equality in both steps of the proof above: the six squares must all vanish, which gives
$a = b = c = -d$, and the rows must be tight, which gives $2a^2 = 1$. So the point of
Theorem~\ref{thm:chsh-tight} and its negative are the only places the bound is reached. A tally's correlators are
ratios of whole numbers and $2a^2 = 1$ has no such solution, so no tally reaches the bound; the integer
check's passing tallies stay strictly below it and come as near as you like. (Take rational
$a, b$ with $a^2 + b^2 = 1$, from a Pythagorean triple, and the correlators $(a, b, b, -a)$: they pass
the check of Section~\ref{sub:chsh-integer} with no slack, and $S = 2(a + b)$, which gets as close to
$2\sqrt 2$ as you like as $a/b$ approaches 1.)
\end{remark}

\begin{lemma}[A rational constant above the bound]\label{lem:chsh-rational}
$\sqrt 8 < 5657/2000$.
\end{lemma}

\begin{proof}
Both numbers are positive, so it suffices to compare squares. $5657^2 = 32{,}001{,}649$ and
$8 \cdot 2000^2 = 32{,}000{,}000$, so $(5657/2000)^2 = 32{,}001{,}649/4{,}000{,}000 > 8 =
(\sqrt 8)^2$. If $\sqrt 8 \ge 5657/2000$ held, squaring both nonnegative sides would give
$8 \ge (5657/2000)^2$, a contradiction.
\end{proof}

The constant is the four-decimal round-up of $2\sqrt 2 = 2.8284271\ldots$; the gap is about
$7.3 \times 10^{-5}$. It lets rational code test ``above the Tsirelson bound'' without a square
root. It plays no part in the integer check of Section~\ref{sub:chsh-integer}.

\subsection{The zero-marginal \texorpdfstring{$5\times5$}{5 by 5} moment matrix}


The four averages fill only the Alice--Bob part of a larger table. Give that
table five labels: a constant reading 1, Alice's two possible readings, and
Bob's two. Each cell records an average product of its row and column
readings. A reading multiplied by itself is 1, so the diagonal is all ones.
The four averages measured fill the cross-player cells and their reflected
partners.

Other cells are still empty. This particular check fills them with zero:
each player's average answer against the constant reading, and the cells
between that player's two settings. These are choices of the check; the
four cross-player averages alone did not supply them. I will later let the
two within-player cells vary and ask whether any filling works. For now,
write down this one fixed filling and test it.

\begin{definition}[The pinned level-1 matrix]
Take rows and columns $I,A_0,A_1,B_0,B_1$. Put 1 on the diagonal and the four correlators in
the Alice--Bob cells, and set the four one-party entries (each player's own average answer) and
the two within-party entries ($A_0$ against $A_1$, $B_0$ against $B_1$) to zero. The resulting
matrix is:
\end{definition}

\[
M(E) \;=\;
\begin{array}{c@{\quad}c}
 & \begin{array}{ccccc} I & A_0 & A_1 & B_0 & B_1 \end{array} \\
\begin{array}{c} I \\ A_0 \\ A_1 \\ B_0 \\ B_1 \end{array} &
\begin{pmatrix}
1 & 0 & 0 & 0 & 0 \\
0 & 1 & 0 & E_{00} & E_{01} \\
0 & 0 & 1 & E_{10} & E_{11} \\
0 & E_{00} & E_{10} & 1 & 0 \\
0 & E_{01} & E_{11} & 0 & 1
\end{pmatrix}
\end{array}
\]

What do I ask of this table? Choose five numbers, one beside each row and
one beside the matching column. For each cell, multiply its entry by its row's
number and its column's number, then add all twenty-five results. Do this
with any five numbers you like, positive, negative or zero. I require the
total never to be negative. That requirement is called \emph{positive
semidefiniteness}, abbreviated PSD.

Write the five chosen numbers as a column $v$. The superscript $\mathsf T$
turns a column into a row, so ordinary matrix multiplication writes that same
weighted total as $v^{\mathsf T}Mv$. The condition is therefore
$v^{\mathsf T}Mv\ge0$ for every choice of $v$. The table is symmetric too:
reflect it across its diagonal and the entries agree.

The first row and column have only one nonzero entry, the corner 1. Their
contribution is the square of the first chosen number. The remaining four
rows and columns can be tested on their own. Grouping them gives

\[
M(E) \;=\; \begin{pmatrix} 1 & 0 \\ 0 & G \end{pmatrix},
\qquad
G \;=\; \begin{pmatrix} I_2 & E \\ E^{\mathsf T} & I_2 \end{pmatrix}.
\]

Here $I_2$ is the two-by-two identity table: ones on its diagonal and zeros
elsewhere. The superscript on $E$ again swaps rows and columns. The weighted
total for $M(E)$ is the first number's square plus the weighted total for $G$.
If every total for $G$ is nonnegative, adding a square keeps it nonnegative.
If some total for $G$ is negative, choose the first number to be zero and the
same negative total appears for $M(E)$. Thus the two tables pass or fail
together.

The next operation is completing a square, with several numbers at once.
For one number, $u^2+2uw+rw^2$ becomes $(u+w)^2+(r-1)w^2$.
The first square cannot hurt, and choosing $u=-w$ makes it vanish.
So the whole expression stays nonnegative for every choice exactly when
the remaining part does. The matrix formula below does the same job for
two groups of numbers.

Here \emph{positive definite} means that the weighted total for $P$ is
strictly positive for every nonzero choice of weights. In particular, $P$
has an inverse, written $P^{-1}$: applying it undoes multiplication by $P$.
The matrix left after completing the square is called the \emph{Schur
complement}. Its name is less useful here than the operation that produces it.

\begin{lemma}[Schur complement, the standard form]
Let $P$ be symmetric positive definite. The symmetric block matrix
$\begin{pmatrix} P & Q \\ Q^{\mathsf T} & R \end{pmatrix}$ is PSD if and only if
$R - Q^{\mathsf T} P^{-1} Q$ is PSD.
\end{lemma}

\begin{proof}
For $u, w$ of the right sizes,
\[
\begin{pmatrix} u \\ w \end{pmatrix}^{\!\mathsf T}
\begin{pmatrix} P & Q \\ Q^{\mathsf T} & R \end{pmatrix}
\begin{pmatrix} u \\ w \end{pmatrix}
= (u + P^{-1} Q w)^{\mathsf T} P \,(u + P^{-1} Q w)
+ w^{\mathsf T} (R - Q^{\mathsf T} P^{-1} Q)\, w.
\]
The first term is never negative because $P$ is positive definite, and it is zero at
$u = -P^{-1} Q w$. So the form is nonnegative for all $(u, w)$ exactly when the second term
is nonnegative for all $w$.
\end{proof}

The identity itself holds for any invertible symmetric $P$. One direction (block PSD gives the
complement PSD) needs only that: put $u = -P^{-1}Qw$. The other direction is where $P$ has to be
positive definite. For the pinned matrix I don't need the lemma in the abstract: I carry out its one instance with
$P = I_2$, $Q = E$, $R = I_2$, where the complement is
\[
I_2 - E^{\mathsf T} E =
\begin{pmatrix}
1 - E_{00}^2 - E_{10}^2 & -(E_{00} E_{01} + E_{10} E_{11}) \\
-(E_{00} E_{01} + E_{10} E_{11}) & 1 - E_{01}^2 - E_{11}^2
\end{pmatrix}.
\]
The diagonal of $E^{\mathsf T} E$ holds the squared lengths of the columns of $E$, which is
why the conditions below are column conditions.

\begin{definition}[Column contractive]
Write
$p = 1 - E_{00}^2 - E_{10}^2$, $q = 1 - E_{01}^2 - E_{11}^2$ and
$s = E_{00} E_{01} + E_{10} E_{11}$. The correlators are column contractive when
\[
p \ge 0, \qquad q \ge 0, \qquad p\,q - s^2 \ge 0.
\]
\end{definition}

\begin{lemma}[The $2\times2$ test]\label{lem:chsh-psd2}
If $\alpha \ge 0$, $\delta \ge 0$ and
$\alpha\delta - \beta^2 \ge 0$, then $\alpha u^2 + 2\beta u w + \delta w^2 \ge 0$ for all real
$u, w$.
\end{lemma}

\begin{proof}
If $w = 0$ the form is $\alpha u^2 \ge 0$. Otherwise put $t = u/w$, so the form is
$w^2(\alpha t^2 + 2\beta t + \delta)$. If $\alpha = 0$, then $-\beta^2 \ge 0$ forces
$\beta = 0$ and the bracket is $\delta \ge 0$. If $\alpha > 0$,
$\alpha t^2 + 2\beta t + \delta = \alpha (t + \beta/\alpha)^2 + (\alpha\delta - \beta^2)/\alpha$,
a sum of two nonnegative terms.
\end{proof}

\begin{theorem}[Contractive implies PSD]\label{thm:chsh-cc-psd}
If $p \ge 0$, $q \ge 0$ and $pq - s^2 \ge 0$, then $M(E)$ is positive semidefinite.
\end{theorem}

\begin{proof}
Write $v = (v_0, x_0, x_1, y_0, y_1)$. Reading off the entries of $M(E)$,
\[
v^{\mathsf T} M(E)\, v = v_0^2 + x_0^2 + x_1^2 + y_0^2 + y_1^2
+ 2x_0(E_{00} y_0 + E_{01} y_1) + 2x_1(E_{10} y_0 + E_{11} y_1).
\]
Complete the squares in $x_0$ and $x_1$:
\[
\begin{aligned}
v^{\mathsf T} M(E)\, v ={}& v_0^2 + (x_0 + E_{00} y_0 + E_{01} y_1)^2
+ (x_1 + E_{10} y_0 + E_{11} y_1)^2 \\
&+ p\, y_0^2 - 2 s\, y_0 y_1 + q\, y_1^2.
\end{aligned}
\]
To check the last three terms, expand the two squares: they contribute
$(E_{00} y_0 + E_{01} y_1)^2 + (E_{10} y_0 + E_{11} y_1)^2$ beyond the cross terms, which is
$(E_{00}^2 + E_{10}^2) y_0^2 + 2 s\, y_0 y_1 + (E_{01}^2 + E_{11}^2) y_1^2$; subtracting it
from $y_0^2 + y_1^2$ leaves $p\, y_0^2 - 2 s\, y_0 y_1 + q\, y_1^2$. The first three terms are
squares. The last is the $2 \times 2$ form with $\alpha = p$, $\beta = -s$, $\delta = q$,
nonnegative by Lemma~\ref{lem:chsh-psd2}.
\end{proof}

\begin{lemma}[Discriminant]\label{lem:chsh-discriminant}
If $a + 2bt + ct^2 \ge 0$ for every real $t$, then
$b^2 \le ac$.
\end{lemma}

\begin{proof}
If $c < 0$ the quadratic tends to $-\infty$, so $c \ge 0$. If $c = 0$, the line
$a + 2bt$ is nonnegative for every $t$ only if $b = 0$, and then $b^2 = 0 = ac$. If $c > 0$,
evaluate at $t = -b/c$: $a - b^2/c \ge 0$, so $b^2 \le ac$.
\end{proof}

\begin{theorem}[PSD implies contractive]\label{thm:chsh-psd-cc}
If $M(E)$ is positive semidefinite, then $p \ge 0$, $q \ge 0$ and $pq - s^2 \ge 0$.
\end{theorem}

\begin{proof}
Feed the quadratic form three test vectors, written in the order
$(I, A_0, A_1, B_0, B_1)$.
Each is the choice $x = -E y$ that zeroes both completed squares in the proof of
Theorem~\ref{thm:chsh-cc-psd}, so only $y^{\mathsf T}(I_2 - E^{\mathsf T} E) y$ is left.

First, $v = (0, -E_{00}, -E_{10}, 1, 0)$ gives $v^{\mathsf T} M v = p$, so $p \ge 0$.

Second, $v = (0, -E_{01}, -E_{11}, 0, 1)$ gives $v^{\mathsf T} M v = q$, so $q \ge 0$.

Third, for each real $t$, $v(t) = (0, -(E_{00} t + E_{01}), -(E_{10} t + E_{11}), t, 1)$
gives $v(t)^{\mathsf T} M v(t) = p\,t^2 - 2 s\,t + q \ge 0$. Apply Lemma~\ref{lem:chsh-discriminant}
with $a = q$, $b = -s$, $c = p$: $s^2 \le pq$.
\end{proof}

\begin{corollary}[The equivalence]
For all real $E_{00},E_{01},E_{10},E_{11}$, the symmetric matrix $M(E)$ is positive
semidefinite if and only if the correlators are column contractive.
\end{corollary}

\begin{proof}
Theorem~\ref{thm:chsh-cc-psd} is one direction and Theorem~\ref{thm:chsh-psd-cc} the other.
\end{proof}

The equivalence above is about the columns. The bound can go through either the rows or the columns: the six-square identity only needs the four squares to add to at most 2, and the two row bounds give that, and so do $p\ge0$ and $q\ge0$.

\begin{theorem}[PSD implies Tsirelson]
If the pinned matrix $M(E)$ is symmetric and positive semidefinite at
$E_{00},E_{01},E_{10},E_{11}$, then $S^2\le8$ and $|S|\le\sqrt8$.
\end{theorem}

\begin{proof}
Feed the form the vector that is 1 at $A_0$, $-E_{00}$ at $B_0$, $-E_{01}$ at $B_1$ and 0
elsewhere. Reading off the entries, the form is
$1 + E_{00}^2 + E_{01}^2 - 2E_{00}^2 - 2E_{01}^2 = 1 - E_{00}^2 - E_{01}^2$ (the $B_0B_1$ cell is
0), so $E_{00}^2 + E_{01}^2 \le 1$. The vector that is 1 at $A_1$, $-E_{10}$ at $B_0$ and
$-E_{11}$ at $B_1$ gives $E_{10}^2 + E_{11}^2 \le 1$ the same way. Then
Theorem~\ref{thm:chsh-row-tsirelson} applies.
\end{proof}

Here are the three conditions and the score at a few points, all by direct arithmetic.
In this table, $a = 3/5$, $b = 4/5$ and $t = 1/\sqrt 2$.

\begingroup\small
\begin{longtable}{@{}>{\raggedright\arraybackslash}p{0.214\linewidth}>{\raggedright\arraybackslash}p{0.122\linewidth}>{\raggedright\arraybackslash}p{0.122\linewidth}>{\raggedright\arraybackslash}p{0.122\linewidth}>{\raggedright\arraybackslash}p{0.122\linewidth}>{\raggedright\arraybackslash}p{0.122\linewidth}@{}}
\toprule
\textbf{point $(E_{00}, E_{01}, E_{10}, E_{11})$} & \textbf{$S$} & \textbf{$p$} & \textbf{$q$} & \textbf{$pq - s^2$} & \textbf{contrac-\newline tive} \\
\midrule
\endhead
\bottomrule
\endlastfoot
$(0, 0, 0, 0)$ & 0 & 1 & 1 & 1 & yes \\
$(0.6, 0.6, 0.6, -0.6)$ & 2.4 & 0.28 & 0.28 & 0.0784 & yes \\
$(a, b, b, -a)$ & 2.8 & 0 & 0 & 0 & yes, with no slack \\
$(t, t, t, -t)$ & $2\sqrt 2$ & 0 & 0 & 0 & yes, with no slack \\
$(1, 1, 1, 1)$, all-ones plan & 2 & $-1$ & $-1$ &  & no \\
$(1, 0, 1, 0)$, a classical mixture & 2 & $-1$ & 1 &  & no \\
$(1, 1, 1, -1)$, PR box & 4 & $-1$ & $-1$ &  & no \\
\end{longtable}
\endgroup

Every deterministic plan has all four correlators equal to $\pm 1$, so $p = 1 - 1 - 1 = -1$
and every one of the sixteen fails.

\begin{remark}[What PSD of the pinned matrix means]\label{rem:chsh-scope-five}
$M(E)$ is the level-1 moment matrix with the two within-party entries $\langle A_0 A_1
\rangle$ and $\langle B_0 B_1 \rangle$ pinned to zero. Its PSD is sufficient for the
correlators to admit some PSD completion (Section~\ref{sub:chsh-elliptope}), and it
is not necessary: the all-ones plan has a completion but fails here. So a check built on
$M(E)$ is sound for the completion set and not complete for it, and it rejects some honest
classical tallies. That correlators with a PSD completion are exactly those reachable by
quantum measurements is Tsirelson's theorem, and both halves are proved here in finite dimension.
For real and for complex amplitudes, a strategy's correlators are inner products of unit vectors
(with complex amplitudes, use real vectors of twice the length) and its level-1 moment matrix is PSD, so they have a PSD completion. The other half, that every PSD completion comes from a quantum
strategy, comes from this construction: the strategy is built on two four-level systems from
the completion's Gram vectors (Section~\ref{sub:four-vectors}).
\end{remark}

\subsection{The integer check}


The machine has whole-number tallies. The three conditions above were written
using averages, which are fractions. You can test the same conditions without
ever dividing: first make sure each denominator is positive, then multiply
it away. Multiplying by a positive number cannot turn a passing inequality
into a failing one or a failing one into a passing one.

For the first column, the condition starts as one minus two squared averages.
Multiplying by the product of their squared denominators leaves three
whole-number products. I call that result $A$. The second column gives $B$
the same way. The cross-column term is a sum of two products of averages;
multiplying by all four denominators leaves a whole number I call $C$.
The last condition becomes $C^2\le AB$, because both sides have had the same
positive denominator removed. Four positive bucket sizes and these three
comparisons make seven checks. Here are the products to calculate.

\begin{definition}[The seven-integer check]\label{def:chsh-check}
From the counters of Definition~\ref{def:chsh-counters}, with all
arithmetic in $\mathbb{Z}$, let
\[
\begin{aligned}
A &= N_{00}^2 N_{10}^2 - D_{00}^2 N_{10}^2 - D_{10}^2 N_{00}^2, \\
B &= N_{01}^2 N_{11}^2 - D_{01}^2 N_{11}^2 - D_{11}^2 N_{01}^2, \\
C &= D_{00} D_{01} N_{10} N_{11} + D_{10} D_{11} N_{00} N_{01}.
\end{aligned}
\]
The check is the conjunction of seven facts, in this order:
$0 < N_{00}$, $0 < N_{01}$, $0 < N_{10}$, $0 < N_{11}$, $0 \le A$, $0 \le B$, and
$C \cdot C \le A \cdot B$.
\end{definition}

There is no division, no real number, no square root and no $S$ in it.

\begin{lemma}[The seven facts are the three conditions with denominators cleared]
Assume $N_{00}, N_{01}, N_{10}, N_{11} > 0$ and put $E_{xy} = D_{xy}/N_{xy}$. Then, with
$p, q, s$ as in Definition~\ref{def:chsh-cc},
\[
\begin{gathered}
p = \frac{A}{N_{00}^2 N_{10}^2}, \qquad q = \frac{B}{N_{01}^2 N_{11}^2}, \qquad
s = \frac{C}{N_{00} N_{01} N_{10} N_{11}}, \\
pq - s^2 = \frac{AB - C^2}{(N_{00} N_{01} N_{10} N_{11})^2}.
\end{gathered}
\]
Because every denominator is positive, $0 \le A \iff p \ge 0$, $0 \le B \iff q \ge 0$ and
$C^2 \le AB \iff s^2 \le pq$.
\end{lemma}

\begin{proof}
For $p$: $1 - D_{00}^2/N_{00}^2 - D_{10}^2/N_{10}^2$; multiply by $N_{00}^2 N_{10}^2$ to get
$N_{00}^2 N_{10}^2 - D_{00}^2 N_{10}^2 - D_{10}^2 N_{00}^2 = A$. For $q$ the same with
columns $01$ and $11$, multiplied by $N_{01}^2 N_{11}^2$. For $s$:
\[
\frac{D_{00}}{N_{00}} \frac{D_{01}}{N_{01}} + \frac{D_{10}}{N_{10}} \frac{D_{11}}{N_{11}}
= \frac{D_{00} D_{01} N_{10} N_{11} + D_{10} D_{11} N_{00} N_{01}}{N_{00} N_{01} N_{10} N_{11}}
= \frac{C}{N_{00} N_{01} N_{10} N_{11}}.
\]
Then
$pq = AB / (N_{00}^2 N_{10}^2 N_{01}^2 N_{11}^2) = AB/(N_{00} N_{01} N_{10} N_{11})^2$ and
$s^2 = C^2/(N_{00} N_{01} N_{10} N_{11})^2$; subtract. Multiplying an inequality by a
positive number preserves it in both directions.
\end{proof}

The check has the following guarantee:

\begin{theorem}[Soundness of the integer check]
If a tally $w$ passes the seven-integer check, then its four correlators
$E_{xy}=D_{xy}/N_{xy}$ are column contractive. Consequently its pinned matrix is positive
semidefinite.
\end{theorem}

\begin{proof}
The first four tests say that every bucket contains a round, so each average
is its signed total divided by a positive count. The next two tests say that
$A$ and $B$ are nonnegative. Dividing each by its positive denominator gives
$p\ge0$ and $q\ge0$. The last test says $AB-C^2\ge0$. Divide that difference
by the positive square of the product of all four counts; the result is
$pq-s^2\ge0$, as the preceding calculation shows. Those are exactly the
three column conditions, and Theorem~\ref{thm:chsh-cc-psd} shows that they
make the matrix pass the nonnegative-total test.
\end{proof}

\begin{corollary}[A passed check implies the Tsirelson bound]
If a tally $w$ passes the seven-integer check, the correlators of $w$ satisfy
$S^2 \le 8$. It is the composition of Theorem~\ref{thm:chsh-check-sound},
Corollary~\ref{cor:chsh-iff} and Theorem~\ref{thm:chsh-psd-tsirelson}.\proofclosed \end{corollary} Two of the seven facts carry the bound by themselves. With every pair sampled, $0\le A$ and $0\le B$ are $p\ge0$ and $q\ge0$, so the four squared correlators add to at most 2 and Lemma~\ref{lem:chsh-sos} gives $S^2\le8$. The fact $C\cdot C\le A\cdot B$ is what positive semidefiniteness adds, and the bound doesn't use it.
For the form this takes on the small machine, where a raised flag has to stand on a passing
check, see Appendix~\ref{sec:chsh}.

A worked tally. Take $(\mathrm{same}, \mathrm{diff}) = (4,1), (4,1), (4,1), (1,4)$ for the
pairs $00, 01, 10, 11$. Then $N_{xy} = 5$ for all four, $D_{00} = D_{01} = D_{10} = 3$,
$D_{11} = -3$, and the correlators are $0.6, 0.6, 0.6, -0.6$ with $S = 2.4$. Compute:
\[
\begin{aligned}
A &= 25 \cdot 25 - 9 \cdot 25 - 9 \cdot 25 = 625 - 225 - 225 = 175, \\
B &= 25 \cdot 25 - 9 \cdot 25 - 9 \cdot 25 = 175, \\
C &= 3 \cdot 3 \cdot 5 \cdot 5 + 3 \cdot (-3) \cdot 5 \cdot 5 = 225 - 225 = 0.
\end{aligned}
\]
All seven facts hold: four positive $N$, $0 \le 175$ twice, and $0 \le 175 \cdot 175$.
Dividing back, $p = 175/625 = 0.28$ and $s = 0$, the second row of the table in Section~\ref{sub:chsh-five}.

Two more, for contrast. One round per pair from the all-ones plan, $(1,0)$ four times, gives
$N = D = 1$ everywhere and $A = 1 - 1 - 1 = -1 < 0$: the check fails. The tally $(4,1), (9,1), (9,1),
(1,4)$ gives correlators $3/5, 4/5, 4/5, -3/5$, with $N = 5, 10, 10, 5$ and
$D = 3, 8, 8, -3$:
$A = 25 \cdot 100 - 9 \cdot 100 - 64 \cdot 25 = 0$,
$B = 100 \cdot 25 - 64 \cdot 25 - 9 \cdot 100 = 0$ and
$C = 3 \cdot 8 \cdot 10 \cdot 5 + 8 \cdot (-3) \cdot 5 \cdot 10 = 0$. It passes with no
slack, at $S = 2.8$.

\subsection{Level 1+AB}


The first table had one place for each player's answer. Enlarge it by giving
each Alice--Bob product a place too. There are four such products, so the
five places become nine. A cell still asks for the average product of its
row and column labels. Some of those products reduce to entries I already
have; five more numbers are needed for the particular filling below.
Call them $\gamma_1$ through $\gamma_5$. The longer labels say which
products they stand for.

The five additional parameters name the moments
$\gamma_1=\langle A_0 A_1 B_0\rangle$,
$\gamma_2=\langle A_0 A_1 B_1\rangle$,
$\gamma_3=\langle A_0 B_0 B_1\rangle$,
$\gamma_4=\langle A_1 B_0 B_1\rangle$ and
$\gamma_5=\langle A_0 A_1 B_0 B_1\rangle$.
These are labels for the entries; the definition below does not construct operators.

\begin{definition}[The level-1+AB matrix]
The $9 \times 9$ real matrix on rows and
columns $I, A_0, A_1, B_0, B_1, A_0B_0, A_0B_1, A_1B_0, A_1B_1$, with the following entries:
\end{definition}

\[
\setlength{\arraycolsep}{3.2pt}
\begin{array}{c|ccccccccc}
 & I & A_0 & A_1 & B_0 & B_1 & A_0B_0 & A_0B_1 & A_1B_0 & A_1B_1 \\ \hline
I      & 1 & 0 & 0 & 0 & 0 & E_{00} & E_{01} & E_{10} & E_{11} \\
A_0    & 0 & 1 & 0 & E_{00} & E_{01} & 0 & 0 & \gamma_1 & \gamma_2 \\
A_1    & 0 & 0 & 1 & E_{10} & E_{11} & \gamma_1 & \gamma_2 & 0 & 0 \\
B_0    & 0 & E_{00} & E_{10} & 1 & 0 & 0 & \gamma_3 & 0 & \gamma_4 \\
B_1    & 0 & E_{01} & E_{11} & 0 & 1 & \gamma_3 & 0 & \gamma_4 & 0 \\
A_0B_0 & E_{00} & 0 & \gamma_1 & 0 & \gamma_3 & 1 & 0 & 0 & \gamma_5 \\
A_0B_1 & E_{01} & 0 & \gamma_2 & \gamma_3 & 0 & 0 & 1 & -\gamma_5 & 0 \\
A_1B_0 & E_{10} & \gamma_1 & 0 & 0 & \gamma_4 & 0 & -\gamma_5 & 1 & 0 \\
A_1B_1 & E_{11} & \gamma_2 & 0 & \gamma_4 & 0 & \gamma_5 & 0 & 0 & 1
\end{array}
\]

Two constraints are fixed by this layout. First, the four-body cells carry opposite signs:
$+\gamma_5$ at $(A_0B_0, A_1B_1)$ and $-\gamma_5$ at $(A_0B_1, A_1B_0)$. In a real
moment matrix those two cells are $\langle A_0 A_1 B_0 B_1 \rangle$ and $\langle A_0 A_1 B_1 B_0
\rangle$, so the tie says $\langle A_0 A_1 (B_0 B_1 + B_1 B_0) \rangle = 0$, which for a real
matrix is $\langle \{A_0, A_1\} \otimes \{B_0, B_1\} \rangle = 0$. The level-1 pinnings don't
imply it. Take a classical strategy: Alice's two answers $a_0, a_1$ are independent fair signs,
and Bob copies them, $b_0 = a_0$ and $b_1 = a_1$. Every pinned cell of its moment matrix is 0 and
both four-body cells are 1, so no choice of the nine numbers gives its matrix. So the matrix
here is level 1+AB cut down to the slice $\langle \{A_0, A_1\} \otimes \{B_0, B_1\} \rangle = 0$,
and that is what ``level 1+AB'' names in this book. The slice holds whenever $A_0$ and $A_1$ anticommute.
The soundness theorems below are about this matrix, so they stand as stated. Second, the matrix
keeps the level-1 pinning: every cell whose label reduces to a marginal or to
$\langle A_0 A_1 \rangle$ or $\langle B_0 B_1 \rangle$ is zero, for instance
$(A_0, A_0B_0)$, $(A_0B_0, A_0B_1)$ and $(A_0B_0, A_1B_0)$. The displayed matrix is symmetric.

I want to test this larger table with the same operation as before:
put a weight beside each row and column and add the weighted cells. The
first three weights will each be collected into one completed square.
Everything involving the other six will be left over. That leftover is
what I call the residual; it is the part whose sign I still need to settle.

Write a vector as $(v_I, v_{A_0}, v_{A_1}, v_{B_0}, v_{B_1}, v_{00}, v_{01}, v_{10}, v_{11})$,
where $v_{xy}$ is the coordinate of $A_xB_y$. Put
\[
\begin{aligned}
L_1 &= E_{00} v_{00} + E_{01} v_{01} + E_{10} v_{10} + E_{11} v_{11}, \\
L_2 &= E_{00} v_{B_0} + E_{01} v_{B_1} + \gamma_1 v_{10} + \gamma_2 v_{11}, \\
L_3 &= E_{10} v_{B_0} + E_{11} v_{B_1} + \gamma_1 v_{00} + \gamma_2 v_{01}.
\end{aligned}
\]

\begin{definition}[The residual]
\[
\begin{aligned}
R ={}& v_{B_0}^2 + v_{B_1}^2 + v_{00}^2 + v_{01}^2 + v_{10}^2 + v_{11}^2 \\
&+ 2 v_{B_0} (\gamma_3 v_{01} + \gamma_4 v_{11}) + 2 v_{B_1} (\gamma_3 v_{00} + \gamma_4 v_{10}) \\
&+ 2 \gamma_5 (v_{00} v_{11} - v_{01} v_{10}) - L_1^2 - L_2^2 - L_3^2.
\end{aligned}
\]
The condition called column contractivity at level 1+AB is $R \ge 0$ for all six real inputs.
\end{definition}

\begin{lemma}[Three squares plus the residual]\label{lem:chsh-nine-sos}
For every vector,
\[
v^{\mathsf T} M_9\, v = (v_I + L_1)^2 + (v_{A_0} + L_2)^2 + (v_{A_1} + L_3)^2 + R.
\]
\end{lemma}

\begin{proof}
The coordinates $v_I$, $v_{A_0}$, $v_{A_1}$ appear in the form only through their diagonal
squares and the cross terms $2 v_I L_1$, $2 v_{A_0} L_2$, $2 v_{A_1} L_3$, read off rows
$I$, $A_0$, $A_1$ of the matrix. Completing each square adds $L_1^2 + L_2^2 + L_3^2$, which
the residual subtracts. What is left in $R$ is the diagonal of the remaining six coordinates
and the cross terms among them, read off the lower right $6 \times 6$ block.
\end{proof}

\begin{theorem}[PSD at level 1+AB]
For all real $E_{xy}$ and $\gamma_1, \ldots, \gamma_5$,
the matrix of Definition~\ref{def:chsh-nine} is positive semidefinite if and only if
the residual $R$ is nonnegative for every choice of its six real inputs.
\end{theorem}

\begin{proof}
If $R \ge 0$ everywhere, Lemma~\ref{lem:chsh-nine-sos} writes the form as three squares plus $R$.
Conversely, given the six residual inputs, choose $v_I = -L_1$, $v_{A_0} = -L_2$,
$v_{A_1} = -L_3$; the three squares vanish and PSD gives $R \ge 0$. This is the level-1
argument of Section~\ref{sub:chsh-five} one size up.
\end{proof}

As a quadratic form in $(v_{B_0},v_{B_1},v_{00},v_{01},v_{10},v_{11})$, numbered 1 to 6,
the residual has the symmetric coefficient matrix $H$ with these upper entries. Each entry
can be read back off $R$ by collecting the coefficient of $v_i v_j$.

\begin{align*}
H_{11} &= 1 - E_{00}^2 - E_{10}^2, & H_{22} &= 1 - E_{01}^2 - E_{11}^2, \\
H_{12} &= -(E_{00}E_{01} + E_{10}E_{11}), & H_{13} &= -E_{10}\gamma_1, \\
H_{14} &= \gamma_3 - E_{10}\gamma_2, & H_{15} &= -E_{00}\gamma_1, \\
H_{16} &= \gamma_4 - E_{00}\gamma_2, & H_{23} &= \gamma_3 - E_{11}\gamma_1, \\
H_{24} &= -E_{11}\gamma_2, & H_{25} &= \gamma_4 - E_{01}\gamma_1, \\
H_{26} &= -E_{01}\gamma_2, & H_{33} &= 1 - E_{00}^2 - \gamma_1^2, \\
H_{34} &= -(E_{00}E_{01} + \gamma_1\gamma_2), & H_{35} &= -E_{00}E_{10}, \\
H_{36} &= \gamma_5 - E_{00}E_{11}, & H_{44} &= 1 - E_{01}^2 - \gamma_2^2, \\
H_{45} &= -\gamma_5 - E_{01}E_{10}, & H_{46} &= -E_{01}E_{11}, \\
H_{55} &= 1 - E_{10}^2 - \gamma_1^2, & H_{56} &= -(E_{10}E_{11} + \gamma_1\gamma_2), \\
H_{66} &= 1 - E_{11}^2 - \gamma_2^2.
\end{align*}

The top left $2 \times 2$ block is $I_2 - E^{\mathsf T} E$ from Section~\ref{sub:chsh-five}.

The full check does not decide $R \ge 0$ directly. It tests the stronger fact that $H$ is
positive definite, by two Schur steps and a Sylvester test, all fraction-free.

\begin{lemma}[Fraction-free Schur step]\label{lem:chsh-schur-step}
For a symmetric $n \times n$ matrix $h$ and any vector $v$, with
$S_{ij} = h_{11} h_{ij} - h_{1i} h_{1j}$ for $i, j \ge 2$ (here $S$ is the Schur matrix;
the score plays no part
in this lemma),
\[
h_{11} \, v^{\mathsf T} h\, v = \Bigl(\sum_{j} h_{1j} v_j\Bigr)^2 + \sum_{i,j \ge 2} S_{ij} v_i v_j.
\]
So if $h_{11} > 0$ and the $(n-1) \times (n-1)$ matrix $S$ gives a nonnegative form, then $h$
does too. The cascade uses this identity at sizes six and five.
\end{lemma}

\begin{proof}
Expand the square: $h_{11}^2 v_1^2 + 2 h_{11} v_1 \sum_{j \ge 2} h_{1j} v_j +
(\sum_{j \ge 2} h_{1j} v_j)^2$. The left side is
$h_{11}^2 v_1^2 + 2 h_{11} v_1 \sum_{j \ge 2} h_{1j} v_j + h_{11} \sum_{i,j \ge 2} h_{ij} v_i v_j$.
The difference is $\sum_{i,j \ge 2} (h_{11} h_{ij} - h_{1i} h_{1j}) v_i v_j$. Dividing by
$h_{11} > 0$ keeps the sign.
\end{proof}

\begin{lemma}[Sylvester for $4\times4$]\label{lem:chsh-sylvester}
Let $h$ be symmetric $4 \times 4$ with leading principal
minors $d_1, d_2, d_3, d_4$. If $d_1, d_2, d_3$ are strictly positive and $d_4 \ge 0$, then $v^{\mathsf T} h\, v \ge 0$ for all $v$.
\end{lemma}

\begin{proof}
Use the fraction-free factorization
\[
d_1^2 d_2 d_3 \, v^{\mathsf T} h\, v = d_1 d_2 d_3 \, P_1^2 + d_1 d_3 \, Q_1^2 + Q_2^2 +
d_1^2 d_2 d_4 \, v_4^2,
\]
where, writing $e_{ij}=h_{11}h_{ij}-h_{1i}h_{1j}$,
\[
\begin{aligned}
P_1 &= \sum_j h_{1j}v_j,\\
Q_1 &= d_2v_2+e_{23}v_3+e_{24}v_4,\\
Q_2 &= d_1d_3v_3+(d_2e_{34}-e_{23}e_{24})v_4.
\end{aligned}
\]
The identity is Lemma~\ref{lem:chsh-schur-step} three times over, each time on the matrix the
last one left. With $f_{ij}=e_{22}e_{ij}-e_{2i}e_{2j}$ for $i,j\ge3$, the lemma gives
$d_1\,v^{\mathsf T}h\,v = P_1^2+\sum_{i,j\ge2}e_{ij}v_iv_j$, then
$e_{22}\sum_{i,j\ge2}e_{ij}v_iv_j = Q_1^2+\sum_{i,j\ge3}f_{ij}v_iv_j$, then
$f_{33}\sum_{i,j\ge3}f_{ij}v_iv_j = Q_2^2+(f_{33}f_{44}-f_{34}^2)\,v_4^2$. The corners are
minors of $h$, by one fact about determinants: the $k\times k$ leading minor of $e$ is
$h_{11}^{k-1}$ times the $(k+1)\times(k+1)$ leading minor of $h$. To see it, multiply rows 2 to
$k+1$ of that block of $h$ by $h_{11}$, which multiplies its determinant by $h_{11}^k$, then
subtract $h_{1i}$ times the first row from row $i$, which changes nothing; those rows now read
$0, e_{i2}, \ldots, e_{i,k+1}$, so the determinant is $h_{11}$ times the leading minor of $e$.
So $e_{22}=d_2$ and $f_{33}=e_{22}e_{33}-e_{23}^2=d_1d_3$, and the same fact applied to $e$,
whose corner is $d_2$, gives $f_{33}f_{44}-f_{34}^2=d_2\det e=d_1^2d_2d_4$. Multiply the first
equation by $d_1d_2d_3$ and substitute the other two. With $d_1, d_2, d_3 > 0$ and $d_4 \ge 0$ the right side is nonnegative and the factor on the left is positive.
\end{proof}

The cascade for the full check runs on $H$. Step one: $H_{11} > 0$ and the $5 \times 5$
matrix $S^{(6)}_{ij} = H_{11} H_{ij} - H_{1i} H_{1j}$, $2 \le i, j \le 6$, gives a nonnegative
form. Step two: $S^{(6)}_{22} > 0$ and the $4 \times 4$ matrix
$S^{(5)}_{ij} = S^{(6)}_{22} S^{(6)}_{ij} - S^{(6)}_{2i} S^{(6)}_{2j}$, $3 \le i, j \le 6$,
gives a nonnegative form. Step three: the four leading minors of $S^{(5)}$ are positive. Together these conditions make the residual nonnegative, and so the nine-by-nine matrix positive semidefinite.

\begin{definition}[What the integer check computes]\label{def:chsh-g12345}
The supplied moments arrive as bucket pairs: $\gamma_k = N_{g_k}/D_{g_k}$ with
$N_{g_k} = \mathrm{same}_{g_k} - \mathrm{diff}_{g_k}$ and
$D_{g_k} = \mathrm{same}_{g_k} + \mathrm{diff}_{g_k}$. Watch the letters: for these five
buckets $N$ names the numerator and $D$ the denominator: the difference and the count, the reverse
of Section~\ref{sub:chsh-notation}. Let
$K = (N_{00} N_{01} N_{10} N_{11} D_{g_1} D_{g_2} D_{g_3} D_{g_4} D_{g_5})^2$. The check computes the 21 integers
$K H_{ij}$, the 15
integers of $S^{(6)}$ at scale $K^2$ and the 10 integers of $S^{(5)}$ at scale $K^4$, each by
$h_{11} h_{ij} - h_{1i} h_{1j}$ on integers. The
decider is the conjunction of: $N_{xy} > 0$ for the four
pairs; $D_{g_k} > 0$ for $k = 1..5$; $-D_{g_k} < N_{g_k} < D_{g_k}$ for $k = 1..5$; and six
strict positivities, namely $K H_{11} > 0$, $K^2 S^{(6)}_{22} > 0$, and the four leading
minors of the integer $S^{(5)}$. The full check
combines the seven-integer check of Definition~\ref{def:chsh-check} with this decider.
\end{definition}

Each cleared integer, read as a real, equals the
corresponding real quantity times the stated positive power of $K$. A $k \times k$ leading
minor of a matrix scaled by $c > 0$ scales by $c^k$. So every integer positivity holds exactly
when the real one does.

Two parts of the full check are there to fail early, and no proof below needs them. A passing
decider makes $H$ positive definite (Remark~\ref{rem:chsh-interior}), and the top left two-by-two
block of $H$ has $p$ and $q$ on its diagonal and $-s$ off it, so $p > 0$, $q > 0$ and $pq > s^2$,
and the seven-integer check passes. And with $H$ positive definite, the nine-by-nine form is zero
only at zero: it is three squares plus $R$ (Lemma~\ref{lem:chsh-nine-sos}), so where it is zero the
six inputs of $R$ are zero, then so are $L_1$, $L_2$ and $L_3$, and then so are $v_I$, $v_{A_0}$
and $v_{A_1}$. So the nine-by-nine matrix is positive definite; with ones on its diagonal, that puts every $\gamma_k$ strictly between
$-1$ and 1, which is the bucket bound $-D_{g_k} < N_{g_k} < D_{g_k}$.

\begin{theorem}[The full 1+AB check implies PSD]\label{thm:chsh-g12345}
Let the counts make every $N_{xy}$ and every $D_{g_k}$ positive, and let
the decider just defined pass on them. Then the matrix of
Definition~\ref{def:chsh-nine} at the four correlators and $\gamma_k = N_{g_k}/D_{g_k}$ is
symmetric and positive semidefinite.
\end{theorem}

\begin{proof}
The bucket conjuncts make every denominator positive, the scaling remark just above carries the
six integer positivities over to the real cascade, and Lemma~\ref{lem:chsh-schur-step} twice plus
Lemma~\ref{lem:chsh-sylvester} give $R \ge 0$, and so positive semidefiniteness by Theorem~\ref{thm:chsh-nine-iff}.
\end{proof}

\begin{remark}[Interior only]\label{rem:chsh-interior}
Every test in the cascade is strict, so a passing check certifies that $H$ is positive
definite, which is stronger than semidefinite. A point where $H$ is PSD but singular cannot pass, and
neither can a moment with $|\gamma_k| = 1$. For example, the tally of Section~\ref{sub:chsh-integer} with correlators
$(3/5, 4/5, 4/5, -3/5)$ has $H_{11} = 1 - 9/25 - 16/25 = 0$, so the full 1+AB check fails on it
for every choice of $\gamma$, though the level-1 check of Section~\ref{sub:chsh-integer} accepts it. The argument establishes soundness only; no converse is stated.
\end{remark}

A worked example of the cascade in exact rational arithmetic. Take the correlators $(0.6, 0.6, 0.6, -0.6)$ and
$\gamma = (0, 0, 0, 0, -1/2)$; as buckets, $\gamma_1$ to $\gamma_4$ are $(1,1)$, giving
$0/2$, and $\gamma_5$ is $(1,3)$, giving $-2/4$, both inside the strict bounds. Then
$H_{11} = H_{22} = 7/25$, $H_{12} = 0$, and every entry $H_{1j}$, $H_{2j}$ with $j \ge 3$ is
zero. So $S^{(6)}_{22} = (7/25)^2 = 49/625$, and $S^{(5)}$ is a positive multiple of the lower
right block of $H$:
\[
S^{(5)} = \frac{343}{781250}
\begin{pmatrix}
32 & -18 & -18 & -7 \\
-18 & 32 & 7 & 18 \\
-18 & 7 & 32 & 18 \\
-7 & 18 & 18 & 32
\end{pmatrix},
\]
where $50 H_{33} = 32$, $50 H_{34} = 50 H_{35} = -18$, $50 H_{36} = 50(-1/2 + 0.36) = -7$,
$50 H_{45} = 50(1/2 - 0.36) = 7$, and $50 H_{46} = 50 H_{56} = 18$. The leading minors of the
integer matrix are $32$, $700$, $15000$ and $140625$, all positive. With
$H_{11} = 7/25 > 0$ and $S^{(6)}_{22} > 0$, all six tests pass, and with Section~\ref{sub:chsh-integer}'s $A = B = 175$,
$C = 0$ the full check passes. At $\gamma_5 = 0$ instead, the same block becomes
$I_4 - \mathbf e \mathbf e^{\mathsf T}$ with $\mathbf e = (0.6, 0.6, 0.6, -0.6)$, whose third
leading minor is negative ($\tfrac{1}{25^3}$ times $-1250$), and the check fails.

That last fact is general. The zero-moment check takes no $\gamma$ buckets; it combines
the level-1 check with the integer form of
$E_{00}^2 + E_{01}^2 + E_{10}^2 + E_{11}^2 \le 1$ with denominators cleared by
$(N_{00} N_{01} N_{10} N_{11})^2$.

\begin{theorem}[The zero-moment check sits inside the classical bound]\label{thm:chsh-gzero}
If the zero-moment check accepts a tally $w$, then the four correlators of $w$ satisfy
$E_{00}^2 + E_{01}^2 + E_{10}^2 + E_{11}^2 \le 1$
 and $|S| \le 2$.
\end{theorem}

\begin{proof}
The check's extra conjunct, divided back by the positive factor
$(N_{00} N_{01} N_{10} N_{11})^2$, is $r = E_{00}^2 + E_{01}^2 + E_{10}^2 + E_{11}^2 \le 1$.
Then Lemma~\ref{lem:chsh-sos} gives $S^2 \le 4 r \le 4$.
\end{proof}

The level-1 half of that check follows from $r \le 1$ on its own: $p$ and $q$ are each at least
$1 - r$, and $pq - s^2 = 1 - r + (E_{00}E_{11} - E_{01}E_{10})^2$.

This is exactly the level-1+AB condition at $\gamma = 0$. There the residual splits into the
level-1 form in $(v_{B_0}, v_{B_1})$ and
$|\mathbf v|^2 - \langle \mathbf e, \mathbf v \rangle^2$ in
$\mathbf v = (v_{00}, v_{01}, v_{10}, v_{11})$, where
$\mathbf e = (E_{00}, E_{01}, E_{10}, E_{11})$. The second piece is nonnegative for every
$\mathbf v$ exactly when $r \le 1$: if $r \le 1$, Cauchy--Schwarz gives
$\langle \mathbf e, \mathbf v \rangle^2 \le r\,|\mathbf v|^2 \le |\mathbf v|^2$; if $r > 1$, the
choice $\mathbf v = \mathbf e$ gives $r - r^2 < 0$.

So the zero-moment check certifies a strict part of the classical region: no point
above $|S| = 2$, and none of the sixteen deterministic points either, since each has
$|\mathbf e|^2 = 4$. A point above 2 needs supplied moments.
The matrix is positive semidefinite at correlators $(3/5, 3/5, 3/5, -3/5)$,
$\gamma = (0,0,0,0,-1/2)$, with $S = 12/5$. The cleanest way to see it is the $\gamma_5$ slice,
$\gamma_1 = \cdots = \gamma_4 = 0$. There the residual splits into the level-1 form in
$(v_{B_0}, v_{B_1})$ and the form of $J - \mathbf e\mathbf e^{\mathsf T}$ in the four $v_{xy}$,
where $J$ is $I_4$ with $\gamma_5$ at $(00, 11)$ and $-\gamma_5$ at $(01, 10)$. For
$-1 < \gamma_5 < 1$, $J$ is positive definite, and the second form is nonnegative exactly when
$\mathbf e^{\mathsf T} J^{-1} \mathbf e \le 1$: apply Lemma~\ref{lem:chsh-schur} to the block
matrix $\begin{pmatrix} 1 & \mathbf e^{\mathsf T} \\ \mathbf e & J \end{pmatrix}$ once with the corner
$1$ as $P$ and once, rows reordered, with $J$ as $P$. Inverting the two $2 \times 2$ blocks of $J$ and multiplying out, that is
\[
\begin{aligned}
&(1 - \gamma_5)\bigl[(E_{00} + E_{11})^2 + (E_{01} - E_{10})^2\bigr] \\
&\quad + (1 + \gamma_5)\bigl[(E_{00} - E_{11})^2 + (E_{01} + E_{10})^2\bigr]
\le 2(1 - \gamma_5^2).
\end{aligned}
\]
At this point the left side is $\tfrac32 \cdot 0 + \tfrac12 \cdot (\tfrac{36}{25} +
\tfrac{36}{25}) = \tfrac{36}{25} = 1.44$ and the right side is $2 \cdot \tfrac34 = 1.5$, and the
level-1 part holds ($p = q = 0.28$, $s = 0$). The slice test is an illustration, not a third
check: on the counters $(4,1), (4,1), (4,1), (1,4)$ and the bucket $(1,3)$, the same comparison
with the denominators multiplied out says yes, and the full check of
Definition~\ref{def:chsh-g12345} accepts the same counters with $\gamma_1$ to $\gamma_4$ given as
the buckets $(1,1)$, as the worked example above shows.

Both 1+AB checks, the full check and the zero-moment check, contain
the seven-integer check as a conjunct, so a passing 1+AB check also passes the
level-1 check and inherits $S^2 \le 8$ from Section~\ref{sub:chsh-integer}. This follows directly by composing the definitions and the bound.

\subsection{The elliptope completion}


Go back to the two within-player cells I filled with zero. Leave them blank
this time. Keep all four measured cross-player averages fixed, and ask whether
there is any way to fill those two blanks that makes the table pass. A failure
with zeros no longer settles the question; another filling may work.

Call the first missing number $x$ and the second $y$, and copy each into its
reflected cell so the table stays symmetric. The collection of averages that
admit a passing filling has a name, the \emph{elliptope} completion set.
Membership means that at least one choice of those two numbers works. It
does not require every choice to work, or require the zero choice to work.

\begin{definition}[Completion and membership]
Complete the matrix of Section~\ref{sub:chsh-five} by setting the
within-party entries to free reals: $x$ in the $(A_0, A_1)$ cells and $y$ in the
$(B_0, B_1)$ cells, marginals still zero. The correlators are elliptope realizable
 when some $x, y$ make the completed matrix symmetric and PSD.
\end{definition}

\[
M(E; x, y) \;=\;
\begin{pmatrix}
1 & 0 & 0 & 0 & 0 \\
0 & 1 & x & E_{00} & E_{01} \\
0 & x & 1 & E_{10} & E_{11} \\
0 & E_{00} & E_{10} & 1 & y \\
0 & E_{01} & E_{11} & y & 1
\end{pmatrix}
\]

Its quadratic form at $v = (v_0, \ldots, v_4)$ is
$\sum_k v_k^2 + 2x v_1 v_2 + 2y v_3 v_4 + 2E_{00} v_1 v_3 + 2E_{01} v_1 v_4 + 2E_{10} v_2 v_3
+ 2E_{11} v_2 v_4$. The pinned matrix is the case $x = y = 0$,
so PSD of $M(E)$ implies membership.

\begin{theorem}[Deterministic plans are inside]\label{thm:chsh-ell-det}
If $a_0^2 = a_1^2 = b_0^2 = b_1^2 = 1$, then
$(a_0 b_0, a_0 b_1, a_1 b_0, a_1 b_1)$ is elliptope realizable, with $x = a_0 a_1$ and
$y = b_0 b_1$.
\end{theorem}

\begin{proof}
With these entries the form equals
$v_0^2 + (a_0 v_1 + a_1 v_2 + b_0 v_3 + b_1 v_4)^2$. Expanding the second square gives
$a_0^2 v_1^2 + \cdots$ on the diagonal, which is $v_1^2 + \cdots$ by the unit squares, and the
cross terms $2 a_0 a_1 v_1 v_2 = 2x v_1 v_2$, $2 a_0 b_0 v_1 v_3 = 2E_{00} v_1 v_3$, and so on
for all six pairs. A sum of two squares is never negative.
\end{proof}

The all-ones plan is one instance. The classical mixture $(1,0,1,0)$ has a completion at
$x=1$, $y=0$, with form $v_0^2+(v_1+v_2+v_3)^2+v_4^2$.

\begin{theorem}[Convexity and finite mixtures]\label{thm:chsh-ell-convex}
For $0 \le \lambda \le 1$, if $E$ and $F$ are elliptope realizable, so is
$\lambda E + (1 - \lambda) F$. For finitely
many members with nonnegative weights summing to at most 1, the weighted sum is realizable. Every finite mixture of deterministic plans is realizable.
\end{theorem}

\begin{proof}
The form is affine in $(E, x, y)$ for each fixed $v$. Take the completion
$(\lambda x_1 + (1-\lambda) x_2, \lambda y_1 + (1-\lambda) y_2)$; the form at the mixture is
the same mixture of two nonnegative forms. For several members, weight each one's completed matrix by its assigned
share and add them. If the shares total less than one, give the missing
share to the zero-correlator point with both missing entries zero; its
completed matrix is the identity. The diagonal of the weighted matrix is
then one, its cross entries are the desired weighted correlators, and every
weighted total is a sum of nonnegative totals with nonnegative weights.
It therefore supplies the required completion.
The last statement combines this with Theorem~\ref{thm:chsh-ell-det}.
\end{proof}

\begin{theorem}[Tsirelson's bound holds on the completion set]\label{thm:chsh-ell-tsirelson}
If $E$ is elliptope realizable then $S^2 \le 8$.
\end{theorem}

\begin{proof}
For a symmetric PSD matrix the bilinear form obeys Cauchy-Schwarz,
$\mathcal B(u, w)^2 \le \mathcal Q(u)\, \mathcal Q(w)$, because
$\mathcal Q(u + t w) = \mathcal Q(u) + 2t\,\mathcal B(u,w) + t^2 \mathcal Q(w) \ge 0$ for all
$t$ and Lemma~\ref{lem:chsh-discriminant} applies. Take $u_1 = e_{A_0}$ and
$w_+ = e_{B_0} + e_{B_1}$: $\mathcal B(u_1, w_+) = E_{00} + E_{01}$, $\mathcal Q(u_1) = 1$,
$\mathcal Q(w_+) = 2 + 2y$. Take $u_2 = e_{A_1}$ and $w_- = e_{B_0} - e_{B_1}$:
$\mathcal B(u_2, w_-) = E_{10} - E_{11}$, $\mathcal Q(w_-) = 2 - 2y$. So
$c_1^2 \le 2 + 2y$ and $c_2^2 \le 2 - 2y$ with $c_1 = E_{00} + E_{01}$,
$c_2 = E_{10} - E_{11}$, and both budgets are nonnegative. Then
$c_1^2 c_2^2 \le (2 + 2y)(2 - 2y) = 4 - 4y^2 \le 4$, so $c_1 c_2 \le 2$, and
$S^2 = (c_1 + c_2)^2 = c_1^2 + 2 c_1 c_2 + c_2^2 \le (2 + 2y) + 4 + (2 - 2y) = 8$.
\end{proof}

The completion can make a row of $E$ longer than 1 (the all-ones plan has rows of squared
length 2), which is why this proof does not use the row bounds of Section~\ref{sub:chsh-five}.

\begin{corollary}[The PR box is outside]\label{cor:chsh-ell-pr}
 $(1, 1, 1, -1)$ is not elliptope realizable, since its
$S^2 = 16 > 8$.\proofclosed
\end{corollary}

\begin{theorem}[Strictly beyond the classical bound]\label{thm:chsh-ell-beyond}
 $(3/5, 3/5, 3/5, -3/5)$ is elliptope realizable and has
$S = 12/5 > 2$.
\end{theorem}

\begin{proof}
At $x = y = 0$ the form equals
$v_0^2 + (v_1 + \tfrac35 v_3 + \tfrac35 v_4)^2 + (v_2 + \tfrac35 v_3 - \tfrac35 v_4)^2 +
\tfrac{7}{25} v_3^2 + \tfrac{7}{25} v_4^2$. Check: the two middle squares give
$v_3^2$ coefficient $\tfrac{9}{25} + \tfrac{9}{25} = \tfrac{18}{25}$, plus $\tfrac{7}{25}$ is
1; the $v_3 v_4$ terms are $+\tfrac{18}{25}$ and $-\tfrac{18}{25}$ and cancel; the cross
terms are $2 \cdot \tfrac35 v_1 v_3$, $2 \cdot \tfrac35 v_1 v_4$, $2 \cdot \tfrac35 v_2 v_3$
and $-2 \cdot \tfrac35 v_2 v_4$, as the entries require.
\end{proof}

By Theorem~\ref{thm:chsh-fine} no finite factorizable model has $S > 2$, so this point is not a finite
mixture of plans, while Theorem~\ref{thm:chsh-ell-convex} puts every such mixture inside. The completion
set then strictly contains the finite classical set, and the PR box lies outside it. It sits
inside the no-signalling set, for the plain reason that every correlator in it is between $-1$
and $1$ (the vector that is 1 at $A_x$ and $\mp 1$ at $B_y$ gives $2 \mp 2E_{xy} \ge 0$), and any
four such numbers come from the no-signalling box that answers $(a, b)$ at settings $(x, y)$ with
probability $(1 + ab\,E_{xy})/4$, each player's own answer alone a fair coin.

The elliptope gate is an integer decider with a soundness proof. Its
completion entries arrive as bucket pairs, and an empty bucket is read as $0/1$,
so unlike Section~\ref{sub:chsh-integer} it does not require every pair to be sampled.

\begin{definition}[The two branches]
The positive-minor branch: with the completion supplied as buckets and all
entries scaled to integers by a positive common factor $K$, the four leading minors of
the $4 \times 4$ block on $A_0, A_1, B_0, B_1$ are strictly positive. The factorization branch: with the completion and a rational certificate supplied (six subdiagonal entries over
a common positive denominator, four pivots $d_k \ge 0$ over another), the ten entry equations
of $M = L D L^{\mathsf T}$ hold in cross-multiplied integer arithmetic. The full gate is their disjunction.
\end{definition}

\begin{theorem}[Soundness of the gate]
If the full gate returns yes on a counter record, any
completion buckets and any certificate, then the four correlators of the record are
elliptope realizable. On the PR box
tally the gate returns no for every choice of the remaining arguments.
\end{theorem}

\begin{proof}
In both branches it is enough to show the $4 \times 4$ block on $A_0, A_1, B_0, B_1$ gives a
nonnegative form, because the $I$ row and column are 1 on the diagonal and 0 elsewhere. The
first branch: dividing by the positive factor $K$ (a $k \times k$ minor scales by $K^k$) gives four
positive leading minors of the real block, and Lemma~\ref{lem:chsh-sylvester} applies. The
second: cross-multiplying by positive denominators changes no equation, so $M = LDL^{\mathsf T}$
holds over the rationals with $L$ unit lower triangular, and the form is
$v^{\mathsf T} L D L^{\mathsf T} v = \sum_k d_k P_k^2$ with $P = L^{\mathsf T} v$ and every
$d_k \ge 0$. The PR statement is soundness followed by Corollary~\ref{cor:chsh-ell-pr}.
\end{proof}

Examples by direct computation: the all-ones plan passes through the factorization branch
at $x=y=1$ with pivots $(1,0,0,0)$; the classical mixture $(1,0,1,0)$ passes through the same
branch at $x=1$, $y=0$ with pivots $(1,0,0,1)$; the $12/5$ point passes through the
positive-minor branch at $x=y=0$; and $(3/5,4/5,4/5,-3/5)$ passes through the factorization
branch at $x=y=0$ with pivots $(1,1,0,0)$, since $B_0 = \tfrac35 A_0 + \tfrac45 A_1$ and
$B_1 = \tfrac45 A_0 - \tfrac35 A_1$ reproduce every entry.

That's all the gate is: no property of the small machine is built from it.

\subsection{The four-vector construction}


Put two equal-length pencils down with their ends together, and draw a line along
one of them. Keep that pencil still and turn the other. Use the square corner of
a sheet of paper to bring the moving pencil's tip straight across to the line,
and mark where it lands. The mark moves from a full pencil-length forward,
through the joined ends when the pencils make a square corner, to a full
pencil-length backward. Call those readings 1, 0 and $-1$. That is the signed
shadow I want. Between those positions the reading is the cosine of the angle.
A dot product of unit vectors measures that shadow.

Now ask for four directions, two labelled Alice and two Bob, with a prescribed
shadow for each Alice--Bob pair. Try this request: Alice's first points exactly
along both of Bob's; Alice's second points along Bob's first and exactly against
his second. Each request alone is easy. Put them together and your hands cannot do
it: Bob's pencils would have to point both the same way and opposite ways. Four
numbers in range need not describe four directions that can coexist.

Write those four readings as $x=E_{00}$, $y=E_{01}$, $z=E_{10}$ and $w=E_{11}$.
I haven't specified the shadow between Alice's own two directions, or between
Bob's own two. Those are the two blanks, $t$ and $u$, in the following table.
Can they be filled so that the whole table comes from directions? The pencils let
you ask it in a plane; the construction will allow four coordinates. A completion is
a positive-semidefinite matrix
\[
G=\begin{pmatrix}
1&t&x&y\\t&1&z&w\\x&z&1&u\\y&w&u&1
\end{pmatrix}.
\]
The extra row and column labelled $I$ in Section~\ref{sub:chsh-elliptope} have diagonal 1 and all other entries
zero, so they do not change this condition.

Why positivity? Lay arrows end to end, allowing a backward arrow when its weight
is negative. Square the length from the start to the finish. Expanding that square
adds up the arrows' pairwise shadows, with their weights. It cannot be negative.
A proposed shadow table that gives a negative answer for some choice of weights
cannot come from arrows. That is the PSD test. A Gram matrix is simply the table
of all those pairwise dot products.

The converse has to build the arrows. Pick the first direction as a baseline.
Give every other arrow the component along it that its shadow demands, then
build the remaining components perpendicular to it. What remains to be fitted is
a smaller shadow table. Here is that construction, including its boundary. If its first diagonal entry is zero, positivity of
$h_{11}a^2+2h_{1j}ab+h_{jj}b^2$ for every $a,b$ forces $h_{1j}=0$:
with $b=1$, a nonzero $h_{1j}$ would let the linear term go negative without bound.
Delete that zero row and column, construct vectors for the smaller matrix, and put back
the zero vector. If $h_{11}>0$, set
\[
K_{ij}=h_{ij}-\frac{h_{1i}h_{1j}}{h_{11}}\quad(i,j>1).
\]
In the quadratic form of $h$, take the first coordinate to be
$-\sum_{j>1}h_{1j}v_j/h_{11}$. The remaining form is $v^{\mathsf T}Kv$, so $K$ is PSD.
Repeat the construction on this smaller matrix. Each repetition removes a row
and column, so it ends at a single nonnegative entry, represented by its square
root, or at the empty table. Working back up supplies vectors $r_i$ with
$r_i\cdot r_j=K_{ij}$. Then
$u_1=(\sqrt{h_{11}},0)$ and $u_i=(h_{1i}/\sqrt{h_{11}},r_i)$ have Gram matrix $h$.
The dimension is at most the number of rows. Applied to $G$, this gives four unit
vectors $a_0,a_1,b_0,b_1$ in $\mathbb R^4$, with the required cross inner products.

Three dimensions are always enough. The table fixes the cross entries and leaves two entries of
$G$ open, $a_0\cdot a_1$ and $b_0\cdot b_1$. Hold the first and let the second be $t$. The values
of $t$ that keep $G$ PSD form an interval: if two values do, so does every value between, because
the quadratic form at a mixture is the same mixture of the two forms; the set is closed, because a
form that is nonnegative at every value of a convergent sequence is nonnegative at its limit; and
it is bounded, because a PSD matrix with ones on its diagonal has every entry between $-1$ and 1.
At an end of that interval $G$ is singular: if it were positive definite there, it would stay
positive definite under every small enough change of $t$, on both sides. Take $G$ at an end. A
singular PSD $G$ has a nonzero $c$ with $c^{\mathsf T}Gc=0$, and $c^{\mathsf T}Gc$ is the squared
length of $\sum_i c_iv_i$ for the four vectors $v_i$, so that sum is zero and the four vectors lie
in a space of at most three dimensions. Coordinates in an orthonormal basis of that space keep every
inner product, so $a_0,a_1,b_0,b_1$ can be taken in $\mathbb R^3$.

The next question is harder: how do directions become a prescription for a
quantum measurement? Start with two signed cards in a row. One move swaps them;
another changes the sign of the card in the right-hand place. Start at $(+,+)$.
Swap then change the right sign: $(+,-)$. Start again, change the right sign then
swap: $(-,+)$. Every sign is reversed. Try unequal numbers on the cards: the two
orders still give negatives of each other. ``Anticommute'' means exactly that
relation between the two moves. Do either move twice and everything is back.

Write the swap as $X$, the sign change as $Z$, and doing nothing as $I$:
\[
X=\begin{pmatrix}0&1\\1&0\end{pmatrix},\qquad
Z=\begin{pmatrix}1&0\\0&-1\end{pmatrix},\qquad I=\begin{pmatrix}1&0\\0&1\end{pmatrix}.
\]
For the larger construction, label four places by the two-bit words 00 through
11 and put a signed number at each. An $X$ in one bit-position swaps each pair of
places differing at that bit; a $Z$ reverses the signs at places whose bit there is 1.
The two positions give the following three moves, written as real $4\times4$
matrices:
\[
\Gamma_1=X\otimes I,\quad
\Gamma_2=Z\otimes X,\quad
\Gamma_3=Z\otimes Z.
\]
The tensor product has entries $(A\otimes B)_{(i,k),(j,l)}=A_{ij}B_{kl}$.
Each $\Gamma_i$ is symmetric and squares to the identity. For $i<j$, the first
position where the two products differ has $X$ against $Z$, and $XZ=-ZX$; all
other positions commute. So $\Gamma_i\Gamma_j=-\Gamma_j\Gamma_i$.
For a unit vector $a$, the symmetric matrix $A(a)=\sum_i a_i\Gamma_i$ then
satisfies $A(a)^2=\sum_i a_i^2 I=I$. Its two-outcome measurement has projectors
$(I+A(a))/2$ and $(I-A(a))/2$: they are symmetric, square to themselves, are
orthogonal, and sum to the identity.

Take
$\lvert\Phi_4\rangle=4^{-1/2}\sum_{k=1}^{4}\lvert k\rangle\otimes\lvert k\rangle$.
Expanding this sum gives
\[
\langle\Phi_4\rvert A\otimes B\lvert\Phi_4\rangle
=\frac14\sum_{i,j}A_{ij}B_{ij}=\frac14\operatorname{tr}(A^{\mathsf T}B).
\]
For $i\ne j$, cyclicity of the trace and anticommutation give
$\operatorname{tr}(\Gamma_i\Gamma_j)=-\operatorname{tr}(\Gamma_i\Gamma_j)=0$;
for $i=j$ the trace is 4. So the correlator of $A(a)$ and $A(b)$ is
$a\cdot b$. Give Alice $A(a_0),A(a_1)$ and Bob $A(b_0),A(b_1)$.
Their strategy realizes $x,y,z,w$. All three $\Gamma_i$ have trace zero, so the
one-party expectations vanish too. The same real matrices and state are also a
complex strategy.

With complex amplitudes two levels on each side are enough. Add
$Y=\begin{pmatrix}0&-i\\ i&0\end{pmatrix}$ to $X$ and $Z$: the three are self-adjoint, square
to $I$ and anticommute in pairs, and they have trace zero. Give Alice
$a_1X+a_2Y+a_3Z$ for her vector $a$, and Bob $b_1X-b_2Y+b_3Z$ for his vector $b$. The expansion
above holds for complex entries too, with $\Phi_2$ and a factor $\frac12$, and the transpose of
Alice's matrix is $a_1X-a_2Y+a_3Z$, since $Y$'s transpose is $-Y$. So the correlator is
$\frac12\operatorname{tr}\bigl((a_1X-a_2Y+a_3Z)(b_1X-b_2Y+b_3Z)\bigr)=a\cdot b$.

What would that mean at a bench? Prepare the stated shared state, choose one of
the two specified measurements on each side, and write down the two signs returned.
For each pair of settings, count agreements as $+1$ and disagreements as $-1$.
Their average is the correlator. The model predicts the four averages $x,y,z,w$;
a finite tally fluctuates around its expectation. The signed-card moves let you
check the measurement algebra by hand. Realizing that state and those measurements
in a device, and comparing its tallies with the predictions, is the physical experiment.

Conversely, in a normalized quantum state $\psi$ with self-adjoint measurements
$A_x^2=B_y^2=I$ on separate systems, put
$a_x=(A_x\otimes I)\psi$ and $b_y=(I\otimes B_y)\psi$. They are unit vectors and
$\operatorname{Re}\langle a_x,b_y\rangle=E_{xy}$. Replace each complex coordinate by
its real and imaginary parts. This preserves those real inner products, and their
Gram matrix is a completion. A mixed state gives the same construction by taking
the direct sum of the vectors of its pure components, weighted by the square roots
of their probabilities. So the completion condition and finite-dimensional quantum
realizability are equivalent.

The necessary direction does not require finite-dimensional vectors. If a positive
normalized expectation $\omega$ on an algebra has self-adjoint $A_x,B_y$, with
$\omega(A_x^2),\omega(B_y^2)\le1$, then the real symmetric matrix with entries $\operatorname{Re}\,\omega(X_iX_j)$ is PSD: its quadratic form is $\omega(C^*C)\ge0$ for the real linear combination $C$
of the four observables. Increasing each diagonal entry to 1 preserves PSD. When
Alice's observables commute with Bob's, its cross entries are precisely the real
correlators. This supplies the same $G$. Including the unit as a fifth observable
also gives the level-1 moment matrix with its marginals. The sufficient direction
already used only two four-level systems.

\subsection{The arcsine boundary, in both directions}


Keep Bob's two directions hinged at one end. Alice's first direction has to make
a specified angle with each. How far can Bob's hinge open while that is possible?
Hold Alice's direction upright and sweep each of Bob's around it at its required
angle. Point them towards the same side to close the hinge as far as it will go;
point them towards opposite sides to open it as far as it will go. There is a
whole interval of openings between those limits. Alice's second direction asks
for another interval.

Draw both intervals on one strip marked from 0 to 180 degrees. If they miss, no
hinge setting satisfies both. If they overlap, choose an opening in the overlap
and construct each Alice direction at that opening. The algebra below permits
each its own perpendicular component in four dimensions. The pencils show why
the shared opening is the question; they need not fit the entire construction
on a tabletop. In the impossible request above, one interval is just 0 degrees
and the other just 180. There is nowhere to put the hinge.

Assume $x,y,z,w\in[-1,1]$, and put
$s_x=\arcsin x$, with analogous notation for the other three entries. A unit vector
making angles $\alpha,\beta\in[0,\pi]$ with two unit vectors can exist at an angle
$\delta$ between those two exactly when
\[
|\alpha-\beta|\le\delta\le
\min(\alpha+\beta,2\pi-\alpha-\beta).
\]
To see this, decompose the two vectors along and perpendicular to the first. Their
inner product ranges from $\cos\alpha\cos\beta-\sin\alpha\sin\beta$ to
$\cos\alpha\cos\beta+\sin\alpha\sin\beta$, by Cauchy--Schwarz; every value in that
interval is attained by choosing the angle between their perpendicular components.
Taking arccosines gives the displayed interval, including zero perpendicular components.

Apply this to Alice's first vector and Bob's two. Since
$\alpha=\pi/2-s_x$ and $\beta=\pi/2-s_y$, Bob's angle must lie in
\[
I_0=[\,|s_x-s_y|,\ \pi-|s_x+s_y|\,].
\]
Alice's second vector gives
$I_1=[\,|s_z-s_w|,\ \pi-|s_z+s_w|\,]$.
Each interval is nonempty because $|a-b|+|a+b|=2\max(|a|,|b|)\le\pi$
for $a,b\in[-\pi/2,\pi/2]$. They intersect exactly when
\[
|s_x-s_y|+|s_z+s_w|\le\pi,\qquad
|s_z-s_w|+|s_x+s_y|\le\pi.
\]
Using $|a|+|b|=\max(|a+b|,|a-b|)$, these are the four inequalities
\[
\begin{aligned}
|s_x+s_y+s_z-s_w|&\le\pi,&|s_x+s_y-s_z+s_w|&\le\pi,\\
|s_x-s_y+s_z+s_w|&\le\pi,&|-s_x+s_y+s_z+s_w|&\le\pi.
\end{aligned}
\]
So a completion implies all four.

For the converse, choose any $\delta\in I_0\cap I_1$ and set
$b_0=e_1$, $b_1=\cos\delta\,e_1+\sin\delta\,e_2$.
If $0<\delta<\pi$, put
\[
\begin{aligned}
c_0&=\frac{y-x\cos\delta}{\sin\delta},&
c_1&=\frac{w-z\cos\delta}{\sin\delta},\\
a_0&=xe_1+c_0e_2+\sqrt{1-x^2-c_0^2}\,e_3,\\
a_1&=ze_1+c_1e_2+\sqrt{1-z^2-c_1^2}\,e_4.
\end{aligned}
\]
The interval conditions say exactly that the two radicands are nonnegative:
after multiplication by $\sin^2\delta$, the first is
$1+2xy\cos\delta-x^2-y^2-\cos^2\delta\ge0$, and likewise for $z,w$.
These four vectors have unit length and the prescribed cross inner products.
If $\delta=0$, the intervals force $x=y$ and $z=w$; take $b_0=b_1=e_1$,
$a_0=xe_1+\sqrt{1-x^2}e_3$, $a_1=ze_1+\sqrt{1-z^2}e_4$.
If $\delta=\pi$, they force $y=-x$ and $w=-z$; use the same Alice vectors and
$b_1=-b_0$. Their Gram matrix supplies the completion in every case.
Together with the preceding construction, this proves the stated quantum characterization.

\subsection{The finite local model, in both directions}


Make sixteen cards, each with four boxes labelled $A_0,A_1,B_0,B_1$, and put every
possible arrangement of plus and minus signs on them, one arrangement per card.
Draw a card before choosing which setting to ask each player about. Alice reads
her chosen box; Bob reads his. Neither needs the other's setting. How often each
card is drawn is the only choice left in this model.

Try to get perfect agreement for the first three pairs and perfect disagreement
for the fourth. Any card that ever gets drawn would need $A_0=B_0$, $A_0=B_1$
and $A_1=B_0$, which makes $A_1=B_1$ too. Changing how often you draw the cards
cannot fix it. For less extreme requests the question is quantitative: can you
assign nonnegative weights to the cards, adding to one, that give all four
requested averages? The following construction supplies those weights exactly
when all eight CHSH inequalities hold and the correlators are in $[-1,1]$.

There is a useful way to assemble the deck. First allow only Bob's first question
and make a distribution of the three answers $A_0,A_1,B_0$. Then do the same for
his second question. Both smaller decks must draw Alice's answer pairs in the
same proportions. The adjustable number $t$ is Alice's agreement-minus-disagreement
average across her two boxes. Choose it in both intervals
\[
J_0=[\,-1+|x+z|,\ 1-|x-z|\,],\qquad
J_1=[\,-1+|y+w|,\ 1-|y-w|\,].
\]
Each interval is nonempty. They overlap exactly when
$|x+z|+|y-w|\le2$ and $|y+w|+|x-z|\le2$, which expand to the eight
CHSH inequalities. For $a_0,a_1,b\in\{-1,1\}$ set
\[
\begin{aligned}
P_0(a_0,a_1,b)&=\frac{1+t a_0a_1+x a_0b+z a_1b}{8},\\
P_1(a_0,a_1,b)&=\frac{1+t a_0a_1+y a_0b+w a_1b}{8}.
\end{aligned}
\]
When $a_0=a_1$, nonnegativity of $P_0$ is $1+t\ge|x+z|$;
when $a_0=-a_1$, it is $1-t\ge|x-z|$. These are precisely $t\in J_0$.
The same check for $P_1$ uses $J_1$. Summing signs shows that both distributions
have mass 1 and the desired pair correlators. Their common Alice marginal is
$P_A(a_0,a_1)=(1+t a_0a_1)/4$.

To join the two smaller decks, first draw Alice's pair with its common probability.
Among the first deck's cards bearing that pair, draw Bob's first answer; among
the second deck's cards bearing the same pair, draw his second. Those conditional
draws give a four-answer card. A pair that never occurs needs no card. Written
as probabilities, this procedure is
\[
P(a_0,a_1,b_0,b_1)=
\begin{cases}
P_0(a_0,a_1,b_0)P_1(a_0,a_1,b_1)/P_A(a_0,a_1),&P_A(a_0,a_1)>0,\\
0,&P_A(a_0,a_1)=0.
\end{cases}
\]
If $P_A=0$, nonnegativity forces both corresponding triple distributions to be zero.
Otherwise summing over either Bob answer cancels the denominator. So $P$ has the
two triple distributions as marginals and total mass 1. Sample one of its sixteen
answer quadruples and let each party answer with its assigned sign. This is a finite
mixture of deterministic local plans with exactly the four required correlators.

The other direction is one line per card. On a single card,
$A_0(B_0+B_1)+A_1(B_0-B_1)$ is $\pm2$, because one bracket is 0 and the other is $\pm2$.
Each of the other seven CHSH expressions is this one with the boxes relabelled or a sign
flipped, so it is $\pm2$ on every card too, and any weighting of the cards keeps all eight
averages between $-2$ and 2 (Theorem~\ref{thm:chsh-fine} is the first of them). So a
deck exists exactly when the correlators are in $[-1,1]$ and the eight inequalities hold.

\subsection{Scope of the constructions}


The constructions above pin down the correlators. What follows is where they stop: they
aren't a completeness claim for the integer check, they say nothing about full behaviours
with marginals, and they say nothing about data from a lab.

The arcsine boundary. For two settings and two outcomes per side, Tsirelson, Landau (1988) and
Masanes (2003) give theorems characterizing the correlators reachable by quantum measurements:
$|\arcsin E_{00} + \arcsin E_{01} + \arcsin E_{10} - \arcsin E_{11}| \le \pi$ and its sign
variants. The characterization holds as an if and only if with the
completion set and, through Tsirelson's representation, with the quantum
set. The point $(3/5, 4/5, 4/5, -3/5)$ is on the boundary:
$\arcsin \tfrac35 + \arcsin \tfrac45 = \tfrac{\pi}{2}$ by the 3-4-5 triangle, so the sum is $\pi$; Tsirelson's point is on it too.
The necessary direction holds in any dimension, for any expectation on an algebra of
observables in which Alice's commute with Bob's.

Equality with the quantum set. That the completion set of Section~\ref{sub:chsh-elliptope} is exactly the set of
quantum correlators is Tsirelson's representation theorem, and it is proved in finite dimension:
for both real and complex amplitudes, with two four-level systems enough for every point, and two
two-level ones with complex amplitudes. That PSD of the level-1 moment matrix is necessary for
quantum behaviour is proved for strategies in finite dimension, with real amplitudes
and with complex ones, and in any dimension, marginals included, for any expectation on
an algebra of observables.

Achievability. Theorem~\ref{thm:chsh-tight} is a point of $\mathbb{R}^4$. A state and
measurements that reach it are explicit: the two-qubit state
$(|00\rangle + |11\rangle)/\sqrt 2$ with $A_0 = Z$, $A_1 = X$, $B_0 = (Z + X)/\sqrt 2$ and
$B_1 = (Z - X)/\sqrt 2$ has correlators $1/\sqrt 2$, $1/\sqrt 2$, $1/\sqrt 2$, $-1/\sqrt 2$ and
$S = 2\sqrt 2$, and every strategy in finite dimension has
$S^2 \le 8$, with real amplitudes and with complex ones. Both are arithmetic about the model.
How close a real laboratory comes, and what that costs in loopholes closed, is the story told in
Appendix~\ref{sec:chsh}.

Completeness of the checks. The 1+AB cascade and the elliptope gate are proved sound. The cascade has no proved converse; the gate's factorization branch accepts every rational PSD completion, by the row-peeling of Section~\ref{sub:four-vectors} with each pivot kept. The level-1 integer check is the exception: on the small machine the same
seven integer facts hold exactly when every pair was sampled and the pinned matrix is positive
semidefinite (Theorem~\ref{thm:chsh-exact}), so that check is
complete for what it tests. The 1+AB cascade is strict and misses boundary points.
The two 1+AB checks don't sit inside one another. The full check accepts the $12/5$ tally above, which the zero-moment check refuses, and it refuses the tally with $E_{00}=1$ and the other three correlators 0, where $H_{11}=0$, which the zero-moment check accepts.

Finite levels. For full behaviours, marginals included, the NPA levels are relaxations.
Chaturvedi (arXiv 2607.14569, 2026) argues, in a preprint, that no finite NPA level characterizes the complete
quantum set even in this simplest Bell scenario. I haven't gone through that proof, and nothing
here rests on it.

Statistics and physics. Nothing here says how many rounds make a tally informative, and
nothing says the counters came from honest play; a program can INC whatever counts it likes.

\subsection{Prior work}


Tsirelson's 1980 paper gives the bound $2\sqrt 2$ for quantum correlations, and his theorem
says that quantum two-outcome correlators are inner products of real unit vectors. Landau (1988)
and Masanes (2003, arXiv quant-ph/0309137) each give a condition that, by their proofs,
characterizes the quantum correlators exactly for two settings and two outcomes. Wehner (2006,
arXiv quant-ph/0510076) gets Tsirelson bounds, CHSH among them, from a semidefinite program on
the $4 \times 4$ Gram matrix of those vectors. Navascués,
Pironio and Acín (2007, arXiv quant-ph/0607119; 2008, arXiv 0803.4290) built the hierarchy of
moment-matrix tests; their 2008 paper applies a Schur complement to the level-1 matrix of this
scenario, recovers the arcsine condition, and discusses the 1+AB level. I went back to each of
those papers for these points, and that is what I found in them. The results are theirs. Two of
them are proved again here, in finite dimension: Tsirelson's representation
and the arcsine characterization; the rest of
their proofs are not redone. Echenim and Mhalla give a detailed derivation of Tsirelson's
bound (arXiv 2306.12535), and exact rational
sum-of-squares certificates go back at least to Peyrl and Parrilo (2008). So the
moment-matrix route to $2\sqrt 2$ is standard. What is specific here is narrower: an integer
certificate computed from outcome counts, and the argument from ``the counts pass'' to PSD
of the pinned matrix and on to $S^2 \le 8$. I did not find that chain elsewhere, and I make no
claim beyond it.
On the host machine, the same chain is earned: a raised flag from a clean start
stands on a passing check of the packed counts, and the check passes exactly when the pinned matrix
is positive semidefinite and every pair was sampled (Appendix~\ref{sec:chsh},
Theorems~\ref{thm:chsh-exact} and~\ref{thm:chsh-earned}). The converse direction of the integer check,
integers from PSD, is proved there too, and there too I make no claim beyond the chain.

%===========================================================================


\section{Further reading}

The works below were read for the book and are not named in its finished text or in this volume. They are kept here with their annotations, so that nothing read for the book is lost.

\subsection{Structure, information and knowledge}



\noindent Hans Albert. Wikipedia. \href{https://en.wikipedia.org/wiki/Hans_Albert}{Article}. German philosopher (1921--2023), with Popper a leading critical rationalist; Traktat über kritische Vernunft (1968).
\par\medskip

\noindent A.~J. Ayer (1956). The Problem of Knowledge. Macmillan, London, p.~34. The other definition Gettier targets.
\par\medskip

\noindent Max Black. Wikipedia. \href{https://en.wikipedia.org/wiki/Max_Black}{Article}. The Russian-born British-American philosopher, at Cornell from 1946, and his 1952 paper on two spheres.
\par\medskip

\noindent Roberto Bonola (1912). Non-Euclidean Geometry: A Critical and Historical Study of Its Development. Open Court; Dover reprint, 1955. The attempts on the parallel postulate from Proclus through the Islamic geometers and Saccheri, and Farkas Bolyai's letter to his son.
\par\medskip

\noindent G. K. Chesterton (1929). The Thing. The fence passage is quoted in the American Chesterton Society's \href{https://www.chesterton.org/lecture-57/}{Lecture 57: The Thing}.
\par\medskip

\noindent G\"osta Grahne (1998). Updates and counterfactuals. Journal of Logic and Computation 8(1), pp. 87--117. \href{https://doi.org/10.1093/logcom/8.1.87}{DOI: 10.1093/logcom/8.1.87}. Revision of a static world versus update of a changing one.
\par\medskip

\noindent Jonathan Jenkins Ichikawa and Matthias Steup (2001, substantive revision 2026). The Analysis of Knowledge. Stanford Encyclopedia of Philosophy. \href{https://plato.stanford.edu/entries/knowledge-analysis/}{Entry}. Dharmottara's mirage and Peter of Mantua's running Plato.
\par\medskip

\noindent Rosalie Iemhoff (2008, substantive revision 2019). Intuitionism in the Philosophy of Mathematics. Stanford Encyclopedia of Philosophy. \href{https://plato.stanford.edu/entries/intuitionism/}{Entry}. Kreisel's creating-subject axioms CS1--CS3 (1967), CS2 being that the creating subject never forgets what it has proven, and Troelstra's strengthening (1969).
\par\medskip

\noindent Petrus Mantuanus (d. 1400). Logica. Medieval manuscripts in Oxford libraries, Bodleian Library. \href{https://medieval.bodleian.ox.ac.uk/catalog/work_3607}{Catalogue entry}. Peter of Mantua as a logician of the late fourteenth century.
\par\medskip

\noindent Per Martin-Löf (1996). On the Meanings of the Logical Constants and the Justifications of the Logical Laws. Nordic Journal of Philosophical Logic 1(1), pp. 11--60. Lectures given in Siena in 1983. The passage on proof as an act is on p.~28. \href{https://archive-pml.github.io/martin-lof/pdfs/Meanings-of-the-Logical-Constants-1983.pdf}{Lecture text}.
\par\medskip

\noindent Meta Platforms (2021). Introducing Meta: A Social Technology Company. Company announcement, 28 October 2021. \href{https://about.fb.com/news/2021/10/facebook-company-is-now-meta/}{Announcement}.
\par\medskip

\noindent Benjamin Morison (2014, revised 2026). Sextus Empiricus. Stanford Encyclopedia of Philosophy. \href{https://plato.stanford.edu/entries/sextus-empiricus/}{Entry}. The source for his dates, his medicine and his skepticism.
\par\medskip

\noindent J. J. O'Connor and E. F. Robertson (n.d.). Theaetetus of Athens. MacTutor History of Mathematics. \href{https://mathshistory.st-andrews.ac.uk/Biographies/Theaetetus/}{Entry}.
\par\medskip

\noindent J. J. O'Connor and E. F. Robertson (n.d.). Paul Richard Halmos. MacTutor History of Mathematics. \href{https://mathshistory.st-andrews.ac.uk/Biographies/Halmos/}{Biography}.
\par\medskip

\noindent Benoît Rittaud and Albrecht Heeffer (2014). The Pigeonhole Principle, Two Centuries before Dirichlet. The Mathematical Intelligencer 36(2), pp. 27--29. Jean Leurechon's use of the principle in 1622.
\par\medskip

\noindent G\"oran Sundholm (1998). Proofs as Acts and Proofs as Objects: Some questions for Dag Prawitz. Theoria 64(2--3), pp. 187--216. \href{https://doi.org/10.1111/j.1755-2567.1998.tb00187.x}{DOI: 10.1111/j.1755-2567.1998.tb00187.x}.
\par\medskip

\noindent G\"oran Sundholm (2014). Constructive Recursive Functions, Church's Thesis, and Brouwer's Theory of the Creating Subject: Afterthoughts on a Parisian Joint Session. In Constructivity and Computability in Historical and Philosophical Perspective, Logic, Epistemology, and the Unity of Science, Springer, pp. 1--35. \href{https://doi.org/10.1007/978-94-017-9217-2_1}{DOI: 10.1007/978-94-017-9217-2\_1}.
\par\medskip

\noindent Timothy Yenter and Ezio Vailati (2003, revised 2024). Samuel Clarke. Stanford Encyclopedia of Philosophy. \href{https://plato.stanford.edu/entries/clarke/}{Entry}. Clarke as the British philosopher and theologian who defended Newton's natural philosophy in his letters with Leibniz.
\par\medskip

\subsection{Cost, records and thermodynamics}



\noindent J\'anos Acz\'el, Bruno Forte and C. T. Ng (1974). Why the Shannon and Hartley entropies are ``natural.'' Advances in Applied Probability 6(1), pp. 131--146. \href{https://doi.org/10.2307/1426210}{DOI: 10.2307/1426210}. Ng's 1971 Waterloo doctorate under Acz\'el: \href{https://mathgenealogy.org/id.php?id=102040}{Mathematics Genealogy Project}.
\par\medskip

\noindent Roberto M. Amadio, Nicolas Ayache, François Bobot, Jaap P. Boender, Brian Campbell, Ilias Garnier, Antoine Madet, James McKinna, Dominic P. Mulligan, Mauro Piccolo, Randy Pollack, Yann Régis-Gianas, Claudio Sacerdoti Coen, Ian Stark and Paolo Tranquilli (2014). Certified Complexity (CerCo). Foundational and Practical Aspects of Resource Analysis, FOPARA 2013, LNCS 8552, pp. 1--18. \href{https://doi.org/10.1007/978-3-319-12466-7_1}{DOI: 10.1007/978-3-319-12466-7\_1}.
\par\medskip

\noindent Walter J. Boyne. Wikipedia entry: United States Air Force officer, aviation historian, director of the National Air and Space Museum (acting from 1982, director 1983--1986). \href{https://en.wikipedia.org/wiki/Walter_J._Boyne}{Entry}.
\par\medskip

\noindent Léon Brillouin (1951). Maxwell's Demon Cannot Operate: Information and Entropy. I. Journal of Applied Physics 22(3), pp. 334--337. \href{https://doi.org/10.1063/1.1699951}{DOI: 10.1063/1.1699951}.
\par\medskip

\noindent Arthur Charguéraud and François Pottier (2019). Verifying the Correctness and Amortized Complexity of a Union-Find Implementation in Separation Logic with Time Credits. Journal of Automated Reasoning 62, pp. 331--365; first published online in 2017. \href{https://doi.org/10.1007/s10817-017-9431-7}{DOI: 10.1007/s10817-017-9431-7}.
\par\medskip

\noindent Gavin E. Crooks (1999). Entropy production fluctuation theorem and the nonequilibrium work relation for free energy differences. Physical Review E 60(3), pp. 2721--2726. \href{https://doi.org/10.1103/PhysRevE.60.2721}{DOI: 10.1103/PhysRevE.60.2721}. A path's chance against its reverse's is $e$ to the dissipated work.
\par\medskip

\noindent Joe DeBlasio (2025). Discussion: removal of one-6962-log requirement. Chromium ct-policy mailing list, 17 November 2025. \href{https://groups.google.com/a/chromium.org/g/ct-policy/c/amrnTiDFdwM/m/mCyPDc4TAwAJ}{Post}.
\par\medskip

\noindent Jérémie Decouchant, David Kozhaya, Vincent Rahli and Jiangshan Yu (2022). Damysus: streamlined BFT consensus leveraging trusted components. Proceedings of the Seventeenth European Conference on Computer Systems (EuroSys '22), pp. 1--16. \href{https://doi.org/10.1145/3492321.3519568}{DOI: 10.1145/3492321.3519568}.
\par\medskip

\noindent Encyclopaedia Britannica. Donation of Constantine. \href{https://www.britannica.com/topic/Donation-of-Constantine}{Web page}. Composed in the eighth century; the consensus view puts it in the 750s or 760s.
\par\medskip

\noindent Encyclopaedia Britannica. Leo Szilard. \href{https://www.britannica.com/biography/Leo-Szilard}{Web page}. Born in Budapest in 1898; Berlin doctorate 1922 and the Institute of Theoretical Physics there.
\par\medskip

\noindent Ethereum Foundation (n.d.). The Merge. ethereum.org roadmap page, on the transition of 15 September 2022. \href{https://ethereum.org/roadmap/merge/}{Page}.
\par\medskip

\noindent Ethereum Foundation (n.d.). Timeline of all Ethereum forks (Tangerine Whistle, EIP-150, and the Berlin upgrade, block 12,244,000). ethereum.org. \href{https://ethereum.org/ethereum-forks/}{Page}.
\par\medskip

\noindent Ethereum Foundation. Devcon editions (Devcon 0, Berlin 2014; Devcon 1, London 2015; Devcon 2, Shanghai 2016). devcon.org. \href{https://devcon.org/past-events}{Page}.
\par\medskip

\noindent Ethereum Foundation. Ethereum execution specifications, Prague fork, src/ethereum/forks/prague/vm/instructions/storage.py and vm/gas.py, at commit ecb68f21a7d28db101c769878184f6facb319da1. \href{https://github.com/ethereum/execution-specs}{Repository}.
\par\medskip

\noindent Federal Aviation Administration (2008). Revisions to Cockpit Voice Recorder and Digital Flight Data Recorder Regulations. Final rule, Federal Register 73(46), 7 March 2008, pp. 12542--12573. The two-hour cockpit voice recorder.
\par\medskip

\noindent Federal Aviation Administration (2026). 25-Hour Cockpit Voice Recorder (CVR) Requirement, New Aircraft Production. Final rule, Federal Register 91(21), 2 February 2026, pp. 4447--4462. \href{https://www.govinfo.gov/content/pkg/FR-2026-02-02/html/2026-02110.htm}{Federal Register text}. The thirty-minute recorders of 1966, the two-hour rule, and the 25-hour standard of ICAO (2016) and EASA (from 2021).
\par\medskip

\noindent Richard FitzNigel (c. 1180). Dialogus de Scaccario (The Dialogue of the Exchequer). Text and translation ed. Charles Johnson, Nelson's Medieval Classics, 1950. The widths of the cuts on Exchequer tallies.
\par\medskip

\noindent Fox-IT (2012). Black Tulip: Report of the investigation into the DigiNotar Certificate Authority breach. Version 1.0. \href{https://roselabs.nl/files/audit_reports/Fox-IT_-_DigiNotar.pdf}{Report}.
\par\medskip

\noindent Tom Furtwangler (2017). The vaccine vial monitor: the world's smartest sticker. PATH, 6 November 2017. \href{https://www.path.org/articles/vaccine-vial-monitor-worlds-smartest-sticker/}{Article}. The stickers first sold in 1996, for oral polio vaccine.
\par\medskip

\noindent General Conference on Weights and Measures (2018). Resolution 1 of the 26th meeting: On the revision of the International System of Units (SI), effective 20 May 2019.
\par\medskip

\noindent C. H. Jenkinson (1910). Early wooden tallies relating to Surrey. Surrey Archaeological Collections 23, pp. 203--207. \href{https://archaeologydataservice.ac.uk/catalogue/adsdata/arch-379-1/dissemination/pdf/vol_23/surreyac023_203-207_jenkinson.pdf}{PDF}. Hazel; the payer kept the stock, the Exchequer the foil; abolished by statute in 1826; the burning that started the fire.
\par\medskip

\noindent Joint Committee for Guides in Metrology (2012). International vocabulary of metrology: basic and general concepts and associated terms (VIM), 3rd edition, JCGM 200:2012, entry 2.41, ``metrological traceability.'' The definition is on p.~29; ``does not ensure \ldots{} that there is an absence of mistakes'' is Note~5 to 2.41, p.~30. \href{https://www.bipm.org/documents/20126/2071204/JCGM_200_2012.pdf}{BIPM text}.
\par\medskip

\noindent Mikhail Kuter, Marina Gurskaya and Albert Gurskii (2025). Genoa 1340, Double Entry, and `Gold Coasting'. De Computis, Revista Española de Historia de la Contabilidad 22(1), pp. 1--33. \href{https://doi.org/10.26784/issn.1886-1881.22.1.11927}{DOI: 10.26784/issn.1886-1881.22.1.11927}.
\par\medskip

\noindent Klaus-J\"orn Lange, Pierre McKenzie and Alain Tapp (2000). Reversible Space Equals Deterministic Space. Journal of Computer and System Sciences 60(2), pp. 354--367. \href{https://doi.org/10.1006/jcss.1999.1672}{DOI: 10.1006/jcss.1999.1672}.
\par\medskip

\noindent Ben Laurie, Adam Langley and Emilia Kasper (2013). Certificate Transparency. RFC 6962. \href{https://www.rfc-editor.org/rfc/rfc6962.html}{RFC 6962}; \href{https://doi.org/10.17487/RFC6962}{DOI: 10.17487/RFC6962}.
\par\medskip

\noindent Ben Laurie, Eran Messeri and Rob Stradling (2021). Certificate Transparency Version 2.0. RFC 9162. \href{https://www.rfc-editor.org/rfc/rfc9162.html}{RFC 9162}; \href{https://doi.org/10.17487/RFC9162}{DOI: 10.17487/RFC9162}.
\par\medskip

\noindent A. Liapounoff (1907). Problème général de la stabilité du mouvement, traduit du russe par Édouard Davaux. Annales de la Faculté des Sciences de Toulouse, 2e série, 9, pp. 203--474. \href{http://www.numdam.org/item?id=AFST_1907_2_9__203_0}{Numdam}. The French translation the 1992 English one was made from.
\par\medskip

\noindent Sinisa Matetic, Mansoor Ahmed, Kari Kostiainen, Aritra Dhar, David Sommer, Arthur Gervais, Ari Juels and Srdjan Capkun (2017). ROTE: rollback protection for trusted execution. Proceedings of the 26th USENIX Security Symposium, pp. 1289--1306. \href{https://www.usenix.org/conference/usenixsecurity17/technical-sessions/presentation/matetic}{USENIX page}.
\par\medskip

\noindent Clyde Lee Miller (2009; revised 2021). Cusanus, Nicolaus [Nicolas of Cusa]. Stanford Encyclopedia of Philosophy. \href{https://plato.stanford.edu/entries/cusanus/}{SEP entry}.
\par\medskip

\noindent Y\^uto Murashita, Ken Funo and Masahito Ueda (2014). Nonequilibrium equalities in absolutely irreversible processes. Physical Review E 90, 042110. \href{https://doi.org/10.1103/PhysRevE.90.042110}{DOI: 10.1103/PhysRevE.90.042110}. The bound on average entropy production by $-\ln(1-\lambda)$, with free expansion as the first example.
\par\medskip

\noindent Satoshi Nakamoto (2008). Bitcoin: A Peer-to-Peer Electronic Cash System. \href{https://bitcoin.org/bitcoin.pdf}{Paper}.
\par\medskip

\noindent Devon O'Brien (2018). Certificate Transparency Enforcement in Google Chrome. Chromium ct-policy mailing list, 6 February 2018. \href{https://groups.google.com/a/chromium.org/g/ct-policy/c/wHILiYf31DE/m/iMFmpMEkAQAJ}{Announcement}.
\par\medskip

\noindent Garson O'Toole (2016). There Ain't No Such Thing as a Free Lunch, TANSTAAFL. Quote Investigator, 27 August 2016. \href{https://quoteinvestigator.com/2016/08/27/free-lunch/}{Article}.
\par\medskip

\noindent Pavel Pakshin, Elena Panteley and Natalia Pakshina (2017). The Legacy of A. M. Lyapunov to Systems and Control Theory: 110 years of the Toulouse edition. IFAC World Congress 2017, open-track proposal. Davaux's translation, reviewed and corrected by Lyapunov.
\par\medskip

\noindent Lídia del Rio, Johan Åberg, Renato Renner, Oscar Dahlsten and Vlatko Vedral (2011). The thermodynamic meaning of negative entropy. Nature 474, pp. 61--63. \href{https://doi.org/10.1038/nature10123}{DOI: 10.1038/nature10123}; \href{https://arxiv.org/abs/1009.1630}{arXiv: 1009.1630}.
\par\medskip

\noindent J\"urgen Schnakenberg (1976). Network theory of microscopic and macroscopic behavior of master equation systems. Reviews of Modern Physics 48, pp. 571--585. \href{https://doi.org/10.1103/RevModPhys.48.571}{DOI: 10.1103/RevModPhys.48.571}. The entropy a Markov chain hands its surroundings, as the net rate across each pair times the log of the forward-to-backward rate ratio, which at a steady state is the whole entropy production; the biased walk is its smallest case.
\par\medskip

\noindent Securities and Exchange Commission (2022). Electronic Recordkeeping Requirements for Broker-Dealers, Security-Based Swap Dealers, and Major Security-Based Swap Participants. Release No. 34-96034, 12 October 2022. \href{https://www.sec.gov/rules/final/2022/34-96034.pdf}{Final rule}. Adds an audit-trail alternative to the write-once requirement of Rule 17a-4(f) in place since 1997.
\par\medskip

\noindent Caroline Shenton (2012). The Day Parliament Burned Down. Oxford University Press. The tallies burned in the furnaces under the House of Lords, 16 October 1834.
\par\medskip

\noindent Michael Stock, Pauline Barat, Richard S. Davis, Alain Picard and Martin J. T. Milton (2015). Calibration campaign against the international prototype of the kilogram in anticipation of the redefinition of the kilogram part I: comparison of the international prototype with its official copies. Metrologia 52(2), pp. 310--316. \href{https://doi.org/10.1088/0026-1394/52/2/310}{DOI: 10.1088/0026-1394/52/2/310}.
\par\medskip

\noindent Wei Tang (2019). EIP-2200: Structured Definitions for Net Gas Metering. Ethereum Improvement Proposals. \href{https://eips.ethereum.org/EIPS/eip-2200}{eips.ethereum.org/EIPS/eip-2200}.
\par\medskip

\noindent This Day in Aviation (n.d.). 30 October 1935: Boeing Model 299 crash at Wright Field. \href{https://www.thisdayinaviation.com/30-october-1935/}{Page}.
\par\medskip

\noindent Kent Weare (2016). Ethereum Security Alert Issued, Ethereum Foundation Responds with ``From Shanghai, With Love''. InfoQ, 19 September 2016. \href{https://www.infoq.com/news/2016/09/Ethereum-DOS-Attack}{Article}.
\par\medskip

\noindent World Health Organization (2011). PQS performance specification: Vaccine Vial Monitor. WHO/PQS/E06/IN05.2.
\par\medskip

\subsection{Computation, self-reference and checking}



\noindent Keith Adams and Ole Agesen (2006). A Comparison of Software and Hardware Techniques for x86 Virtualization. Proceedings of the 12th International Conference on Architectural Support for Programming Languages and Operating Systems (ASPLOS XII), pp. 2--13. \href{https://doi.org/10.1145/1168857.1168860}{DOI: 10.1145/1168857.1168860}. Binary translation of instructions such as POPF, which on x86 runs in user mode without trapping. Intel's VT-x and AMD-V in 2005 and 2006: Wikipedia, \href{https://en.wikipedia.org/wiki/X86_virtualization}{x86 virtualization}.
\par\medskip

\noindent American Standards Association (1963). American Standard Code for Information Interchange, ASA X3.4-1963.
\par\medskip

\noindent Henk Barendregt. Wikipedia. \href{https://en.wikipedia.org/wiki/Henk_Barendregt}{Article}. A Dutch logician known for his work in lambda calculus and type theory.
\par\medskip

\noindent Gerd Behrmann, Ansgar Fehnker, Thomas Hune, Kim Larsen, Paul Pettersson, Judi Romijn and Frits Vaandrager (2001). Minimum-cost reachability for priced timed automata. Hybrid Systems: Computation and Control (HSCC 2001), Lecture Notes in Computer Science 2034, pp. 147--161.
\par\medskip

\noindent N. G. de Bruijn (1972). Lambda calculus notation with nameless dummies, a tool for automatic formula manipulation, with application to the Church--Rosser theorem. Indagationes Mathematicae (Proceedings) 75(5), pp. 381--392. \href{https://doi.org/10.1016/1385-7258(72)90034-0}{DOI: 10.1016/1385-7258(72)90034-0}. Variables written as numbers that count the binders back, now called de Bruijn indices.
\par\medskip

\noindent Stanley Burris and Marcel Jackson (2010, substantive revision 2026). George Boole. Stanford Encyclopedia of Philosophy. \href{https://plato.stanford.edu/entries/boole/}{Entry}.
\par\medskip

\noindent Stanley Burris and Javier Legris (2009, substantive revision 2021). The Algebra of Logic Tradition. Stanford Encyclopedia of Philosophy. \href{https://plato.stanford.edu/entries/algebra-logic-tradition/}{Entry}.
\par\medskip

\noindent G. K. Chesterton. Wikipedia. \href{https://en.wikipedia.org/wiki/G._K._Chesterton}{Article}. An English writer, referred to as the ``prince of paradox''.
\par\medskip

\noindent Donald W. Davies (2004). Corrections to Turing's Universal Computing Machine. In B. Jack Copeland (ed.), The Essential Turing, pp. 103--124. Oxford University Press. Davies at the National Physical Laboratory from 1947 and the word \emph{packet}: Wikipedia, \href{https://en.wikipedia.org/wiki/Donald_Davies}{Donald Davies}.
\par\medskip

\noindent Charles Dickens (1843). A Christmas Carol. Chapman and Hall, London. Stave Three, the game of Yes and No.
\par\medskip

\noindent William Henry Eccles and Frank Wilfred Jordan (1920). Improvements in ionic relays. British patent GB148582A, filed 21 June 1918, published 5 August 1920. \href{https://patents.google.com/patent/GB148582A/en}{Patent text}. The 1918 date in the text is the filing date. The circuit as the flip-flop, the basis for electronic memory: History of Information, \href{https://www.historyofinformation.com/detail.php?id=3606}{entry}; and Wikipedia, \href{https://en.wikipedia.org/wiki/Flip-flop_(electronics)}{Flip-flop (electronics)}, on the Eccles--Jordan trigger circuit as the first electronic latch, later the basic storage element of digital circuitry.
\par\medskip

\noindent László Fejes Tóth (1953). Lagerungen in der Ebene, auf der Kugel und im Raum. Berlin: Springer.
\par\medskip

\noindent László Fejes Tóth (1964). Regular Figures. Oxford: Pergamon Press.
\par\medskip

\noindent Jean-Christophe Filliâtre, Léon Gondelman and Andrei Paskevich (2014). The spirit of ghost code. Computer Aided Verification (CAV 2014), Lecture Notes in Computer Science 8559, pp. 1--16. \href{https://doi.org/10.1007/978-3-319-08867-9_1}{DOI: 10.1007/978-3-319-08867-9\_1}.
\par\medskip

\noindent Nicholas Griffin, ed. (2003). The Cambridge Companion to Bertrand Russell. Cambridge University Press. G. G. Berry as a junior librarian at the Bodleian, p. 63.
\par\medskip

\noindent Jean-Pierre Marquis (1996, substantive revision 2019). Category Theory. Stanford Encyclopedia of Philosophy. \href{https://plato.stanford.edu/entries/category-theory/}{Entry}. The borrowed words ``category'' and ``functor''.
\par\medskip

\noindent Ross M. McConnell, Kurt Mehlhorn, Stefan Näher and Pascal Schweitzer (2011). Certifying algorithms. Computer Science Review 5(2), pp. 119--161. \href{https://doi.org/10.1016/j.cosrev.2010.09.009}{DOI: 10.1016/j.cosrev.2010.09.009}.
\par\medskip

\noindent Jason Megill. The Lucas-Penrose Argument about Gödel's Theorem. Internet Encyclopedia of Philosophy. \href{https://iep.utm.edu/lp-argue/}{Article}.
\par\medskip

\noindent Glen Mével, Jacques-Henri Jourdan and François Pottier (2019). Time credits and time receipts in Iris. ESOP 2019, LNCS 11423, pp. 3--29. \href{https://doi.org/10.1007/978-3-030-17184-1_1}{DOI: 10.1007/978-3-030-17184-1\_1}.
\par\medskip

\noindent Jeff Miller. Earliest Uses of Symbols of Set Theory and Logic. MacTutor History of Mathematics. \href{https://mathshistory.st-andrews.ac.uk/Miller/mathsym/set/}{Page}. Peano's epsilon for \emph{est}, 1889.
\par\medskip

\noindent J. J. O'Connor and E. F. Robertson (n.d.). The four colour theorem. MacTutor History of Mathematics. \href{https://mathshistory.st-andrews.ac.uk/HistTopics/The_four_colour_theorem/}{Entry}. (De Morgan's letter to Hamilton of 23 October 1852 is quoted from here.)
\par\medskip

\noindent Panu Raatikainen (2013, revised 2025). Gödel's Incompleteness Theorems. Stanford Encyclopedia of Philosophy. \href{https://plato.stanford.edu/entries/goedel-incompleteness/}{Entry}. Source for the cautions about what the Gödel sentence says.
\par\medskip

\noindent Erich Reck (2008, substantive revision 2025). Dedekind's Contributions to the Foundations of Mathematics. Stanford Encyclopedia of Philosophy. \href{https://plato.stanford.edu/entries/dedekind-foundations/}{Entry}. Peano acknowledged Dedekind's priority for the axioms.
\par\medskip

\noindent Neil Robertson, Daniel Sanders, Paul Seymour and Robin Thomas (1997). The four-colour theorem. Journal of Combinatorial Theory, Series B 70(1), pp. 2--44.
\par\medskip

\noindent Jürgen Schmidhuber (2003). Gödel machines: self-referential universal problem solvers making provably optimal self-improvements. Technical Report IDSIA-19-03. \href{https://arxiv.org/abs/cs/0309048}{arXiv: cs/0309048}.
\par\medskip

\noindent Marcel-Paul Schützenberger (1961). On the definition of a family of automata. Information and Control 4(2--3), pp. 245--270.
\par\medskip

\noindent Dana Scott (1962). A system of functional abstraction. Lectures delivered at the University of California, Berkeley, 1962--63; unpublished. The origin of the Scott encoding, as given in Mogensen--Scott encoding, Wikipedia, \href{https://en.wikipedia.org/wiki/Mogensen\%E2\%80\%93Scott_encoding}{Article}.
\par\medskip

\noindent Dana Scott. Carnegie Mellon University's School of Computer Science, \href{https://www.cs.cmu.edu/~scsfacts/turing.html}{faculty Turing Award recipients}: the 1976 award jointly with Michael Rabin for their work on nondeterministic machines. His doctorate at Princeton under Alonzo Church, domain theory and models of the lambda calculus: Wikipedia, \href{https://en.wikipedia.org/wiki/Dana_Scott}{Dana Scott}.
\par\medskip

\noindent Emin G\"un Sirer, Willem de Bruijn, Patrick Reynolds, Alan Shieh, Kevin Walsh, Dan Williams and Fred B. Schneider (2011). Logical attestation: an authorization architecture for trustworthy computing. SOSP 2011, pp. 249--264. \href{https://doi.org/10.1145/2043556.2043580}{DOI: 10.1145/2043556.2043580}.
\par\medskip

\noindent Ken Thompson. Heidelberg Laureate Forum, \href{https://www.heidelberg-laureate-forum.org/laureate/ken-thompson/}{laureate profile}: the 1983 Turing Award with Dennis Ritchie for operating systems theory and the implementation of Unix. Thompson's design and implementation of the original Unix at Bell Labs: Wikipedia, \href{https://en.wikipedia.org/wiki/Ken_Thompson}{Ken Thompson}.
\par\medskip

\noindent Alfred North Whitehead and Bertrand Russell (1910). Principia Mathematica, vol. I. Cambridge University Press. Introduction, chapter II.
\par\medskip

\noindent Robin Wilson (2002). Four Colours Suffice? Gresham College lecture, 21 October 2002. \href{https://www.gresham.ac.uk/node/11353}{Lecture}. De Morgan's letter of 23 October 1852 to the Irish mathematician William Rowan Hamilton.
\par\medskip

\subsection{Correlations and environmental records}



\noindent M.-C. Chen et al. (2019). Emergence of classical objectivity of quantum Darwinism in a photonic quantum simulator. Science Bulletin 64, p. 580.
\par\medskip

\noindent M. A. Ciampini, G. Pinna, P. Mataloni and M. Paternostro (2018). Experimental signature of quantum Darwinism in photonic cluster states. Physical Review A 98, 020101.
\par\medskip

\noindent Mauro Dorato and Jan Faye (2026). Niels Bohr as a Physicist of Principle. \href{https://arxiv.org/abs/2609.04297}{arXiv: 2609.04297}. Quotes Bohr's Unity of Knowledge (1954) on irreversible amplification and permanent marks on a photographic plate, citing p. 73.
\par\medskip

\noindent Lluís Masanes (2003). Necessary and sufficient condition for quantum-generated correlations. \href{https://arxiv.org/abs/quant-ph/0309137v1}{arXiv: quant-ph/0309137v1}. The paper gives his affiliation as the Universitat de Barcelona.
\par\medskip

\noindent Miguel Navascués, Stefano Pironio and Antonio Acín (2007). Bounding the Set of Quantum Correlations. Physical Review Letters 98, 010401. \href{https://doi.org/10.1103/PhysRevLett.98.010401}{DOI: 10.1103/PhysRevLett.98.010401}; \href{https://arxiv.org/abs/quant-ph/0607119v2}{arXiv: quant-ph/0607119v2}.
\par\medskip

\noindent Miguel Navascués, Stefano Pironio and Antonio Acín (2008). A convergent hierarchy of semidefinite programs characterizing the set of quantum correlations. New Journal of Physics 10, 073013. \href{https://arxiv.org/abs/0803.4290v1}{arXiv: 0803.4290v1}.
\par\medskip

\noindent Marcin Pawłowski, Tomasz Paterek, Dagomir Kaszlikowski, Valerio Scarani, Andreas Winter and Marek Żukowski (2009). Information causality as a physical principle. Nature 461, pp. 1101--1104. Pawłowski's affiliation, the University of Gdańsk, is on the paper: \href{https://arxiv.org/abs/0905.2292}{arXiv: 0905.2292}.
\par\medskip

\noindent Erwin Schrödinger. Wikipedia. \href{https://en.wikipedia.org/wiki/Erwin_Schr\%C3\%B6dinger}{Article}. Austrian physicist; the wave equation of 1926.
\par\medskip

\noindent B. S. Tsirel'son (1987). Quantum analogues of the Bell inequalities. The case of two spatially separated domains. Journal of Soviet Mathematics 36(4), pp. 557--570. \href{https://doi.org/10.1007/BF01663472}{DOI: 10.1007/BF01663472}.
\par\medskip

\noindent T. Unden, D. Louzon, M. Zwolak, W. H. Zurek and F. Jelezko (2019). Revealing the emergence of classicality using nitrogen-vacancy centers. Physical Review Letters 123, 140402.
\par\medskip

\cleardoublepage
\thispagestyle{empty}
\vspace*{154pt}
\begingroup\setlength{\parindent}{0pt}\setlength{\parskip}{7pt}
\noindent\BookLabel{Colophon}

\vspace{14pt}
This book is set in Charter, drawn by Matthew Carter in 1987 for the laser printers of its day,
10.5 on 14 point, on a page 189 by 246 millimetres. It was built from its source with pdfTeX.
The cells are drawn in TikZ.

Its own count, taken by the build: \ClaimTotal\ numbered statements, of which \PaidTotal\ are
closed by a proof on the page, and the rest stated and not proved: \OpenClaimList. One cell for
each, in the order they come, filled where the proof closed:

\vspace{7pt}
\noindent\LedgerGrid
\endgroup
\clearpage
\thispagestyle{empty}
\vspace*{0pt}
\noindent\BookLabel{\phantom{Devon Thiele}}

\vspace{14pt}
\PageCellFilled
\par\medskip

\end{document}
